% !TEX TS-program = lualatex
% !TEX encoding = UTF-8

\documentclass[antiphonaire_ordinaire.tex]{subfiles}

\ifcsname preamble@file\endcsname
  \setcounter{page}{\getpagerefnumber{M-ao_sanctis}}
\fi

\begin{document}

\feast{3}{0114}{saint Hilaire, évêque, confesseur et docteur de l'Église}{14}
\saintindex{Hilárii}{Hilaire}
\saintindex{Felicis}{Félix}

\oratio{Deus, qui pópulo tuo ætérnæ salútis beátum Hilárium minístrum tribuísti:\pscross{} præsta, quǽsumus; ut, quem Doctórem vitæ habúimus in terris,\psstar{} intercessórem habére mereámur in cælis. \perdominumla}{Ô Dieu, qui avez donné à votre peuple le bienheureux Hilaire comme ministre du salut éternel, nous vous en prions : faites que, l’ayant eu comme maître de vie sur la terre, nous méritions de l’avoir comme intercesseur dans le ciel. \perdominumfr}

\rubrica{À Laudes, comm. de saint Félix, prêtre et martyr, \aarub \scorename{UMEXLAB},% Qui odit
\vvrub \normaltext{Justus ut palma}.}

\oratio{Concede, quǽsumus, omnípotens Deus:\psstar{} ut ad meliórem vitam Sanctórum tuórum exémpla nos próvocent; quátenus, quorum solémnia ágimus, étiam actus imitémur. \perdominumla}{Faites, nous vous en prions, Dieu tout-puissant, que les exemples de Vos saints nous entraînent vers une vie plus parfaite, de telle sorte que, célébrant leur fête, nous imitions aussi leurs actions. \perdominumfr}

\mtv

\rubrica{À Vêpres, \aarub \scorename{CODOVAM}% O doctor optime
avec le nom \normaltext{Hilári}.}

%%%%%%%%%%%%%%%%%%%%%%%%%%%%%%%%%%%%%%%%%%%%%%%%%%%%%%

\feast{3}{0115}{saint Paul, premier ermite, confesseur}{15}
\saintindex{Pauli Primi Eremitæ}{Paul Premier Ermite}
\saintindex{Mauri}{Maur}

\rubrica{Oraison du commun, avec le nom \normaltext{Pauli}.

\rubrica{À Laudes, comm. de saint Maur, abbé, \aarub \scorename{CONPLAB},% Similabo
\vvrub \normaltext{Amávit eum}, et oraison du commun, avec le nom \normaltext{Mauri}.}

\mtv

%%%%%%%%%%%%%%%%%%%%%%%%%%%%%%%%%%%%%%%%%%%%%%%%%%%%%%

\feast{3}{0116}{saint Marcel, pape et martyr}{16}
\saintindex{Marcélli}{Marcel}

\oratio{Preces pópuli tui, quǽsumus, Dómine, cleménter exáudi: ut beáti Marcélli Mártyris tui atque Pontíficis méritis adjuvémur, cujus passióne lætámur. \perdominumla}{Exaucez dans Votre clémence, Seigneur, s'il Vous plaît, les prières de Votre peuple afin que nous soyons assistés par les mérites du bienheureux Marcel, Votre martyr et pontife, dont nous fêtons la passion. \perdominumfr}

%%%%%%%%%%%%%%%%%%%%%%%%%%%%%%%%%%%%%%%%%%%%%%%%%%%%%%

\feast{3}{0117}{saint Antoine, abbé}{17}

\saintindex{Antónii}{Antoine}

\oratio{Intercéssio nos, quǽsumus Dómine, beáti Antónii Abbátis comméndet:\pscross{} ut quod nostris méritis non valémus,\psstar{} ejus patrocínio assequámur. \perdominumla}{Que l’intercession du bienheureux Antoine, Abbé, nous soit recommandation, nous vous le demandons, Seigneur, pour que nous obtenions, par son patronage, ce qui dépasse le pouvoir de nos mérites. \perdominumfr}

%%%%%%%%%%%%%%%%%%%%%%%%%%%%%%%%%%%%%%%%%%%%%%%%%%%%%%

\feast{4}{0118}{sainte Prisque, vierge et martyre}{18}
\saintindex{Priscæ}{Prisque}

\oratio{Da, quǽsumus, omnípotens Deus:\psstar{} ut, qui beátæ Priscæ Vírginis et Mártyris tuæ natalícia cólimus; et ánnua solemnitáte lætémur, et tantæ fídei proficiámus exémplo. \perdominumla}{Accordez à notre demande, ô Dieu tout-puissant, que célébrant la naissance céleste de la bienheureuse Prisque, Vierge et Martyre, nous nous réjouissions de sa fête annuelle et profitions de l’exemple d’une si grande foi. \perdominumfr}

%%%%%%%%%%%%%%%%%%%%%%%%%%%%%%%%%%%%%%%%%%%%%%%%%%%%%%

\feast{4}{0119}{saints Maris, Marthe, Audifax et Abacum, martyrs}{19}
\saintindex{Marii}{Maris}
\saintindex{Marthæ (mart.)}{Marthe (mart.)}
\saintindex{Audifacis}{Audifax}
\saintindex{Abachum}{Abacum}
\saintindex{Canuti}{Knut}

\oratio{Exáudi, Dómine, pópulum tuum cum Sanctórum tuórum patrocínio supplicántem:\psstar{} ut et temporális vitæ nos tríbuas pace gaudére; et ætérnæ reperíre subsídium. \perdominumla}{Exaucez, Seigneur, les supplications de votre peuple, jointes au patronage de vos saints, afin que vous nous accordiez de jouir de la paix en la vie temporelle et d’obtenir le secours pour celle de l’éternité. \perdominumfr}

\rubrica{À Laudes, comm. de saint Knut, roi et martyr, \aarub \scorename{UMEXLAB},% Qui odit
\vvrub \normaltext{Justus ut palma}.}

\oratio{Deus, qui ad illustrándam Ecclésiam tuam beátum Canútum regem martýrii palma et gloriósis miráculis decoráre dignátus es:\pscross{} concéde propítius; ut, sicut ipse Domínicæ passiónis imitátor fuit, ita nos, per ejus vestígia gradiéntes,\psstar{} ad gáudia sempitérna perveníre mereámur. \pereumdemla}{Ô Dieu qui, pour glorifier votre Église, avez daigné honorer le bienheureux roi Canut de la palme du martyre et d’éclatants miracles, accordez-nous miséricordieusement, que, comme il a lui-même imité la passion du Seigneur, nous aussi, marchant sur ses traces, nous méritions d’arriver aux joies éternelles. \pereumdemfr}

%%%%%%%%%%%%%%%%%%%%%%%%%%%%%%%%%%%%%%%%%%%%%%%%%%%%%%

\feast{3}{0120}{saints Fabien, pape, et Sébastien, martyrs}{20}
\saintindex{Fabiani}{Fabien}
\saintindex{Sebastiani}{Sébastien}

\oratio{Infirmitátem nostram réspice, omnípotens Deus:\pscross{} et, quia pondus própriæ actiónis gravat,\psstar{} beatórum Mártyrum tuórum Fabiáni et Sebastiáni intercéssio gloriósa nos prótegat. \perdominumla}{Regardez notre faiblesse, Dieu tout-puissant ; et puisque le poids de notre propre activité nous alourdit, faites que l’intercession glorieuse de Vos bienheureux Martyrs Fabien et Sébastien nous protège. \perdominumfr}


%%%%%%%%%%%%%%%%%%%%%%%%%%%%%%%%%%%%%%%%%%%%%%%%%%%%%%

\feast{2}{0121}{sainte Agnès, vierge et martyre}{21 janvier}
\saintindex{Agnetis}{Agnès}

\rubrica{Matines, Laudes et Vêpres à l'antiphonaire. Petites heures au psautier: capitules propres, répons brefs et versets du Commun.}

\oratio{Omnípotens sempitérne Deus, qui infírma mundi éligis, ut fórtia quæque confúndas:\pscross{} concéde propítius; ut, qui beátæ Agnétis Vírginis et Mártyris tuæ solémnia cólimus,\psstar{} ejus apud te patrocínia sentiámus. \perdominumla}{Ô Dieu tout-puissant et éternel, qui choisissez ce qui est faible en ce monde pour confondre tout ce qui est fort, faites dans votre bonté, que, célébrant les fêtes de la bienheureuse Agnès, votre Vierge et Martyre, nous ressentions auprès de vous les effets de sa protection. \perdominumfr}

\titulumrubrica{À Tierce}
\capitulum{Si. 51: 1-2}{Confitébor tibi, Dómine, Rex, et collaudábo te Deum Salvatórem meum.\pscross{} Confitébor nómini tuo: quóniam adjútor et protéctor factus es mihi,\psstar{} et liberásti corpus meum a perditióne.}{ Je veux te rendre grâce, Seigneur mon roi, je te louerai, Dieu mon sauveur. Je rends grâce à ton nom : tu as été pour moi un défenseur et un soutien, tu as délivré mon corps de la perdition.}

\titulumrubrica{À Sexte}
\capitulum{Si. 51: 3-4}{Liberásti me secúndum multitúdinem misericórdiæ nóminis tui a rugiéntibus, præparátis ad escam,\pscross{} de mánibus quæréntium ánimam meam,\psstar{} et de multis tribulatiónibus quæ circumdedérunt me.}{Par la grandeur de ta miséricorde et de ton nom, tu m’as délivré des mâchoires qui allaient me dévorer, des mains qui menaçaient ma vie, de tourments innombrables.}

\titulumrubrica{À None}
\capitulum{Si. 51: 5; 51:8}{Laudábit usque ad mortem ánima mea Dóminum,\pscross{} quóniam éruis sustinéntes te, et líberas eos de manu angústiæ,\psstar{} Dómine, Deus noster.}{Du gouffre profond de la mort, mon âme louera le Seigneur, car tu retires de la peine ceux qui t’attendent, tu les sauves de la main des ennemis, Seigneur, notre Dieu.}

%%%%%%%%%%%%%%%%%%%%%%%%%%%%%%%%%%%%%%%%%%%%%%%%%%%%%%

\feast{3}{0122}{saints Vincent et Anastase, martyrs}{22}
\saintindex{Vincentii}{Vincent}
\saintindex{Anastasii}{Anastase}

\oratio{Adésto, Dómine, supplicatiónibus nostris:\pscross{} ut, qui ex iniquitáte nostra reos nos esse cognóscimus,\psstar{} beatórum Mártyrum tuórum Vincéntii et Anastásii intercessióne liberémur. \perdominumla}{Ecoutez, Seigneur, nos supplications, afin que, reconnaissant l’iniquité qui nous rend coupables, nous en soyons délivrés par l’intercession de vos bienheureux Martyrs, Vincent et Anastase. \perdominumfr}

%%%%%%%%%%%%%%%%%%%%%%%%%%%%%%%%%%%%%%%%%%%%%%%%%%%%%%

\feast{3}{0123}{saint Raymond de Peñafort, confesseur}{23}
\saintindex{Raymundi}{Raymond}
\saintindex{Emerentianæ}{Émérentienne}

\oratio{Deus, qui beátum Raymúndum pœniténtiæ sacraménti insígnem minístrum elegísti, et per maris undas mirabíliter traduxísti:\pscross{} concéde; ut ejus intercessióne dignos pœniténtiæ fructus fácere,\psstar{} et ad ætérnæ salútis portum perveníre valeámus. \perdominumla}{Ô Dieu, qui avez choisi le bienheureux Raymond pour ministre insigne du sacrement de pénitence et lui avez fait traverser les eaux de la mer miraculeusement, accordez-nous de pouvoir, par son intercession, opérer de dignes fruits de pénitence et ainsi, parvenir au port du salut éternel. \perdominumfr}

\rubrica{Comm. de sainte Émérentienne, vierge et martyre.}

\oratio{Indulgéntiam nobis, quǽsumus, Dómine, beáta Emerentiána Virgo et Martyr implóret:\psstar{} quæ tibi grata semper éxstitit, et mérito castitátis, et tuæ professióne virtútis. \perdominumla}{Que la bienheureuse Vierge et Martyre Émérentienne, Seigneur, implore pour nous le pardon, elle qui toujours vous fut agréable, tant par le mérite de sa chasteté que par le témoignage qu'elle rendit à votre puissance. \perdominumfr}

\mtv

%%%%%%%%%%%%%%%%%%%%%%%%%%%%%%%%%%%%%%%%%%%%%%%%%%%%%%

\feast{3}{0124}{saint Timothée, évêque et martyr}{24}
\saintindex{Timóthei}{Timothée}

\rubrica{Oraison du commun, avec le nom \normaltext{Timóthei}.}

\oratio{Infirmitátem nostram réspice, omnípotens Deus:\pscross{} et, quia pondus própriæ actiónis gravat,\psstar{} beáti Timóthei Mártyris tui atque Pontíficis intercéssio gloriósa nos prótegat. \perdominumla}{Voyez notre infirmité, Dieu tout-puissant : et qu’à cause du poids de notre activité personnelle qui nous alourdit, la glorieuse intercession de votre bienheureux Timothée, Martyr et Pontife, nous protège. \perdominumfr}

%%%%%%%%%%%%%%%%%%%%%%%%%%%%%%%%%%%%%%%%%%%%%%%%%%%%%%

\feast{2}{0125}{Conversion de saint Paul, apôtre}{25 janvier}
\saintindex{Pauli!Conversio}{Paul!Conversion}

\rubrica{Matines, Laudes et Vêpres à l'antiphonaire. Petites heures au psautier: capitules propres, répons brefs et versets du Commun.}

\oratio{Deus, qui univérsum mundum beáti Pauli Apóstoli prædicatióne docuísti:\pscross{} da nobis, quǽsumus; ut, qui ejus hódie Conversiónem cólimus,\psstar{} per ejus ad te exémpla gradiámur. \perdominumla}{Ô Dieu, qui avez instruit le monde entier par la parole du bienheureux Apôtre Paul, accordez-nous, à nous qui célébrons aujourd’hui sa Conversion, de marcher vers vous en suivant ses exemples. \perdominumfr}

\titulumrubrica{À Tierce}
\capitulum{Ac. 9: 1-2}{Saulus autem adhuc spirans minárum et cædis in discípulos Dómini, accéssit ad príncipem sacerdótum,\pscross{} et pétiit ab eo epístolas in Damáscum ad synagógas:\psstar{} ut si quos invenísset hujus viæ viros ac mulíeres, vinctos perdúceret in Jerúsalem.}{Saul était toujours animé d’une rage meurtrière contre les disciples du Seigneur. Il alla trouver le grand prêtre et lui demanda des lettres pour les synagogues de Damas, afin que, s’il trouvait des hommes et des femmes qui suivaient le Chemin du Seigneur, il les amène enchaînés à Jérusalem.}

\titulumrubrica{À Sexte}
\capitulum{Ac. 9: 8-9}{Surréxit autem Saulus de terra, apertísque óculis nihil vidébat.\pscross{} Ad manus autem illum trahéntes, introduxérunt Damáscum.\psstar{} Et erat ibi tribus diébus non videns, et non manducávit neque bibit.}{Saul se releva de terre et, bien qu’il eût les yeux ouverts, il ne voyait rien. Ils le prirent par la main pour le faire entrer à Damas. Pendant trois jours, il fut privé de la vue et il resta sans manger ni boire.}

\titulumrubrica{À None}
\capitulum{Ac. 9: 22}{Saulus autem multo magis convalescébat,\pscross{} et confundébat Judǽos qui habitábant Damásci,\psstar{} affírmans quóniam hic est Christus.}{Saul, avec une force de plus en plus grande, réfutait les Juifs qui habitaient Damas, en démontrant que Jésus est le Christ.}

%%%%%%%%%%%%%%%%%%%%%%%%%%%%%%%%%%%%%%%%%%%%%%%%%%%%%%

\feast{3}{0126}{saint Polycarpe, évêque et martyr}{26}
\saintindex{Polycarpi}{Polycarpe}

\rubrica{Oraison du commun, avec le nom \normaltext{Polycárpi}.}

\oratio{Deus, qui nos beáti Polycárpi Mártyris tui atque Pontíficis ánnua solemnitáte lætíficas:\pscross{} concéde propítius; ut, cujus natalítia cólimus,\psstar{} de ejúsdem étiam protectióne gaudeámus. \perdominumla}{Dieu, qui nous réjouissez par la fête annuelle de votre bienheureux Polycarpe, Martyr et Pontife, faites-nous cette faveur que, célébrant sa naissance céleste, nous nous réjouissions aussi de sa protection. \perdominumfr}

%%%%%%%%%%%%%%%%%%%%%%%%%%%%%%%%%%%%%%%%%%%%%%%%%%%%%%

\feast{3}{0127}{saint Jean Chrysostome, évêque, confesseur et docteur de l'Église}{27}
\saintindex{Joannis Chrysostomi}{Jean Chrysostome}

\oratio{Ecclésiam tuam, quǽsumus, Dómine, grátia cæléstis amplíficet:\psstar{} quam beáti Joánnis Chrysóstomi Confessóris tui atque Pontíficis illustráre voluísti gloriósis méritis et doctrínis. \perdominumla}{Que la grâce céleste, nous vous en prions, Seigneur, fasse grandir votre Église que vous avez voulu illustrer par les enseignements et les glorieux mérites du bienheureux Jean Chrysostome, votre Confesseur et Pontife. \perdominumfr}

\mtv

\rubrica{À Vêpres, \aarub \scorename{CODOVAM}% O doctor optime
avec le nom \normaltext{beáte Joánnes Chrysóstome}.}

%%%%%%%%%%%%%%%%%%%%%%%%%%%%%%%%%%%%%%%%%%%%%%%%%%%%%%

\feast{3}{0128}{saint Pierre Nolasque, confesseur}{28 janvier}
\saintindex{Petri Nolasci}{Pierre Nolasque}
\saintindex{Agnetis!Octava}{Agnès!Octave}

\oratio{Deus, qui in tuæ caritátis exémplum ad fidélium redemptiónem sanctum Petrum Ecclésiam tuam nova prole fecundáre divínitus docuísti:\pscross{} ipsíus nobis intercessióne concéde; a peccáti servitúte solútis,\psstar{} in cælésti pátria perpétua libertáte gaudére: \quivivisla}{Ô Dieu qui, en exemple de votre charité, et pour le rachat des fidèles, avez divinement inspiré saint Pierre de doter votre Église d’une nouvelle famille ; accordez-nous, par son intercession, que, délivrés de la servitude du péché, nous jouissions de la liberté perpétuelle dans la céleste patrie : \quivivisfr}

\rubrica{Commémoraison de la seconde fête de sainte Agnès:}

\gscore{0121LAB}{}{Je vois maintenant ce que j’ai ardemment désiré ; ce que j’ai espéré, je le possède enfin ; je suis unie dans les cieux à celui que, sur terre, j’ai aimé de toute mon âme.}

%% TODO verset diffusa est

\oratio{Deus, qui nos ánnua beátæ Agnétis Vírginis et Mártyris tuæ solemnitáte lætíficas:\pscross{} da, quǽsumus; ut, quam venerámur offício,\psstar{} étiam piæ conversatiónis sequámur exémplo. \perdominumla}{Ô Dieu, qui nous réjouissez par la solennité annuelle de la bienheureuse Agnès, votre Vierge et Martyre ; faites, nous vous en supplions, que celle dont nous célébrons l’office avec vénération, nous l’imitions par l’exemple d’une sainte vie. \perdominumfr}

\rubrica{Si on commémore aussi à Vêpres la seconde fête de sainte Agnès, on emploie l'antienne ci-dessous avec le verset \normaltext{Spécie tua}:}

\gscore{0128VAX}{}{Debout à sa droite, Agneau plus blanc que neige, le Christ se l’est consacrée épouse et martyre.}

%%%%%%%%%%%%%%%%%%%%%%%%%%%%%%%%%%%%%%%%%%%%%%%%%%%%%%

\feast{3}{0129}{saint François de Sales, évêque, confesseur et docteur de l'Église}{29}
\saintindex{Francisci Salesii}{François de Sales}

\oratio{Deus, qui ad animárum salútem beátum Francíscum Confessórem tuum atque Pontíficem ómnibus ómnia factum esse voluísti:\pscross{} concéde propítius; ut caritátis tuæ dulcédine perfúsi,\psstar{} ejus dirigéntibus mónitis ac suffragántibus méritis, ætérna gáudia consequámur. \perdominumla}{Ô Dieu qui, pour le salut des âmes, avez voulu que le bienheureux François votre Confesseur et Pontife se fît tout à tous, accordez-nous dans votre bonté que, pénétrés de la douceur de sa charité, dirigés par ses enseignements et aidés de ses mérites, nous obtenions les joies éternelles. \perdominumfr}

\mtv

\rubrica{À Vêpres, \aarub \scorename{CODOVAM}% O doctor optime
avec le nom \normaltext{beáte Francísce}.}

%%%%%%%%%%%%%%%%%%%%%%%%%%%%%%%%%%%%%%%%%%%%%%%%%%%%%%

\feast{0130}{S. Martinæ Virginis et Martyris}{30 janvier}
\saintindex{Martinæ}{Martine}

\rubrica{À Laudes et Vêpres:}

\gscore{0130H}{}{Protège le sol qui t'a vu naître ;
Accorde un repos paisible à la terre des Chrétiens ;
Renvoie sur le pays infidèle des Thraces
Le bruit des armes et les cruels combats.

Rassemble tous les rois avec leurs bataillons.
Sous l'étendard de la croix, délivre Jérusalem de la captivité,
Venge le sang innocent,
Et renverse à jamais les remparts du Turc notre ennemi.

Ô Vierge, notre appui, notre gloire éclatante,
Reçois l'hommage de nos cœurs.
Agrée les vœux de Rome
Qui te chante et t'honore dans son amour.

Eloignez de nous les jours dangereux,
Ô Dieu, dont le bras soutient les Martyrs ;
Unité, Trinité, donnez à vos serviteurs cette lumière
Par laquelle vous daignez faire le bonheur des âmes.}

\oratio{Deus, qui inter cétera poténtiæ tuæ mirácula étiam in sexu frágili victóriam martýrii contulísti:\pscross{} concéde propítius; ut, qui beátæ Martínæ Vírginis et Mártyris tuæ natalícia cólimus,\psstar{} per ejus ad te exémpla gradiámur. \perdominumla}{Ô Dieu, qui, entre autres miracles de Votre puissance, avez accordé même au sexe faible la victoire du martyre, faites, dans Votre bonté, qu'honorant la naissance de Votre bienheureuse Vierge et Martyre Martine, nous allions à Vous en imitant ses exemples. \perdominumfr}

%%%%%%%%%%%%%%%%%%%%%%%%%%%%%%%%%%%%%%%%%%%%%%%%%%%%%%

\feast{3}{0131}{saint Jean Bosco, confesseur}{31}
\saintindex{Joannis Bosco}{Jean Bosco}

\oratio{Deus, qui sanctum Joánnem Confessórem tuum adolescéntium patrem et magístrum excitásti, ac per eum, auxiliatríce Vírgine María, novas in Ecclésia tua famílias floréscere voluísti:\pscross{} concéde, quǽsumus;\psstar{} ut eódem caritátis igne succénsi, ánimas quǽrere, tibíque soli servíre valeámus. \perdominumla}{Ô Dieu, qui avez suscité saint Jean, votre Confesseur, comme père et maître des adolescents, et par lui, avec le secours de la Vierge Marie, avez voulu faire fleurir dans votre Église de nouvelles familles, accordez à notre demande, qu’enflammés du même feu de charité, nous puissions vous chercher les âmes et vous servir vous seul. \perdominumfr}

%%%%%%%%%%%%%%%%%%%%%%%%%%%%%%%%%%%%%%%%%%%%%%%%%%%%%%

\feast{0201}{S. Ignatii Episcopi et Martyris}{1 février}
\saintindex{Ignátii}
\saintindex{Ignace}
\oratio{}{Voyez notre infirmité, Dieu tout-puissant, et qu’à cause du poids de notre propre activité qui nous alourdit, la glorieuse intercession de votre bienheureux Ignace, Martyr et Pontife, nous protège. \perdominumfr}
\feast{0202}{}{2 février}
\saintindex{}
\saintindex{}
\feast{0202}{}{2 février}
\saintindex{}
\saintindex{}
\oratio{Omnípotens sempitérne Deus, majestátem tuam súpplices exorámus: ut, sicut unigénitus Fílius tuus hodiérna die cum nostræ carnis substántia in templo est præsentátus; ita nos fácias purificátis tibi méntibus præsentári. \pereumdemla}{Dieu tout-puissant et éternel, nous implorons suppliants votre majesté pour que, de même que votre Fils unique a été présenté aujourd’hui revêtu de notre chair dans votre temple, ainsi, nous puissions vous être présentés avec des cœurs purifiés. \pereumdemfr}
\feast{0203}{S. Blasii Episcopi et Martyris}{3 février}
\saintindex{Blásii}
\saintindex{Blaise}
\oratio{}{Ô Dieu, qui nous réjouissez par la fête annuelle de Votre bienheureux Blaise, Martyr et Pontife, accordez-nous cette faveur qu’en célébrant sa naissance céleste, nous nous réjouissions aussi de sa protection. \perdominumfr}
\feast{0204}{S. Andreæ Corsini Episcopi et Confessoris}{4 février}
\saintindex{Andréæ}
\saintindex{André}
\oratio{Deus, qui in Ecclésia tua, nova semper instáuras exémpla virtútum: da pópulo tuo beáti Andréæ Confessóris tui atque Pontíficis ita sequi vestígia; ut assequátur et prǽmia. \perdominumla}{Ô Dieu, qui suscitez toujours dans votre Église de nouveaux exemples de vertu, accordez à votre peuple, de suivre si bien les traces du bienheureux André votre Confesseur et Pontife, qu’il en obtienne aussi les récompenses. \perdominumfr}
\feast{0205}{S. Agathæ Virginis et Martyris}{5 février}
\saintindex{Agathæ}
\saintindex{Agathe}
\oratio{}{Ô Dieu qui, parmi les autres miracles de Votre puissance, avez donné au sexe faible la victoire du martyre, accordez-nous miséricordieusement que, fêtant la naissance céleste de la bienheureuse Agathe, Votre Vierge et Martyre, nous nous servions de ses exemples pour monter vers Vous. \perdominumfr}
\feast{0206}{S. Titi Episcopi et Confessoris}{6 février}
\saintindex{Titum}
\saintindex{Tite}
\oratio{Deus, qui beátum Titum Confessórem tuum atque Pontíficem apostólicis virtútibus decorásti: ejus méritis et intercessióne concéde; ut juste et pie vivéntes in hoc sǽculo, ad cæléstem pátriam perveníre mereámur. \perdominumla}{Ô Dieu qui avez orné le bienheureux Tite votre Confesseur et Pontife, des vertus apostoliques accordez-nous, par ses mérites et son intercession, que, vivant justement et pieusement en ce monde, nous méritions de parvenir à la céleste patrie. \perdominumfr}
\feast{0206}{S. Dorotheæ Virginis et Martyris}{6 février}
\saintindex{Doróthea}
\saintindex{Dorothée}
\feast{0207}{S. Romualdi Abbatis}{7 février}
\saintindex{Romuáldi}
\saintindex{Romuald}
\feast{0208}{S. Joannis de Matha Confessoris}{8 février}
\saintindex{}
\saintindex{}
\oratio{Deus, qui per sanctum Joánnem órdinem sanctíssimæ Trinitátis ad rediméndum de potestáte Saracenórum captívos cǽlitus institúere dignátus es: præsta, quǽsumus; ut, ejus suffragántibus méritis, a captivitáte córporis et ánimæ, te adjuvánte, liberémur. \perdominumla}{Ô Dieu qui, par saint Jean, avez daigné instituer miraculeusement l'Ordre de la Très Sainte Trinité pour racheter les captifs du pouvoir des Sarrasins ; faites, nous vous le demandons, que, par les suffrages de ses mérites, nous soyons délivrés, avec votre secours, de la captivité du corps et de l’âme. \perdominumfr}
\feast{0209}{S. Cyrilli Episc. Alexandrini Confessoris et Ecclesiæ Doctoris}{9 février}
\saintindex{Cyrílle}
\saintindex{Cyrílle}
\oratio{Deus, qui beátum Cyríllum Confessórem tuum atque Pontíficem divínæ maternitátis beatíssimæ Vírginis Maríæ assertórem invíctum effecísti: concéde, ipso intercedénte; ut, qui vere eam Genetrícem Dei crédimus, matérna eíusdem protectióne salvémur. \pereumdemla}{Ô Dieu qui avez fait du bienheureux Cyrille, votre Confesseur et Pontife, le victorieux défenseur de la divine maternité de la bienheureuse Vierge Marie, accordez-nous par son intercession, à nous qui la croyons vraiment Mère de Dieu, d’être sauvés par sa maternelle protection. \pereumdemfr}
\feast{0209}{S. Apolloniæ Virginis et Martyris}{9 février}
\saintindex{Apollóniæ}
\saintindex{Apolline}
\feast{0210}{S. Scholasticæ Virginis}{10 février}
\saintindex{}
\saintindex{}
\oratio{Deus, qui ánimam beátæ Vírginis tuæ Scholásticæ ad ostendéndam innocéntiæ viam in colúmbæ spécie cælum penetráre fecísti: da nobis ejus méritis et précibus ita innocénter vívere; ut ad ætérna mereámur gáudia perveníre. \perdominumla}{Ô Dieu qui, pour manifester la voie d’innocence suivie par votre bienheureuse Vierge Scholastique, avez fait monter au ciel son âme sous la forme d’une colombe, donnez-nous, par ses mérites et ses prières, de vivre avec une telle innocence que nous méritions de parvenir aux joies éternelles. \perdominumfr}
\feast{0211}{}{11 février}
\saintindex{}
\saintindex{}
\oratio{Deus, qui per immaculátam Vírginis Conceptiónem dignum Fílio tuo habitáculum præparásti: súpplices a te quǽsumus; ut, ejúsdem Vírginis Apparitiónem celebrántes, salútem mentis et córporis consequámur. \pereumdemla}{Ô Dieu, qui par l’immaculée Conception de la Vierge avez préparé à votre Fils une habitation digne de lui, faites, nous vous demandons suppliants, que, célébrant l’Apparition de cette même Vierge, nous obtenions la santé de l’âme et du corps. \pereumdemfr}
\feast{0212}{Ss. Septem Fundatorum Ordinis Servorum B. M. V.}{12 février}
\saintindex{}
\saintindex{}
\oratio{Dómine Jesu Christe, qui ad recoléndam memóriam dolórum sanctíssimæ Genetrícis tuæ, per septem beátos Patres nova Servórum ejus família Ecclésiam tuam fecundásti: concéde propítius; ita nos eórum consociári flétibus, ut perfruámur et gáudiis.
$Qui vivis}{Seigneur Jésus-Christ qui, pour honorer le souvenir des souffrances de votre très sainte Mère, avez, par sept bienheureux Pères, enrichi votre Église de la nouvelle famille de ses Serviteurs, accordez-nous la faveur de nous associer à leurs larmes pour que nous jouissions de leurs joies.
$Qui vivis}
\feast{0214}{S. Valentini Presbyteri et Martyris}{14 février}
\saintindex{Valentini}
\saintindex{Valentini}
\oratio{Præsta, quǽsumus, omnípotens Deus: ut, qui beáti Valentíni, Mártyris tui, natalítia cólimus, a cunctis malis imminéntibus, ejus intercessióne, liberémur. \perdominumla}{Faites, nous vous en prions, Dieu tout-puissant, que nous, qui honorons la naissance du bienheureux Valentin, votre Martyr, nous soyons délivrés par son intercession, de tous les maux qui nous menacent. \perdominumfr
}
\feast{0215}{SS. Faustini et Jovitæ}{15 février}
\saintindex{Faustíni et Jovítæ}
\saintindex{Faustin et Jovite}
\oratio{}{Ô Dieu, qui nous réjouissez par la solennité annuelle de Vos saints Martyrs Faustin et Jovite, accordez-nous, dans Votre bonté, qu’en nous réjouissant de leurs mérites, nous soyons enflammés par leurs exemples. \perdominumfr}
\feast{0218}{S. Simeonis Episcopi et Martyris}{18 février}
\saintindex{Simeónis}
\saintindex{Siméon}
\oratio{}{Voyez notre faiblesse, Dieu tout-puissant, et, parce que le poids de notre action propre nous accable, faites que l’intercession glorieuse du bienheureux Siméon, votre Martyr et Pontife, nous protège. \perdominumfr}
\feast{0222}{In Cathedra S. Petri Apostoli Antiochiæ}{22 février}
\saintindex{}
\saintindex{}
\oratio{Deus, qui beáto Petro Apóstolo tuo, collátis clávibus regni cæléstis, ligándi atque solvéndi pontifícium tradidísti: concéde; ut, intercessiónis ejus auxílio, a peccatórum nostrórum néxibus liberémur:
$Qui vivis
_
@Sancti/02-22:Commemoratio4}{Ô Dieu qui, en confiant au bienheureux Pierre votre Apôtre les clefs du royaume céleste, lui avez confié le pouvoir de lier et de délier, accordez-nous d’être, par le secours de son intercession, libérés des liens de nos péchés.
$Qui vivis
_
@Sancti/02-22:Commemoratio4}
\feast{0223}{In Vigilia S. Matthiæ Ap.}{23 février}
\saintindex{Matthíæ}
\saintindex{Matthias}
\oratio{}{Faites, nous vous en prions, Dieu tout-puissant, que la vénérable solennité anticipée du bienheureux Matthias, votre Apôtre, augmente en nous la dévotion et Tassurance de notre salut. \perdominumfr
}
\feast{0223}{S. Petri Damiani Episcopi Confessoris et Ecclesiæ Doctoris}{23 février}
\saintindex{Petre}
\saintindex{Petre}
\oratio{Concéde nos, quǽsumus, omnípotens Deus: beáti Petri Confessóris tui atque Pontíficis mónita et exémpla sectári; ut per terréstrium rerum contémptum ætérna gáudia consequámur. \perdominumla}{Accordez-nous, nous vous le demandons, Dieu tout-puissant, de suivre les leçons et les exemples du bienheureux Pierre votre Confesseur et Pontife, afin que, par le mépris des choses terrestres, nous obtenions les joies éternelles. \perdominumfr}
\feast{0224}{S. Matthiæ Apostoli}{24 février}
\saintindex{Matthíam}
\saintindex{Matthias}
\oratio{Deus, qui beátum Matthíam Apostolórum tuórum collégio sociásti: tríbue, quǽsumus; ut, ejus interventióne, tuæ circa nos pietátis semper víscera sentiámus. \perdominumla}{0 Dieu, qui avez associé le bienheureux Matthias au collège de vos Apôtres, accordez-nous, s’il vous plaît, que, par son intercession, nous ressentions toujours les effets de vôtre miséricorde à notre égard. \perdominumfr}
\feast{0227}{S. Gabrielis a Virgine Perdolente Confessoris}{27 février}
\saintindex{}
\saintindex{}
\oratio{Deus, qui beátum Gabriélem dulcíssimæ Matris tuæ dolóres assídue recólere docuísti, ac per illam sanctitátis et miraculórum glória sublimásti: da nobis, ejus intercessióne et exémplo; ita Genetrícis tuæ consociári flétibus, ut matérna ejúsdem protectióne salvémur:
$Qui vivis}{Ô Dieu, qui avez appris au bienheureux Gabriel à vénérer assidûment les douleurs de votre très douce Mère et qui par elle l’avez élevé à la gloire de la sainteté et des miracles, accordez-nous, par son intercession et son exemple, de si bien nous associer aux larmes de votre Mère, que nous soyons sauvés par son intercession.
$Qui vivis}
\feast{0304}{S. Casimiri Confessoris}{4 mars}
\saintindex{}
\saintindex{}
\oratio{Deus, qui inter regáles delícias et mundi illécebras sanctum Casimírum virtúte constántiæ roborásti: quǽsumus; ut ejus intercessióne fidéles tui terréna despíciant, et ad cæléstia semper aspírent. \perdominumla}{Ô Dieu, qui, parmi les délices royales et les séductions du monde, avez fortifié saint Casimir, par la vertu de constance, faites, nous vous demandons, que, par son intercession, vos fidèles méprisent les biens de la terre et aspirent toujours à ceux du ciel. \perdominumfr}
\feast{0304}{S. Lucii I Papæ et Martyris}{4 mars}
\saintindex{Lúcii}
\saintindex{Lucien}
\feast{0306}{Ss. Perpetuæ et Felicitatis Martyrum}{6 mars}
\saintindex{Perpétuæ et Felicitátis}
\saintindex{Perpétue et Félicité}
\oratio{}{Accordez-nous, nous vous en prions, Seigneur notre Dieu, de vénérer avec une constante dévotion les triomphes de vos saintes Martyres Perpétue et Félicité, afin que ne pouvant les célébrer avec un cœur digne d'elles, nous les poursuivions du moins de nos humbles hommages. \perdominumfr}
\feast{0307}{S. Thomæ de Aquino Confessoris et Ecclesiæ Doctoris}{7 mars}
\saintindex{Thoma}
\saintindex{Thomas}
\oratio{Deus, qui Ecclésiam tuam beáti Thomæ Confessóris tui mira eruditióne claríficas, et sancta operatióne fœcúndas: da nobis, quǽsumus; et quæ dócuit, intelléctu conspícere, et quæ egit, imitatióne complére. \perdominumla}{Ô Dieu qui illustrez votre Église par l'admirable science du bienheureux Thomas votre Confesseur et la fécondez par sa sainte opération, donnez-nous, nous vous en prions, de contempler par l'intelligence ce qu’il a enseigné et d’accomplir, en l’imitant, ce qu’il a fait. \perdominumfr}
\feast{0308}{S. Joannis de Deo Confessoris}{8 mars}
\saintindex{}
\saintindex{}
\oratio{Deus, qui beátum Joánnem, tuo amóre succénsum, inter flammas innóxium incédere fecísti, et per eum Ecclésiam tuam nova prole fecundásti: præsta, ipsíus suffragántibus méritis; ut igne caritátis tuæ vítia nostra curéntur, et remédia nobis ætérna provéniant. \perdominumla.}{Ô Dieu, qui avez fait marcher indemne au milieu des flammes le bienheureux Jean embrasé de votre amour et qui, par lui, avez enrichi votre Église d’une nouvelle famille, faites que par l’intercession de ses mérites, le feu de votre charité guérisse nos misères, et nous obtienne les remèdes éternels. \perdominumfr.}
\feast{0309}{S. Franciscæ Romanæ Viduæ}{9 mars}
\saintindex{}
\saintindex{}
\oratio{Deus, qui beátam Francíscam fámulam tuam, inter cétera grátiæ tuæ dona, familiári Angeli consuetúdine decorásti: concéde, quǽsumus; ut, intercessiónis ejus auxílio, Angelórum consórtium cónsequi mereámur. \perdominumla}{Ô Dieu, qui, entre autres dons de votre grâce, avez honoré la bienheureuse Françoise, votre servante, de l'habituelle familiarité d’un Ange, accordez-nous, nous vous en prions, de mériter par le secours de son intercession, d’être introduits dans la société des Anges. \perdominumfr}
\feast{0310}{Ss. Quadraginta Martyrum}{10 mars}
\saintindex{}
\saintindex{}
\oratio{Præsta, quǽsumus, omnípotens Deus: ut, qui gloriósos Mártyres fortes in sua confessióne cognóvimus, pios apud te in nostra intercessióne sentiámus. \perdominumla}{Faites, nous vous le demandons, Dieu tout-puissant, qu’ayant reconnu la force de ces glorieux Martyrs dans la confession de leur foi, nous sentions l’effet de leur piété dans leur intercession pour nous. \perdominumfr}
\feast{0312}{S. Gregorii Papæ Confessoris et Ecclesiæ Doctoris}{12 mars}
\saintindex{Gregóri}
\saintindex{Gregori}
\oratio{Deus, qui ánimæ fámuli tui Gregórii ætérnæ beatitúdinis prǽmia contulísti: concéde propítius; ut, qui peccatórum nostrórum póndere prémimur, ejus apud te précibus sublevémur. \perdominumla}{Ô Dieu, qui avez accordé à l’âme de votre serviteur Grégoire la récompense de la béatitude éternelle, faites-nous cette miséricorde, qu’accablés sous le poids de nos péchés, nous soyons soulagés par ses prières auprès de vous. \perdominumfr}
\feast{0317}{S. Patricii Episcopi et Confessoris}{17 mars}
\saintindex{Patrícium}
\saintindex{Patrick}
\oratio{Deus, qui ad prædicándam géntibus glóriam tuam beátum Patrícium Confessórem atque Pontíficem míttere dignátus es: ejus méritis et intercessióne concéde; ut, quæ nobis agénda prǽcipis, te miseránte adimplére possímus. \perdominumla}{Ô Dieu, qui avez daigné envoyer le bienheureux Patrick, Confesseur et Pontife, prêcher votre gloire aux nations, accordez-nous, par ses mérites et son intercession, de pouvoir accomplir, grâce à votre miséricorde, ce que vous nous ordonnez de faire. \perdominumfr}
\feast{0318}{S. Cyrilli Episcopi Hierosolymitani Confessoris et Ecclesiæ Doctoris}{18 mars}
\saintindex{Cyrille}
\saintindex{Cyrille}
\oratio{Da nobis, quǽsumus, omnípotens Deus, beato Cyríllo Pontífice intercedénte: te solum verum Deum, et quem misísti Jesum Christum ita cognóscere; ut inter oves, quæ vocem ejus áudiunt, perpétuo connumerári mereámur. \pereumdemla}{Donnez-nous, nous vous le demandons, Dieu tout-puissant, par l’intercession du bienheureux Pontife Cyrille, de vous connaître, vous, le seul vrai Dieu et celui que vous avez envoyé Jésus-Christ, de telle façon que nous méritions d’être comptés éternellement parmi les brebis qui écoutent toujours sa voix. \pereumdemfr}
\feast{0319}{S. Joseph Sponsi B.M.V. Confessoris}{19 mars}
\saintindex{}
\saintindex{}
\oratio{Sanctíssimæ Genetrícis tuæ sponsi, quǽsumus, Dómine, méritis adjuvémur; ut quod possibílitas nostra non óbtinet, ejus nobis intercessióne donétur.
$Qui vivis}{Faites, Seigneur, s’il Vous plaît, que nous soyons secourus par les mérites de l’époux de Votre très sainte Mère afin que ce que nous ne pouvons obtenir de nous-mêmes nous soit donné par son intercession.
$Qui vivis}
\feast{0319}{S. Joseph Sponsi B.M.V. Confessoris}{19 mars}
\saintindex{}
\saintindex{}
\oratio{@Sancti/03-19:Oratio}{}
\feast{0321}{S. Benedicti Abbatis}{21 mars}
\saintindex{Benedícti}
\saintindex{Benoît}
\feast{0323}{S. Turibius of Mogrovejo, bishop}{23 mars}
\saintindex{}
\saintindex{}
\feast{0324}{S. Gabrielis Archangeli}{24 mars}
\saintindex{}
\saintindex{}
\oratio{Deus, qui inter céteros Angelos, ad annuntiándum Incarnatiónis tuæ mystérium, Gabriélem Archángelum elegísti: concéde propítius; ut qui festum ejus celebrámus in terris, ipsíus patrocínium sentiámus in cælis:
$Qui vivis}{Ô Dieu, qui avez choisi l’Archange Gabriel entre tous les Anges, pour annoncer le mystère de Votre Incarnation, accordez-nous, dans Votre bonté, qu’après avoir célébré sa fête sur la terre, nous goûtions dans le ciel les effets de sa protection.
$Qui vivis}
\oratio{Deus, qui per Archángelum Gabriëlem Salvatórem mundi Sacratíssimæ Vírgini concipiéndum nuntiásti: da, ut eúndem et mente pura concipiámus, et férvido imitémur afféctu.
$Qui tecum}{}
\feast{0325}{}{25 mars}
\saintindex{}
\saintindex{}
\oratio{Deus, qui de beátæ Maríæ Vírginis útero Verbum tuum, Angelo nuntiánte, carnem suscípere voluísti: præsta supplícibus tuis; ut, qui vere eam Genetrícem Dei crédimus, ejus apud te intercessiónibus adjuvémur. \pereumdemla}{Ô Dieu, qui avez voulu que votre Verbe prît un corps humain à la parole de l’Ange dans le sein de la bienheureuse Vierge Marie ; accordez à ceux qui vous en supplient que, nous qui la croyons véritablement Mère de Dieu, nous soyons secourus auprès de vous grâce à son intercession. \pereumdemfr}
\feast{0327}{S. Joannis Damasceni Confessoris et Ecclesiæ Doctoris}{27 mars}
\saintindex{Joannes}
\saintindex{Joannes}
\oratio{Omnípotens sempitérne Deus, qui ad cultum sacrárum imáginum asseréndum, beátum Joánnem cælésti doctrína et admirábili spíritus fortitúdine imbuísti: concéde nobis ejus intercessióne et exémplo; ut, quorum cólimus imágines, virtútes imitémur et patrocínia sentiámus. \perdominumla}{Dieu tout-puissant et éternel, qui avez donné au bienheureux Jean une science toute céleste et une admirable force d’âme pour défendre le culte des saintes images, accordez-nous, par son intercession et à son exemple, d’imiter les vertus de ceux dont nous honorons les images, et de ressentir les effets de leur protection. \perdominumfr}
\feast{0328}{S. Joannis a Capistrano Confessoris}{28 mars}
\saintindex{}
\saintindex{}
\oratio{Deus, qui per beátum Joánnem fidéles tuos in virtúte sanctíssimi nóminis Jesu de crucis inimícis triumpháre fecísti: præsta, quǽsumus; ut, spirituálium hóstium, ejus intercessióne, superátis insídiis, corónam justítiæ a te accípere mereámur. \pereumdemla}{Ô Dieu, qui, par le bienheureux Jean, avez fait triompher vos fidèles des ennemis de la croix, grâce à la vertu du très saint nom de Jésus, faites nous vous en supplions, qu’ayant surmonté par son intercession les embûches de nos ennemis spirituels, nous méritions de recevoir de vous la couronne de justice. \pereumdemfr}
\feast{0401}{}{1 avril}
\saintindex{}
\saintindex{}
\oratio{Deus omípotens, cujus sérvitus summa et plena est felícitas: * concéde nobis, intercedénte beáto Hugóne, Abbáte, § in sanctitáte et justítia sine timóre tibi servíre. \perdominumla}{}
\feast{0402}{S. Francisci de Paula Confessoris}{2 avril}
\saintindex{}
\saintindex{}
\oratio{Deus, humílium celsitúdo, qui beátum Francíscum Confessórem Sanctórum tuórum glória sublimásti: tríbue, quǽsumus; ut, ejus méritis et imitatióne, promíssa humílibus prǽmia felíciter consequámur. \perdominumla.}{Ô Dieu, grandeur des humbles, qui avez élevé le bienheureux François, Confesseur, à la gloire de vos Saints, accordez, nous vous en prions, que, par ses mérites et par l’imitation de ses vertus, nous ayons le bonheur d’obtenir les récompenses promises à ceux qui sont humbles. \perdominumfr.}
\feast{0404}{S. Isidori Episcopi Confessoris et Ecclesiæ Doctoris}{4 avril}
\saintindex{Isidórum}
\saintindex{Isidóre}
\oratio{}{Ô Dieu qui avez fait à votre peuple la grâce d’avoir le bienheureux Isidore, pour ministre du salut éternel, faites, nous vous en prions, que nous méritions d’avoir pour intercesseur dans les cieux celui qui nous a donné sur terre la doctrine de vie. \perdominumfr}
\feast{0405}{S. Vincentii Ferrerii Confessoris}{5 avril}
\saintindex{}
\saintindex{}
\oratio{Deus, qui Ecclésiam tuam beáti Vincéntii Confessóris tui méritis et prædicatióne illustráre dignátus es: concéde nobis fámulis tuis; ut et ipsíus instruámur exémplis et ab ómnibus ejus patrocínio liberémur advérsis. \perdominumla}{Dieu, qui avez daigné illustrer Votre Église par Les mérites et la prédication du bienheureux Vincent, Votre confesseur, accordez à nous, Vos familiers, que nous soyons instruits par ses exemples et libérés de toutes adversités par son patronage. \perdominumfr}
\feast{0411}{S. Leonis I Papæ Confessoris et Ecclesiæ Doctoris}{11 avril}
\saintindex{Leónem}
\saintindex{Léon}
\oratio{@Commune/C4-1}{Nous vous supplions, Seigneur, d’exaucer les prières que nous vous adressons en la solennité du bienheureux Léon, votre Confesseur et Pontife, et de nous accorder, grâce aux mérites et à l’intercession de celui qui vous a si dignement servi, le pardon de tous nos péchés. \perdominumfr}
\feast{0413}{S. Hermenegildi Martyris}{13 avril}
\saintindex{}
\saintindex{}
\oratio{Deus, qui beátum Hermenegíldum Mártyrem tuum cælésti regno terrénum postpónere docuísti: da quǽsumus nobis; ejus exémplo cadúca despícere, atque ætérna sectári. \perdominumla}{Ô Dieu, qui avez appris au Bienheureux Herménégilde à faire passer le royaume de ce monde après celui des cieux, accordez-nous, s’il Vous plaît, de mépriser, à son exemple, les biens périssables et de poursuivre les éternels. \perdominumfr}
\feast{0414}{S. Justini Martyris}{14 avril}
\saintindex{}
\saintindex{}
\oratio{Deus, qui per stultítiam crucis eminéntem Jesu Christi sciéntiam beátum Justínum Mártyrem mirabíliter docuísti: ejus nobis intercessióne concéde; ut, errórum circumventióne depúlsa, fídei firmitátem consequámur. \pereumdemla}{Ô Dieu qui, par la folie de la croix, avez merveilleusement enseigné au bienheureux Martyr Justin la suprême science de Jésus-Christ, accordez-nous par Son intercession que, repoussant tout piège d’erreur, nous obtenions la fermeté de la foi. \pereumdemfr}
\feast{0414}{Ss. Tiburtii, Valeriani, et Maximi Martyrum}{14 avril}
\saintindex{}
\saintindex{}
\oratio{Præsta, quǽsumus, omnípotens Deus: ut, qui sanctórum Mártyrum tuórum Tiburtii, Valeriáni et Maximi solemnia colimus; eórum étiam virtútes imitemur. \perdominumla}{Faites, nous Vous en supplions, Seigneur, que nous qui célébrons la fête de Vos saints Martyrs Tiburce, Valérien et Maxime, nous imitions aussi leurs vertus. \perdominumfr}
\feast{0417}{S. Aniceti Papæ et Martyris}{17 avril}
\saintindex{Anicéti}
\saintindex{Anicet}
\feast{0421}{S. Anselmi Episcopi Confessoris et Ecclesiæ Doctoris}{21 avril}
\saintindex{Ansélmum}
\saintindex{Anselme}
\feast{0422}{SS. Soteris et Caji Summorum Pontificum et Martyrum}{22 avril}
\saintindex{Sotéris et Caji}
\saintindex{Soter et Caius}
\feast{0423}{S. Georgii Martyris}{23 avril}
\saintindex{}
\saintindex{}
\oratio{Deus, qui nos beáti Geórgii Mártyris tui méritis et intercessióne lætíficas: concéde propítius; ut, qui tua per eum benefícia póscimus, dono tuæ grátiæ consequámur. \perdominumla}{Ô Dieu qui nous réjouissez par les mérites et l’intercession du Bienheureux Georges, Votre Martyr, accordez-nous, dans Votre bonté, que, Vous demandant par lui Vos bienfaits, nous les obtenions par le don de Votre grâce. \perdominumfr}
\feast{0424}{S. Fidelis de Sigmaringa Martyris}{24 avril}
\saintindex{}
\saintindex{}
\oratio{Deus, qui beátum Fidélem, seráphico spíritus ardóre succénsum, in veræ fídei propagatióne martýrii palma et gloriósis miráculis decoráre dignátus es: ejus, quǽsumus, méritis et intercessióne, ita nos per grátiam tuam in fide et caritáte confírma; ut in servítio tuo fidéles usque ad mortem inveníri mereámur. \perdominumla}{Ô Dieu qui, après avoir embrasé le Bienheureux Fidèle de l’ardeur de l’esprit séraphique, avez daigné l’illustrer par la palme du martyre et par de glorieux miracles, nous Vous demandons, par ses mérites et son intercession, de si bien nous confirmer, par Votre grâce, dans la foi et la charité, que nous méritions d’être trouvés fidèles dans Votre service jusqu’à la mort. \perdominumfr}
\feast{0425}{S. Marci Evangelistæ}{25 avril}
\saintindex{}
\saintindex{}
\oratio{Deus, qui beátum Marcum Evangelístam tuum evangélicæ prædicatiónis grátia sublimásti: tríbue, quǽsumus; ejus nos semper et eruditióne profícere et oratióne deféndi. \perdominumla}{Ô Dieu, qui avez élevé si haut le bienheureux Marc, Votre Evangéliste, par la grâce de la prédication évangélique, accordez-nous, s’il Vous plaît, de profiter toujours de sa doctrine et d’être toujours défendus par sa prière. \perdominumfr}
\feast{0426}{SS. Cleti et Marcellini Summorum Pontificum et Martyrum}{26 avril}
\saintindex{Cleti et Marcellíni}
\saintindex{Clet et Marcellin}
\feast{0427}{S. Petri Canisii Confessoris et Ecclesiæ Doctoris}{27 avril}
\saintindex{Petrum}
\saintindex{Pierre}
\oratio{Deus, qui ad tuéndam cathólicam fidem beátum Petrum Confessórem tuum virtúte et doctrína roborásti: concéde propítius; ut ejus exémplis et mónitis errántes ad salútem resipíscant, et fidéles in veritátis confessióne persevérent. \perdominumla}{Ô Dieu, qui pour la protection de la foi catholique, avez affermi dans la vertu et la doctrine le Bienheureux Pierre, Votre Confesseur, accordez, dans Votre clémence, qu’instruits par ses exemples et ses enseignements, les égarés reviennent à la voie du salut et que les fidèles persévèrent dans la confession de la vérité. \perdominumfr}
\feast{0428}{S. Pauli a Cruce Confessoris}{28 avril}
\saintindex{}
\saintindex{}
\oratio{Dómine Jesu Christe, qui, ad mystérium crucis prædicándum, sanctum Paulum singulári caritáte donásti, et per eum novam in Ecclésia famíliam floréscere voluísti: ipsíus nobis intercessióne concéde; ut, passiónem tuam júgiter recoléntes in terris, ejúsdem fructum cónsequi mereámur in cælis:
$Qui vivis}{Seigneur Jésus-Christ qui, pour la prédication du mystère de la Croix, avez gratifié saint Paul d’une charité sans pareille et qui avez voulu, par lui, faire fleurir une nouvelle famille dans l’Eglise, accordez-nous, par son intercession, qu’en nous remémorant continuellement Votre passion sur Terre, nous méritions d’en recueillir les fruits au ciel,
$Qui vivis}
\feast{0428}{S. Peter Chanel, priest and martyr}{28 avril}
\saintindex{Petri}
\saintindex{Pierre}
\feast{0428}{S. Vitalis Martyris}{28 avril}
\saintindex{Vitális}
\saintindex{Vital}
\feast{0429}{S. Petri Martyris}{29 avril}
\saintindex{}
\saintindex{}
\oratio{Præsta, quǽsumus, omnípotens Deus: ut beáti Petri Mártyris tui fidem cóngrua devotióne sectémur; qui, pro ejúsdem fídei dilatatióne, martýrii palmam méruit obtinére. \perdominumla}{Faites, nous Vous le demandons, ô Dieu tout-puissant, que nous imitions avec le zèle qui convient, la foi du bienheureux Pierre, Votre Martyr qui, pour l’expansion de cette même foi, a mérité d’obtenir la palme du martyre. \perdominumfr}
\feast{0430}{S. Catharinæ Senensis Virginis}{30 avril}
\saintindex{}
\saintindex{}
\oratio{Da, quǽsumus, omnípotens Deus: ut, qui beátæ Catharínæ Vírginis tuæ natalícia cólimus; et ánnua solemnitáte lætémur, et tantæ virtútis proficiámus exémplo. \perdominumla}{Accordez-nous, s’il Vous plaît, Dieu tout-puissant, qu’honorant la naissance au ciel de Votre bienheureuse Vierge Catherine, nous nous réjouissions de cette solennité annuelle et progressions par l’exemple d’une si grande vertu. \perdominumfr}
\feast{0501}{Ss. Philippi et Jacobi Apostolorum}{1 mai}
\saintindex{}
\saintindex{}
\oratio{Deus, qui nos ánnua Apostolórum tuórum Philíppi et Jacóbi solemnitáte lætíficas: præsta, quǽsumus; ut, quorum gaudémus méritis, instruámur exémplis. \perdominumla}{Dieu qui nous réjouissez en la solennité annuelle de Vos Apôtres Philippe et Jacques ; faites, nous Vous en prions, que nous soyons instruits aux exemples de ceux dont les mérites nous remplissent d’allégresse. \perdominumfr}
\feast{0501}{S. Joseph Sponsi B.M.V. Confessoris}{1 mai}
\saintindex{}
\saintindex{}
\oratio{Deus, qui ineffábili providéntia beátum Joseph sanctíssimæ Genitrícis tuæ sponsum eligere dignátus es: præsta, quǽsumus; ut, quem protectórem venerámur in terris, intercessórem habere mereámur in cælis:
$Qui vivis}{Dieu créateur de toutes choses, qui avez imposé au genre humain la loi du travail : accordez-nous favorablement ; grâce à l’exemple et au patronage de saint Joseph, d’accomplir parfaitement le travail que vous nous fixez, et d’obtenir les récompenses que vous nous promettez.
$Qui vivis}
\feast{0501}{S. Joseph Opificis}{1 mai}
\saintindex{}
\saintindex{}
\oratio{Rerum cónditor Deus, qui legem labóris humáno géneri statuísti: concéde propítius; ut, sancti Joseph exémplo et patrocínio, ópera perficiámus quæ prǽcipis, et prǽmia consequámur quæ promíttis. \perdominumla}{Créateur du monde, ô Dieu, qui avez imposé la loi du travail au genre humain, accordez-nous avec bonté que, à l’exemple et sous le patronage de saint Joseph, nous accomplissions les travaux que vous ordonnez et nous obtenions les récompenses que vous promettez. \perdominumfr}
\feast{0502}{S. Athanasii Episcopi Confessoris et Ecclesiæ Doctoris}{2 mai}
\saintindex{Athanásii}
\saintindex{Athanase}
\oratio{@Commune/C4-1}{}
\feast{0503}{}{3 mai}
\saintindex{}
\saintindex{}
\oratio{Deus, qui in præclára salutíferæ Crucis Inventióne, passiónis tuæ mirácula suscitásti: concéde; ut vitális ligni prétio, ætérnæ vitæ suffrágia consequámur:
$Qui vivis}{Dieu, qui lors de la glorieuse Découverte de la Croix, instrument de notre salut, avez renouvelé les miracles de votre passion : faites qu’au prix de cet arbre de vie, nous méritions d’obtenir la vie éternelle.
$Qui vivis}
\feast{0503}{Ss. Alexandri I Papæ, Eventii et Theoduli Mm. ac Juvenalis Ep. et Conf.}{3 mai}
\saintindex{}
\saintindex{}
\oratio{Præsta, quǽsumus, omnípotens Deus: ut, qui sanctórum tuórum Alexándri, Evéntii, Theodúli atque Juvenális natalícia cólimus; a cunctis malis imminéntibus, eórum intercessiónibus, liberémur.  \perdominumla}{Accordez-nous, s’il vous plaît, ô Dieu tout-puissant, que, célébrant la naissance au ciel de vos saints Alexandre, Eventius, Théodule et Juvénal, nous soyons délivrés, grâce à leur intercession, de tous les maux qui nous menacent. \perdominumfr}
\feast{0503}{}{3 mai}
\saintindex{}
\saintindex{}
\feast{0504}{S. Monicæ Viduæ}{4 mai}
\saintindex{Mónicæ}
\saintindex{Monique}
\oratio{Deus, mæréntium consolátor et in te sperántium salus, qui beátæ Mónicæ pias lácrimas in conversióne fílii sui Augustíni misericórditer suscepísti: da nobis utriúsque intervéntu; peccáta nostra deploráre, et grátiæ tuæ indulgéntiam inveníre. \perdominumla}{Ô Dieu, consolateur des affligés et salut de ceux qui espèrent en Vous, Vous avez miséricordieusement exaucé les pieuses larmes versées par la bienheureuse Monique pour la conversion de son fils Augustin ; donnez-nous, par l'intercession de tous deux, de pleurer nos péchés et de trouver l'indulgence de Votre grâce. \perdominumfr}
\feast{0505}{S. Pii V Papæ et Confessoris}{5 mai}
\saintindex{Pium}
\saintindex{Pie}
\oratio{Deus, qui, ad conteréndos Ecclésiæ tuæ hostes et ad divínum cultum reparándum, beátum Pium Pontíficem Máximum elígere dignátus es: fac nos ipsíus deféndi præsídiis et ita tuis inhærére obséquiis; ut, ómnium hóstium superátis insídiis, perpétua pace lætémur. \perdominumla}{Ô Dieu qui, pour écraser les ennemis de Votre Eglise et restaurer le culte divin, avez daigné choisir le bienheureux Pie pour Souverain Pontife, faites que, défendus par son secours pour demeurer très attachés à Votre service, nous triomphions des pièges de tous nos ennemis et jouissions de l'éternelle paix. \perdominumfr}
\feast{0506}{S. Joannis Apostoli ante Portam Latinam}{6 mai}
\saintindex{}
\saintindex{}
\oratio{Deus, qui conspicis, quia nos undique mala nostra perturbant: præsta, quǽsumus; ut beáti Joánnis Apóstoli tui et Evangelistæ intercéssio gloriósa nos prótegat. \perdominumla}{Ô Dieu qui voyez comment nous sommes troublés par les maux qui nous arrivent de partout, faites, nous vous le demandons, que nous soyons protégés par la glorieuse intercession du bienheureux Jean, votre Apôtre et Évangéliste.  \perdominumfr}
\feast{0507}{S. Stanislai Episcopi et Martyris}{7 mai}
\saintindex{Stanisláus}
\saintindex{Stanislas}
\oratio{Deus, pro cujus honóre gloriósus Póntifex Stanisláus gládiis impiórum occúbuit: præsta, quǽsumus; ut omnes, qui ejus implórant auxílium, petitiónis suæ salutárem consequántur efféctum. \perdominumla}{Ô Dieu, c'est pour Votre honneur que le glorieux Pontife Stanislas a succombé sous le glaive des impies ; accordez-nous, s'il Vous plaît, que tous ceux qui implorent son aide, obtiennent pour leur salut l'effet de leur demande. \perdominumfr}
\feast{0508}{In Apparitione S. Michaëlis Archangeli}{8 mai}
\saintindex{}
\saintindex{}
\oratio{Deus, qui, miro órdine, Angelórum ministéria hominúmque dispénsas: concéde propítius; ut, a quibus tibi ministrántibus in cælo semper assístitur, ab his in terra vita nostra muniátur.  \perdominumla}{Ô Dieu, qui distribuez selon un ordre admirable les ministères des Anges et des hommes, accordez-nous miséricordieusement que ceux qui, dans le ciel, Vous entourent d'un continuel service, soient la protection de notre vie sur terre. \perdominumfr}
\feast{0509}{S. Gregorii Nazianzeni Episcopi Confessoris et Ecclesiæ Doctoris}{9 mai}
\saintindex{Gregórium}
\saintindex{Grégoire }
\feast{0509}{}{9 mai}
\saintindex{}
\saintindex{}
\feast{0510}{S. Antonini Episcopi et Confessoris}{10 mai}
\saintindex{}
\saintindex{}
\oratio{Sancti Antoníni, Dómine, Confessóris tui atque Pontíficis méritis adjuvémur: ut, sicut te in illo mirábilem prædicámus, ita in nos misericórdem fuísse gloriémur. \perdominumla}{Que les mérites de saint Antonin, votre Confesseur et Pontife, nous soient en aide, ô Seigneur ; et comme nous vous proclamons admirable dans votre serviteur, faites que nous puissions aussi nous glorifier de votre miséricorde à notre égard. \perdominumfr}
\feast{0510}{Ss. Gordiani et Epimachi Martyrum}{10 mai}
\saintindex{}
\saintindex{}
\oratio{Da, quǽsumus, omnípotens Deus: ut, qui beatórum Mártyrum tuórum Gordiáni et Epímachi solémnia cólimus, eórum apud te intercessiónibus adjuvémur. \perdominumla}{Faites, s’il vous plaît, Dieu tout-puissant, qu’en célébrant la fête de vos bienheureux martyrs Gordien et Épimaque, nous soyons secourus par leur intercession auprès de vous.  \perdominumfr}
\feast{0511}{}{11 mai}
\saintindex{}
\saintindex{}
\feast{0512}{Ss. Nerei, Achillei et Domitillæ Virg. atque Pancratii Martyrum}{12 mai}
\saintindex{}
\saintindex{}
\oratio{Semper nos, Dómine, Mártyrum tuórum Nérei, Achíllei, Domitíllæ atque Pancrátii fóveat, quǽsumus, beáta solémnitas: et tuo dignos reddat obséquio. \perdominumla}{Que toujours, nous Vous en supplions, Seigneur, l’heureuse fête de Vos Martyrs Nérée, Achillée, Domitille et Pancrace nous soutienne et nous rende dignes de Votre service. \perdominumfr}
\feast{0513}{S. Roberti Bellarmino Episcopi Confessoris et Ecclesiæ Doctoris}{13 mai}
\saintindex{Robérte}
\saintindex{Robert}
\oratio{Deus, qui ad errórum insídias repelléndas et apostólicæ Sedis jura propugnánda, beátum Robértum Pontíficem tuum atque Doctórem mira eruditióne et virtúte decorásti: ejus méritis et intercessióne concéde; ut nos in veritátis amóre crescámus et errántium corda ad Ecclésiæ tuæ rédeant unitátem. \perdominumla}{Ô Dieu, pour repousser les pièges de l’erreur et défendre les droits du Siège apostolique, Vous avez doté le bienheureux Robert, Votre Pontife et Docteur, d’une science et d’une force admirable ; faites, par ses mérites et son intercession, que nous grandissions dans l’amour de la vérité et que les cœurs des égarés reviennent à l’unité de Votre Église. \perdominumfr}
\feast{0513}{}{13 mai}
\saintindex{}
\saintindex{}
\feast{0514}{S. Bonifatii Martyris}{14 mai}
\saintindex{}
\saintindex{}
\oratio{Da, quǽsumus, omnípotens Deus: ut qui beáti Bonifátii Mártyris tui solémnia cólimus, ejus apud te intercessiónibus adjuvémur. \perdominumla}{Faites, nous Vous en prions, Dieu tout-puissant, que nous, qui célébrons la fête de Votre bienheureux Martyr Boniface, nous recevions l'assistance de son intercession. \perdominumfr}
\feast{0515}{S. Joannis Baptistæ de la Salle Confessoris}{15 mai}
\saintindex{}
\saintindex{}
\oratio{Deus, qui, ad christiánam páuperum eruditiónem et ad juvéntam in via veritátis firmándam, sanctum Joánnem Baptístam Confessórem excitásti, et novam per eum in Ecclésia famíliam collegísti: concéde propítius; ut ejus intercessióne et exémplo, stúdio glóriæ tuæ in animárum salúte fervéntes, ejus in cælis corónæ partícipes fíeri valeámus. \perdominumla}{Ô Dieu qui, pour l’instruction chrétienne des pauvres et pour la confirmation de la jeunesse dans la voie de la vérité, avez suscité Votre confesseur Jean-Baptiste et avez par lui rassemblé dans l’Église une nouvelle famille, accordez-nous dans Votre bonté, qu’à son exemple et par son intercession, brûlants de zèle pour procurer Votre gloire au moyen du salut des âmes, nous puissions dans les cieux partager sa récompense. \perdominumfr}
\feast{0516}{S. Ubaldi Episcopi et Confessoris}{16 mai}
\saintindex{}
\saintindex{}
\oratio{Auxílium tuum nobis, Dómine, quǽsumus, placátus impénde: et, intercessióne beáti Ubáldi Confessóris tui atque Pontíficis, contra omnes diáboli nequítias déxteram super nos tuæ propitiatiónis exténde. \perdominumla}{Laissez-Vous fléchir, Seigneur : que l’intercession du bienheureux Ubald, Votre Confesseur et Pontife, nous obtienne Votre secours ; étendez sur nous Votre main miséricordieuse pour nous défendre contre toutes les perfidies du démon. \perdominumfr}
\feast{0516}{}{16 mai}
\saintindex{}
\saintindex{}
\oratio{Deus, qui ob invictum beati Joannis sacramentale silentium, nova Ecclesiam tuam martyrii corona decorasti: da nobis, ejus intercessione et exemplo, linguam caute custodire; ac omnia potius mala, quam animæ detrimentum, in hoc sæculo tolerare. \perdominumla.}{Deus, qui ob invictum beati Joannis sacramentale silentium, nova Ecclesiam tuam martyrii corona decorasti: da nobis, ejus intercessione et exemplo, linguam caute custodire; ac omnia potius mala, quam animæ detrimentum, in hoc sæculo tolerare. \perdominumfr.}
\feast{0517}{S. Paschalis Baylon Confessoris}{17 mai}
\saintindex{}
\saintindex{}
\oratio{Deus, qui beátum Paschálem Confessórem tuum mirífica erga Córporis et Sánguinis tui sacra mystéria dilectióne decorásti: concéde propítius; ut, quam ille ex hoc divíno convívio spíritus percépit pinguédinem, eándem et nos percípere mereámur:
$Qui vivis}{Ô Dieu, qui avez orné l’âme du bienheureux Pascal, votre Confesseur, d’un admirable et tendre amour pour les mystères sacrés de votre corps et de votre sang, accordez-nous, dans votre bonté, que nous méritions de retirer de ce banquet divin la même abondance de grâces qu’il y a trouvée.
$Qui vivis}
\feast{0518}{S. Venantii Martyris}{18 mai}
\saintindex{}
\saintindex{}
\oratio{Deus, qui hunc diem beáti Venántii Mártyris tui triúmpho consecrásti: exáudi preces pópuli tui et præsta: ut, qui ejus mérita venerámur, fídei constántiam imitémur.  \perdominumla}{Ô Dieu, qui avez consacré ce jour par le triomphe du bienheureux Venant, Votre Martyr, exaucez les prières de Votre peuple et faites qu’honorant ses mérites, nous imitions la constance de sa foi. \perdominumfr}
\feast{0518}{S. John I, pope and martyr}{18 mai}
\saintindex{}
\saintindex{}
\feast{0519}{S. Petri Celestini Papæ et Confessoris}{19 mai}
\saintindex{}
\saintindex{}
\oratio{Deus, qui beátum Petrum Cælestínum ad summi pontificátus ápicem sublimásti, quique illum humilitáti postpónere docuísti: concéde propítius; ut ejus exémplo cuncta mundi despícere, et ad promíssa humílibus prǽmia perveníre felíciter mereámur. \perdominumla}{Dieu, qui avez élevé le bienheureux Pierre Célestin à l’éminente dignité du souverain pontificat et qui lui avez appris à mettre l’humilité au-dessus de cette élévation, accordez-nous, dans Votre bonté, la grâce de mépriser, à son exemple, tous les biens de ce monde et de parvenir heureusement à la possession des récompenses promises à ceux qui sont humbles. \perdominumfr}
\oratio{@Commune/C4-1:Oratio:s/N./Petri Cœlestíni/}{}
\feast{0519}{S. Pudentianæ Virginis}{19 mai}
\saintindex{Pudentiánæ}
\saintindex{Pudentienne}
\feast{0520}{S. Bernardini Senensis Confessoris}{20 mai}
\saintindex{}
\saintindex{}
\oratio{Dómine Jesu, qui beáto Bernardíno Confessóri tuo exímium sancti nóminis tui amórem tribuísti: ejus, quǽsumus, méritis et intercessióne, spíritum nobis tuæ dilectiónis benígnus infúnde:
$Qui vivis}{Seigneur Jésus, qui avez accordé au bienheureux Bernardin, Votre Confesseur, un très ardent amour pour Votre saint Nom, nous Vous supplions, par ses mérites et son intercession, de daigner, dans Votre bonté, répandre en nous l’esprit de Votre charité.
$Qui vivis}
\feast{0521}{S. Christopher Magallanes and companions, martyrs}{21 mai}
\saintindex{}
\saintindex{}
\feast{0522}{S. Ritæ de Cascia}{22 mai}
\saintindex{}
\saintindex{}
\oratio{Deus, qui sanctæ Ritæ tantam grátiam conférre dignátus es, ut inimícos dilígeret et in corde ac fronte caritátis et passiónis tuæ signa portáret: da nobis, quǽsumus, ejus intercessióne et méritis; inimícis nostris sic párcere et passiónis tuæ dolóres contemplári, ut promíssa mítibus ac lugéntibus prǽmia consequámur:
$Qui vivis
}{}
\feast{0523}{S. Joannis Baptistæ de Rossi Confessoris}{23 mai}
\saintindex{}
\saintindex{}
\oratio{Deus, qui sanctum Joánnem Baptistam Confessorem tuum in evangelizandis pauperibus caritate et patientia decorasti: concede quǽsumus; ut cujus pia merita veneramur, virtutum quoque imitemur exempla. \perdominumla}{Dieu, qui avez orné saint Jean Baptiste, Votre Confesseur, de la charité et de la patience dans l’évangélisation des pauvres, accordez-nous, nous Vous le demandons, de vénérer ses pieux mérites et d’imiter les exemples de ses vertus. \perdominumfr}
\feast{0523}{S. John Baptist Rossi Confessor}{23 mai}
\saintindex{}
\saintindex{}
\oratio{Deus, qui sanctum Joánnem Baptistam Confessorem tuum in evangelizandis pauperibus caritate et patientia decorasti: concede qusesumus; ut cujus pia merita veneramur, virtutum quoque imitemur exempla. \perdominumla}{Dieu, qui avez orné saint Jean Baptiste, Votre Confesseur, de la charité et de la patience dans l’évangélisation des pauvres, accordez-nous, nous Vous le demandons, de vénérer ses pieux mérites et d’imiter les exemples de ses vertus. \perdominumfr}
\feast{0525}{S. Gregorii VII Papæ et Confessoris}{25 mai}
\saintindex{}
\saintindex{}
\oratio{Deus, in te sperántium fortitúdo, qui beátum Gregórium Confessórem tuum atque Pontíficem, pro tuénda Ecclésiæ libertáte, virtúte constántiæ roborásti: da nobis, ejus exémplo et intercessióne, ómnia adversántia fórtiter superáre. \perdominumla}{Ô Dieu, force de ceux qui espèrent en Vous, Vous qui avez fortifié le bienheureux Grégoire, Votre Confesseur et Pontife par la vertu de constance dans la défense de la liberté de l'Eglise, accordez-nous, par son exemple et son intercession, de triompher courageusement de toute adversité. \perdominumfr}
\feast{0525}{S. Urbanis Papæ et Martyris}{25 mai}
\saintindex{Urbáno}
\saintindex{Urbain}
\oratio{Da, quǽsumus, omnípotens Deus: ut qui beáti Urbáni Mártyris tui atque Pontíficis solémnia cólimus, ejus apud te intercessiónibus adjuvémur. \perdominumla}{Faites, nous Vous en prions, Dieu tout-puissant, que, nous qui célébrons la fête de Votre bienheureux Martyr et Pontife Urbain, nous recevions l’assistance de son intercession. \perdominumfr}
\oratio{@Commune/C2b:Oratio pro commemoratio
(sed rubrica 196)
@Commune/C2b}{Ô Dieu, sur le roc inébranlable des apôtres, vous avez fondé votre Eglise qui n’a plus a redouter la puissance des enfers ; faites qu’à la prière du bienheureux Urbain, votre Martyr et Souverain Pontife, elle reste ferme dans votre vérité et sous votre garde, elle soit en sûreté pour toujours. \perdominumfr
(sed rubrica 196 loco horum versum dicitur)
@Commune/C2b}
\feast{0526}{S. Philippi Neri Confessoris}{26 mai}
\saintindex{}
\saintindex{}
\oratio{Deus, qui beátum Philíppum Confessórem tuum Sanctórum tuórum glória sublimásti: concéde propítius; ut, cujus solemnitáte lætámur, ejus virtútum proficiámus exémplo. \perdominumla}{Dieu, qui avez élevé le bienheureux Philippe, Votre Confesseur, à la gloire de Vos Saints, accordez-nous, dans Votre miséricorde que, célébrant avec joie cette solennité, nous mettions à profit l’exemple de ses vertus. \perdominumfr}
\feast{0526}{S. Eleutherii Papæ et Martyris}{26 mai}
\saintindex{Eleuthérii}
\saintindex{Eleuthère}
\feast{0527}{S. Bedæ Venerabilis Confessoris et Ecclesiæ Doctoris}{27 mai}
\saintindex{Bedæ}
\saintindex{Bède}
\oratio{Deus, qui Ecclésiam tuam beáti Bedæ Confessóris tui atque Doctóris eruditióne claríficas: concéde propítius fámulis tuis; ejus semper illustrári sapiéntia et méritis adjuvári. \perdominumla}{Dieu, qui illustrez Votre Église par la science du bienheureux Bède, Votre Docteur et Confesseur, accordez dans Votre bonté à Vos serviteurs d’être toujours éclairés de sa sagesse et aidés de ses mérites. \perdominumfr}
\feast{0527}{S. Joannis Papæ et Martyris}{27 mai}
\saintindex{Joannis}
\saintindex{Jean}
\feast{0528}{S. Augustini Episcopi et Confessoris}{28 mai}
\saintindex{}
\saintindex{}
\oratio{Deus, qui Anglórum gentes, prædicatióne et miráculis beáti Augustíni Confessóris tui atque Pontíficis, veræ fídei luce illustráre dignátus es: concéde; ut, ipso interveniénte, errántium corda ad veritátis tuæ rédeant unitátem, et nos in tua simus voluntáte concórdes. \perdominumla}{Ô Dieu, qui, par la prédication et les miracles du bienheureux Augustin, Votre Confesseur et Pontife, avez daigné éclairer de la lumière de la vraie foi la nation anglaise, faites que, par son intercession, les cœurs égarés reviennent à l’unité de Votre vérité, et que nous soyons tous unis de cœur en Votre volonté. \perdominumfr}
\feast{0529}{S. Mariæ Magdalenæ de Pazzis Virginis}{29 mai}
\saintindex{}
\saintindex{}
\oratio{Deus, virginitátis amátor, qui beátam Maríam Magdalénam Vírginem, tuo amóre succénsam, cæléstibus donis decorásti: da; ut, quam festíva celebritáte venerámur, puritáte et caritáte imitémur. \perdominumla}{Dieu, qui aimez la virginité et qui avez orné des dons célestes Marie-Madeleine, cette vierge embrasée de Votre amour, donnez-nous d’imiter, dans sa pureté et sa charité, celle que nous vénérons en célébrant sa fête. \perdominumfr}
\feast{0530}{S. Felicis I Papæ et Martyris}{30 mai}
\saintindex{Felícis}
\saintindex{Félix}
\feast{0531}{}{31 mai}
\saintindex{}
\saintindex{}
\oratio{Concéde nobis, quǽsumus, Dómine: ut, qui solemnitátem beatæ Maríæ Vírginis Regínæ nostræ celebrámus; ejus muníti præsídio, pacem in præsénti et glóriam in futúro cónsequi mereámur. \perdominumla}{Nous célébrons la solennité de la Bienheureuse Vierge Marie, notre Reine : accordez-nous, Seigneur, de pouvoir, avec l'appui de son secours, obtenir la paix dans le présent et la gloire dans l'avenir. \perdominumfr}
\feast{0531}{S. Petronillæ Virginis}{31 mai}
\saintindex{Petroníllæ}
\saintindex{Petronille}
\feast{0601}{S. Angelæ Mericiæ Virginis}{1 juin}
\saintindex{}
\saintindex{}
\oratio{Deus, qui novum per beátam Angelam sacrárum Vírginum collégium in Ecclésia tua floréscere voluísti: da nobis, ejus intercessióne, angélicis móribus vívere; ut, terrénis ómnibus abdicátis, gáudiis pérfrui mereámur ætérnis. \perdominumla}{Ô Dieu, qui, par la bienheureuse Angèle, avez voulu qu’une nouvelle société de vierges saintes fleurît dans votre Eglise, faites-nous, par son intercession, la grâce de mener une vie angélique, afin que, renonçant à toutes les choses de la terre, nous méritions de jouir des joies éternelles. \perdominumfr}
\feast{0602}{Ss. Marcellini, Petri, atque Erasmi, Episcopi, Martyrum}{2 juin}
\saintindex{}
\saintindex{}
\oratio{Deus, qui nos ánnua beatórum Mártyrum tuórum Marcellíni, Petri atque Erásmi solemnitáte lætíficas: præsta, quǽsumus; ut, quorum gaudémus méritis, accendámur exémplis. \perdominumla}{Ô Dieu, qui nous faites trouver un sujet de joie dans la solennité de Vos bienheureux Martyrs Marcellin, Pierre et Erasme, accordez-nous, s’il Vous plaît, la grâce d’être enflammés d’ardeur par les exemples de ceux dont les mérites nous réjouissent.  \perdominumfr}
\feast{0603}{Ss. Caroli Lwanga, Matthiæ Murumba et Sociorum Martyrum}{3 juin}
\saintindex{}
\saintindex{}
\oratio{Omnípotens et miséricors Deus, qui univérsum mundum lúmine veritátis illústras, et per gloriósos Mártyres Matthíam, Cárolum et Sócios fídei prodígia in Africæ regiónibus renováre dignátus es: largíre supplícibus tuis; ut, ipsis intercedéntibus, in nobis grátiæ tuæ dona custódias, et inter gentes omnes mérito et número pópulus tibi sérviens augeátur. \perdominumla}{}
\feast{0604}{S. Francisci Caracciolo Confessoris}{4 juin}
\saintindex{}
\saintindex{}
\oratio{Deus, qui beátum Francíscum, novi órdinis institutórem, orándi stúdio et pœniténtiæ amóre decorásti: da fámulis tuis in ejus imitatióne ita profícere; ut, semper orántes et corpus in servitútem redigéntes, ad cæléstem glóriam perveníre mereántur. \perdominumla}{Ô Dieu, qui avez suscité le bienheureux François pour être le fondateur d’un nouvel Ordre et l’avez admirablement doué de zèle pour la prière et d’amour de la pénitence, accordez à Vos serviteurs de profiter si bien de ses exemples que, s’appliquant toujours à prier et à réduire leur corps en servitude, ils méritent de parvenir à la gloire céleste. \perdominumfr}
\feast{0605}{S. Bonifatii Episcopi et Martyris}{5 juin}
\saintindex{}
\saintindex{}
\oratio{Deus, qui multitúdinem populórum, beáti Bonifátii Mártyris tui atque Pontíficis zelo, ad agnitiónem tui nóminis vocáre dignátus es: concéde propítius; ut, cujus solémnia cólimus, étiam patrocínia sentiámus. \perdominumla}{Ô Dieu, qui avez daigné appeler une multitude de peuples à la connaissance de Votre nom par le zèle du bienheureux Boniface, Votre Martyr et Pontife, accordez-nous, dans Votre bonté que, célébrant sa fête, nous ressentions les effets de sa protection. \perdominumfr}
\feast{0606}{S. Norberti Episcopi et Confessoris}{6 juin}
\saintindex{}
\saintindex{}
\oratio{Deus, qui beátum Norbértum Confessórem tuum atque Pontíficem verbi tui præcónem exímium effecísti, et per eum Ecclésiam tuam nova prole fecundásti: præsta, quǽsumus; ut, ejúsdem suffragántibus méritis, quod ore simul et ópere dócuit, te adjuvánte, exercére valeámus. \perdominumla.}{Dieu, Vous avez fait du bienheureux Norbert, Votre Confesseur et Pontife, un excellent prédicateur de Votre parole et Vous avez donné par lui à Votre Église une nouvelle famille : faites, nous Vous en supplions, qu’aidés de ses mérites, nous puissions, grâce à Votre secours, mettre en pratique ce qu’il a enseigné par ses paroles et par ses œuvres. \perdominumfr.}
\feast{0609}{Ss. Primi et Feliciani Martyrum}{9 juin}
\saintindex{}
\saintindex{}
\oratio{Fac nos, quǽsumus, Dómine, sanctórum Mártyrum tuórum Primi et Feliciáni semper festa sectári: quorum suffrágiis protectiónis tuæ dona sentiámus. \perdominumla}{Nous Vous en supplions, Seigneur, faites que nous célébrions toujours fidèlement la fête de Vos saints Martyrs Prime et Félicien, afin que, grâce à leur intercession, nous éprouvions les bienfaits de Votre protection. \perdominumfr}
\feast{0610}{S. Margaritæ Reginæ Viduæ}{10 juin}
\saintindex{}
\saintindex{}
\oratio{Deus, qui beátam Margarítam regínam exímia in páuperes caritáte mirábilem effecísti: da; ut, ejus intercessióne et exémplo, tua in córdibus nostris cáritas júgiter augeátur. \perdominumla}{Dieu, Vous avez rendu admirable la bienheureuse reine Marguerite en lui inspirant une extrême charité pour les pauvres : faites que, par son intercession et à son exemple, Votre charité croisse continuellement dans nos cœurs. \perdominumfr}
\feast{0611}{S. Barnabæ Apostoli}{11 juin}
\saintindex{}
\saintindex{}
\oratio{Deus, qui nos beáti Bárnabæ Apóstoli tui méritis et intercessióne lætíficas: concéde propítius; ut, qui tua per eum benefícia póscimus, dono tuæ grátiæ consequámur. \perdominumla}{Dieu, Vous nous donnez un motif de joie dans les mérites et l’intercession du bienheureux Barnabé, Votre Apôtre : accordez-nous, avec bonté, qu’en recourant à cette intercession pour solliciter Vos bienfaits, nous les obtenions au moyen de Votre grâce. \perdominumfr}
\feast{0612}{S. Joannis a S. Facundo Confessoris}{12 juin}
\saintindex{}
\saintindex{}
\oratio{Deus, auctor pacis et amátor caritátis, qui beátum Joánnem Confessórem tuum mirífica dissidéntes componéndi grátia decorásti: ejus méritis et intercessióne concéde; ut, in tua caritáte firmáti, nullis a te tentatiónibus separémur. \perdominumla}{Dieu, auteur de la paix et amant de la charité, Vous avez orné le bienheureux Jean, Votre Confesseur, d’un merveilleux don du ciel pour apaiser les différends : accordez-nous, par ses mérites et son intercession, d’être tellement affermis dans Votre amour, que nous ne soyons plus séparés de Vous par aucune tentation. \perdominumfr}
\feast{0612}{Ss. Basilidis, Cyrini, Naboris and Nazarii Martyrum}{12 juin}
\saintindex{}
\saintindex{}
\oratio{Sanctórum Mártyrum tuórum Basílidis, Cyríni, Náboris atque Nazárii, quǽsumus, Dómine, natalítia nobis votíva respléndeant: et, quod illis cóntulit excelléntia sempitérna, frúctibus nostræ devotiónis accréscat. \perdominumla}{Faites, s’il Vous plaît, Seigneur, que l’anniversaire de la naissance au ciel de Vos saints Martyrs Basilide, Cyrin, Nabor et Nazaire répande la lumière dans nos âmes et que les biens spirituels dont le bonheur éternel leur a procuré la plénitude, s’accroissent en nous grâce aux fruits de notre dévotion. \perdominumfr}
\feast{0613}{S. Antonii de Padua Confessoris}{13 juin}
\saintindex{Antóni}
\saintindex{Antoine}
\oratio{@:OratioText:s/ atque Doctóris//
(sed rubrica 195 aut rubrica 196 aut rubrica altovadensis)
@:OratioText}{@:OratioText:s/ atque Doctóris//
(sed rubrica 195 aut rubrica 196)
@:OratioText}
\oratio{Ecclésiam tuam, Deus, beáti Antónii Confessóris tui atque Doctóris solémnitas votíva lætíficet: ut spirituálibus semper muniátur auxíliis et gáudiis pérfrui mereátur ætérnis. \perdominumla}{Que la solennité annuelle de Votre Confesseur et Docteur, le bienheureux Antoine, réjouisse Votre Église, ô Dieu, afin qu’elle soit toujours munie des secours spirituels et qu’elle mérite de goûter les joies éternelles. \perdominumfr}
\feast{0614}{S. Basilii Magni, Episcopis Confessoris et Ecclesiæ Doctoris}{14 juin}
\saintindex{Basilíi}
\saintindex{Basile}
\oratio{@Commune/C4-1}{Nous Vous supplions, Seigneur, d’exaucer les prières que nous Vous adressons en la solennité du bienheureux Basile, Votre Confesseur et Pontife, et de nous accorder, grâce aux mérites et à l’intercession de celui qui Vous a si dignement servi, le pardon de tous nos péchés. \perdominumfr}
\feast{0615}{Ss. Viti, Modesti atque Crescentiæ Martyrum}{15 juin}
\saintindex{}
\saintindex{}
\oratio{Da Ecclésiæ tuæ, quǽsumus, Dómine, sanctis Martýribus tuis Vito, Modésto atque Crescéntia intercedéntibus, supérbe non sápere, sed tibi plácita humilitáte profícere: ut, prava despíciens, quæcúmque recta sunt, líbera exérceat caritáte. \perdominumla}{Nous Vous en prions, Seigneur, faites que, par l’intercession de Vos saints Martyrs Vite, Modeste et Crescence, Votre Église, éloignée de tout sentiment d’orgueil, professe l’humilité qui a le don de Vous plaire, afin que, méprisant ce qui est mal, elle pratique avec amour et liberté tout ce qui est bien. \perdominumfr}
\feast{0617}{S. Gregorii Barbadici Episcopi et Confessoris}{17 juin}
\saintindex{Gregórium}
\saintindex{Grégoire}
\oratio{Deus, qui beátum Gregórium Confessórem tuum atque Pontíficem pastoráli sollicitúdine, et páuperum miseratióne claréscere voluísti: concéde propítius; ut, cujus mérita celebrámus, caritátis imitémur exémpla. \perdominumla}{Seigneur, Vous avez voulu illuminer le bienheureux Grégoire, Votre Confesseur et Pontife, par le zèle pastoral et la tendresse pour les pauvres : accordez-nous favorablement qu’en célébrant ses mérites, nous prenions exemple sur sa charité. \perdominumfr}
\feast{0618}{S. Ephræm Syri Confessoris et Ecclesiæ Doctoris}{18 juin}
\saintindex{Ephræm}
\saintindex{Éphrem}
\oratio{Deus, qui Ecclésiam tuam beáti Ephræm Confessóris tui et Doctóris mira eruditióne et præcláris vitæ méritis illustráre voluísti: te súpplices exorámus; ut, ipso intercedénte, eam advérsus erróris et pravitátis insídias perénni tua virtúte deféndas. \perdominumla}{Ô Dieu, vous avez voulu illustrer votre Église par l’admirable érudition et les mérites éclatants du bienheureux Éphrem, votre Confesseur et Docteur : nous vous en supplions, daignez à son intercession la défendre par votre constant secours contre les embûches de l’erreur et de la dépravation. \perdominumfr}
\feast{0618}{SS. Marci et Marcelliani Martyrum}{18 juin}
\saintindex{}
\saintindex{}
\oratio{Præsta, quǽsumus, omnípotens Deus: ut, qui sanctórum Mártyrum tuórum Marci et Marcelliáni natalítia cólimus; a cunctis malis imminéntibus eórum intercessiónibus liberémur. \perdominumla}{Accordez-nous, s’il vous plaît, ô Dieu tout-puissant, que, célébrant la naissance au ciel de vos bienheureux Martyrs Marc et Marcellien, nous soyons délivrés, grâce à leur intercession, de tous les maux qui nous menacent. \perdominumfr}
\feast{0619}{S. Julianæ de Falconeriis Virginis}{19 juin}
\saintindex{}
\saintindex{}
\oratio{Deus, qui beátam Juliánam Vírginem tuam extrémo morbo laborántem, pretióso Fílii tui Córpore mirabíliter recreáre dignátus es: concéde, quǽsumus; ut, ejus intercedéntibus méritis, nos quoque eódem in mortis agóne refécti ac roboráti, ad cæléstem pátriam perducámur. \pereumdemla}{Dieu, vous avez daigné réconforter admirablement par le Corps précieux de votre Fils la bienheureuse Julienne, votre Vierge, peinant sous le poids de sa dernière maladie : accordez-nous, nous vous en prions, par l’intercession de ses mérites, d’être nous aussi nourris et fortifiés dans l’agonie de la mort et ainsi de parvenir à la patrie céleste. \pereumdemfr}
\feast{0619}{Ss. Gervasii et Protasii Martyrum}{19 juin}
\saintindex{Gervásii et Protásii}
\saintindex{Gervais et Protais}
\feast{0620}{S. Silverii Papæ et Martyris}{20 juin}
\saintindex{Silvérii}
\saintindex{Silvère}
\feast{0620}{}{20 juin}
\saintindex{}
\saintindex{}
\feast{0621}{S. Aloisii Gonzagæ Confessoris}{21 juin}
\saintindex{}
\saintindex{}
\oratio{Cæléstium donórum distribútor, Deus, qui in angélico júvene Aloísio miram vitæ innocéntiam pari cum pœniténtia sociásti: ejus méritis et précibus concéde; ut, innocéntem non secúti, pœniténtem imitémur. \perdominumla}{Ô Dieu, vous distribuez les biens célestes, et vous avez réuni dans le jeune et angélique Louis, une merveilleuse innocence à la pratique de la mortification : faites qu’en nous appuyant sur ses mérites et son intercession, si nous n’avons pas sa pureté, nous imitions au moins sa pénitence. \perdominumfr}
\feast{0622}{S. Paulini Episcopi et Confessoris}{22 juin}
\saintindex{}
\saintindex{}
\oratio{Deus, qui ómnia pro te in hoc sǽculo relinquéntibus, céntuplum in futúro et vitam ætérnam promisísti: concéde propítius; ut, sancti Pontíficis Paulíni vestígiis inhæréntes, valeámus terréna despícere et sola cæléstia desideráre:
$Qui vivis}{Dieu, Vous avez promis à ceux qui abandonnent tout en ce siècle pour Vous, le centuple dans le siècle à venir et la vie éternelle : accordez-nous, dans Votre bonté que, suivant fidèlement les traces du saint Pontife Paulin, nous ayons la force de mépriser les biens de la terre et de désirer les seuls biens du ciel.
$Qui vivis}
\feast{0622}{Ss. John Fisher and Thomas More, martyrs}{22 juin}
\saintindex{}
\saintindex{}
\feast{0622}{S. Paulini Episcopi et Confessoris}{22 juin}
\saintindex{Paulíni}
\saintindex{Paulíni}
\feast{0623}{In Vigilia S. Joannis Baptistæ}{23 juin}
\saintindex{}
\saintindex{}
\oratio{Præsta, quǽsumus, omnípotens Deus: ut família tua per viam salútis incédat; et beáti Joánnis Præcursóris hortaménta sectándo, ad eum quem prædíxit, secúra pervéniat, Dóminum nostrum Jesum Christum, Fílium tuum:
$Qui tecum}{Qu’il Vous plaise, ô Dieu tout-puissant, d’accorder à Votre famille de marcher dans la voie du salut ; afin que, fidèle aux enseignements du bienheureux Jean, le Précurseur, elle parvienne sûrement jusqu’à Celui qu’il eut mission d’annoncer, notre Seigneur Jésus-Christ, Votre Fils.
$Qui tecum}
\feast{0624}{In Nativitate S. Joannis Baptistæ}{24 juin}
\saintindex{}
\saintindex{}
\oratio{Deus, qui præséntem diem honorábilem nobis in beáti Joánnis nativitáte fecísti: da pópulis tuis spirituálium grátiam gaudiórum; et ómnium fidélium mentes dírige in viam salútis ætérnæ. \perdominumla}{Dieu, vous nous avez rendu ce jour vénérable par la nativité du bienheureux Jean : accordez à votre peuple la grâce des joies spirituelles ; et dirigez les âmes de tous les fidèles dans la voie du salut éternel. \perdominumfr}
\feast{0625}{S. Gulielmi Abbatis}{25 juin}
\saintindex{}
\saintindex{}
\oratio{Deus, qui infirmitáti nostræ ad teréndam salútis viam in Sanctis tuis exémplum et præsídium collocásti: da nobis, ita beáti Guliélmi abbátis mérita venerári; ut ejúsdem excipiámus suffrágia, et vestígia prosequámur. \perdominumla}{Ô Dieu, pour aplanir à notre faiblesse la voie du salut, Vous nous avez donné l’exemple et la protection de Vos Saints : faites que nous honorions les mérites du bienheureux Guillaume, Abbé, de manière à mériter le secours de ses prières et à marcher sur ses traces. \perdominumfr}
\feast{0625}{Die secunda infra Octavam Nativitatis S. Joannis Baptistæ}{25 juin}
\saintindex{}
\saintindex{}
\feast{0626}{Ss. Joannis et Pauli Martyrum}{26 juin}
\saintindex{}
\saintindex{}
\oratio{Quǽsumus, omnípotens Deus: ut nos gemináta lætítia hodiérnæ festivitátis excípiat, quæ de beatórum Joánnis et Pauli glorificatióne procédit; quos éadem fides et pássio vere fecit esse germános. \perdominumla}{Nous vous prions, Dieu tout-puissant, faites-nous entrer dans la joie de cette double fête, joie qui provient de la glorification des bienheureux Jean et Paul qu’une même foi et un même martyre ont rendus vraiment frères. \perdominumfr}
\feast{0626}{Die tertia infra Octavam Nativitatis S. Joannis Baptistæ}{26 juin}
\saintindex{}
\saintindex{}
\feast{0627}{Quarta die infra Octavam Nativitatis S. Joannis Baptistæ}{27 juin}
\saintindex{}
\saintindex{}
\feast{0628}{S. Irenæi Episcopi et Martyris}{28 juin}
\saintindex{}
\saintindex{}
\oratio{Deus, qui beáto Irenǽo Mártyri tuo atque Pontífici tribuísti, ut et veritáte doctrínæ expugnáret hǽreses, et pacem Ecclésiæ felíciter confirmáret: da, quǽsumus, plebi tuæ in sancta religióne constántiam; et pacem tuam nostris concéde tempóribus. \perdominumla}{Dieu, vous avez accordé au bienheureux Irénée, votre Martyr et Pontife, de réprimer les hérésies par la vérité de sa doctrine et d’affermir la paix de l’Église : nous vous en supplions, donnez à votre peuple la constance dans la sainte religion, et à nos temps la paix. \perdominumfr}
\feast{0628}{Die quinta infra Octavam Nativitatis S. Joannis Baptistæ}{28 juin}
\saintindex{}
\saintindex{}
\feast{0628}{In Vigilia Ss. Petri et Pauli Apostolorum}{28 juin}
\saintindex{}
\saintindex{}
\oratio{Præsta, quǽsumus, omnípotens Deus: ut nullis nos permíttas perturbatiónibus cóncuti; quos in apostólicæ confessiónis petra solidásti. \perdominumla}{Daignez, nous vous en supplions, ô Dieu tout-puissant, ne point permettre qu’aucun trouble nous ébranle, après que vous nous avez établis sur la pierre solide de la foi des Apôtres. \perdominumfr}
\feast{0628}{}{28 juin}
\saintindex{}
\saintindex{}
\feast{0629}{SS. Apostolorum Petri et Pauli}{29 juin}
\saintindex{}
\saintindex{}
\oratio{Deus, qui hodiérnam diem Apostolórum tuórum Petri et Pauli martýrio consecrásti: da Ecclésiæ tuæ, eórum in ómnibus sequi præcéptum; per quos religiónis sumpsit exórdium. \perdominumla}{Dieu, vous avez consacré ce jour par le martyre de vos Apôtres saint Pierre et saint Paul : faites la grâce à votre Église de suivre en tout le précepte de ceux par qui la religion a commencé. \perdominumfr}
\feast{0629}{Die sexta infra Octavam Nativitatis S. Joannis Baptistæ}{29 juin}
\saintindex{}
\saintindex{}
\feast{0630}{In Commemoratione S. Pauli Apostoli}{30 juin}
\saintindex{}
\saintindex{}
\oratio{Deus, qui multitúdinem géntium beáti Pauli Apóstoli prædicatióne docuísti: da nobis, quǽsumus; ut, cujus natalícia cólimus, ejus apud te patrocínia sentiámus. \perdominumla
_
@Sancti/01-25:Commemoratio4}{Ô Dieu, vous avez instruit la multitude des nations par la parole de l’Apôtre saint Paul, accordez à notre demande, qu’en fêtant sa naissance céleste, nous sentions l’effet de son patronage auprès de vous. \perdominumfr
_
@Sancti/01-25:Commemoratio4}
\feast{0630}{First Martyrs of the Church of Rome}{30 juin}
\saintindex{}
\saintindex{}
\feast{0630}{Die septima infra Octavam Nativitatis S. Joannis Baptistæ}{30 juin}
\saintindex{}
\saintindex{}
\feast{0701}{Pretiosissimi Sanguinis Domini Nostri Jesu Christi}{1 juillet}
\saintindex{}
\saintindex{}
\oratio{Omnípotens sempitérne Deus, qui unigénitum Fílium tuum mundi Redemptórem constituísti, ac ejus Sánguine placári voluísti: concéde, quǽsumus, salútis nostræ prétium solémni cultu ita venerári, atque a præséntis vitæ malis ejus virtúte deféndi in terris; ut fructu perpétuo lætémur in cælis. \pereumdemla}{Dieu tout-puissant et éternel, Vous avez établi Votre Fils unique rédempteur du monde et Vous avez voulu que Votre justice soit apaisée par Son sang : faites-nous la grâce, nous Vous en prions, de vénérer d’un culte solennel ce prix de notre salut et d’être ici-bas préservés par Sa vertu des maux de la vie présente de manière à jouir éternellement de Ses fruits dans les cieux. \pereumdemfr}
\feast{0701}{Tertia die infra Octavam Ss. Apostolorum Petri et Pauli}{1 juillet}
\saintindex{}
\saintindex{}
\oratio{@Sancti/06-29}{}
\feast{0701}{In Octava S. Joannis Baptistæ}{1 juillet}
\saintindex{}
\saintindex{}
\feast{0702}{}{2 juillet}
\saintindex{}
\saintindex{}
\oratio{Fámulis tuis, quǽsumus, Dómine, cæléstis grátiæ munus impertíre: ut, quibus beátæ Vírginis partus éxstitit salútis exórdium; Visitatiónis ejus votíva solémnitas, pacis tríbuat increméntum. \perdominumla}{Seigneur, nous Vous prions d’accorder à Vos serviteurs le don de la grâce céleste : et, comme l’enfantement de la bienheureuse Vierge a été le principe de leur salut, qu’ainsi la pieuse solennité de sa Visitation leur procure un accroissement de paix. \perdominumfr}
\oratio{Omnípotens sempitérne Deus, qui ex abundántia charitátis, beátam Maríam tuo Fílio fœcundátam ad salutatiónem Elísabeth inspirásti: præsta quæsumus; * ut per ejus Visitatiónem donis cœléstibus repleámur, et ab ómnibus adversitátibus eruámur. \pereumdemla}{}
\feast{0702}{Ss. Processi et Martiniani Martyrum}{2 juillet}
\saintindex{}
\saintindex{}
\oratio{v. Deus, qui nos sanctórum Mártyrum tuórum Procéssi et Martiniáni gloriósis confessiónibus circúmdas et prótegis: da nobis et eórum imitatióne profícere, et intercessióne gaudére. \perdominumla}{v. Dieu, Vous nous donnez dans la glorieuse profession de foi de Vos saints Martyrs Processus et Martinien un gage de Votre secours et de Votre protection : accordez-nous de profiter de leur exemple et de nous réjouir de leur intercession. \perdominumfr}
\feast{0702}{}{2 juillet}
\saintindex{}
\saintindex{}
\feast{0702}{Quarta die infra Octavam Ss. Apostolorum Petri et Pauli}{2 juillet}
\saintindex{}
\saintindex{}
\oratio{@Sancti/07-03oct}{}
\feast{0703}{S. Leonis Papæ et Confessoris}{3 juillet}
\saintindex{Leónem}
\saintindex{Léon}
\oratio{Deus, qui beátum Leónem Pontíficem sanctórum tuórum méritis coæquásti: concéde propítius; ut, qui commemoratiónis ejus festa percólimus, vitæ quoque imitémur exémpla. \perdominumla}{Dieu, Vous avez égalé le bienheureux Pontife Léon en mérites à Vos Saints, accordez-nous, dans Votre bonté que, comme nous célébrons la mémoire de sa fête, nous imitions aussi les exemples de sa vie. \perdominumfr}
\feast{0703}{Quinta die infra Octavam Ss. Apostolorum Petri et Pauli}{3 juillet}
\saintindex{}
\saintindex{}
\oratio{@Sancti/06-29}{}
\feast{0703}{}{3 juillet}
\saintindex{}
\saintindex{}
\feast{0704}{Sexta die infra Octavam Ss. Apostolorum Petri et Pauli}{4 juillet}
\saintindex{}
\saintindex{}
\oratio{@Sancti/06-29}{}
\feast{0705}{S. Antonii Mariæ Zaccaria Confessoris}{5 juillet}
\saintindex{}
\saintindex{}
\oratio{Fac nos, Dómine Deus, supereminéntem Jesu Christi sciéntiam, spíritu Pauli Apóstoli, edíscere: qua beátus Antónius María mirabíliter erudítus, novas in Ecclésia tua clericórum et vírginum famílias congregávit. \pereumdemla}{Faites-nous, Seigneur, la grâce d’apprendre selon l’esprit de l’Apôtre Paul, la science suréminente de Jésus-Christ dont le bienheureux Antoine-Marie fut merveilleusement instruit, lui qui rassembla dans Votre Église de nouvelles familles de clercs et de religieuses. \pereumdemfr}
\feast{0705}{Septima die infra Octavam Ss. Apostolorum Petri et Pauli}{5 juillet}
\saintindex{}
\saintindex{}
\oratio{@Sancti/06-29}{}
\feast{0706}{In Octava Ss. Apostolorum Petri et Pauli}{6 juillet}
\saintindex{}
\saintindex{}
\oratio{Deus, cujus déxtera beátum Petrum ambulántem in fluctibus, ne mergerétur, eréxit, et coapóstolum ejus Paulum, tértio naufragántem, de profúndo pélagi liberávit: exáudi nos propítius, et concéde; ut, ambórum meritis, æternitátis glóriam consequámur:
$Qui vivis}{Dieu, Dont la droite a relevé le bienheureux Pierre marchant sur les flots pour l'empêcher de couler et a retiré de la profondeur de la mer Paul, son collègue d'apostolat, à son troisième naufrage, exaucez-nous dans Votre bonté et accordez-nous d'arriver, par leurs mérites à tous les deux, à la gloire de l'éternité.
$Qui vivis}
\feast{0706}{S. Maríæ Goretti Virginis et Martyris}{6 juillet}
\saintindex{Maríæ}
\saintindex{}
\feast{0706}{}{6 juillet}
\saintindex{}
\saintindex{}
\feast{0707}{Ss. Cyrilli et Methodii Pont. et Conf.}{7 juillet}
\saintindex{}
\saintindex{}
\oratio{Omnípotens sempitérne Deus, qui Slavóniæ gentes per beátos Confessóres tuos atque Pontífices Cyríllum et Methódium ad agnitiónem tui nóminis veníre tribuísti: præsta; ut, quorum festivitáte gloriámur, eórum consórtio copulémur. \perdominumla}{Dieu tout puissant et éternel, Vous avez accordé aux peuples slaves d’arriver à la connaissance de Votre nom, par le ministère de Vos bienheureux Pontifes et Confesseurs Cyrille et Méthode : faites que, mettant notre gloire à célébrer leur fête, nous soyons associés à leur sort commun. \perdominumfr}
\feast{0708}{S. Elisabeth Reg. Portugaliæ Viduæ}{8 juillet}
\saintindex{}
\saintindex{}
\oratio{Clementíssime Deus, qui beátam Elísabeth regínam, inter céteras egrégias dotes, béllici furóris sedándi prærogatíva decorásti: da nobis, ejus intercessióne; post mortális vitæ, quam supplíciter pétimus, pacem, ad ætérna gáudia perveníre. \perdominumla}{Dieu très clément, parmi tant d’autres qualités éminentes, Vous avez donné à la bienheureuse reine Élisabeth la vertu d’apaiser les fureurs de la guerre : accordez-nous, par son intercession, qu’après avoir, pendant cette vie mortelle, joui de la paix que nous Vous demandons humblement, nous parvenions aux joies éternelles. \perdominumfr}
\feast{0708}{}{8 juillet}
\saintindex{}
\saintindex{}
\feast{0709}{S. Augustine Zhao Rong and companions, martyrs}{9 juillet}
\saintindex{}
\saintindex{}
\feast{0710}{Ss. Septem Fratrum Martyrum, ac Rufinæ et Secundæ Virginum et Martyrum}{10 juillet}
\saintindex{}
\saintindex{}
\oratio{Præsta, quǽsumus, omnípotens Deus: ut, qui gloriósos Mártyres fortes in sua confessióne cognóvimus, pios apud te in nostra intercessióne sentiámus. \perdominumla}{Dieu tout puissant, faites, nous Vous en prions, qu’ayant connu combien Vos glorieux Martyrs ont été fermes dans la confession de leur foi, nous ressentions les effets de leur pieuse intercession auprès de Vous. \perdominumfr}
\feast{0711}{S. Pii I Papæ et Martyris}{11 juillet}
\saintindex{Pii}
\saintindex{Pie}
\feast{0711}{}{11 juillet}
\saintindex{}
\saintindex{}
\feast{0712}{S. Joannis Gualberti Abbatis}{12 juillet}
\saintindex{Joánnis}
\saintindex{Jean}
\feast{0712}{Ss. Naboris et Felicis Martyrum}{12 juillet}
\saintindex{}
\saintindex{}
\oratio{Præsta, quæsumus, Domine: ut, sicut nos sanctorum Martyrum tuorum Naboris et Felicis natalitia celebranda non deserunt; ita jugiter suffragiis comitentur. \perdominumla}{Faites, Seigneur, s’il Vous plaît, que, comme nous célébrons fidèlement la naissance au ciel de Vos saints Martyrs Nabor et Félix, ainsi nous soyons constamment aidés de leur prières. \perdominumfr}
\feast{0713}{S. Anacleti Papæ et Martyris}{13 juillet}
\saintindex{Anacléti}
\saintindex{Anaclet}
\feast{0714}{S. Bonaventuræ Episcopi Confessoris et Ecclesiæ Doctoris}{14 juillet}
\saintindex{Bonaventúram}
\saintindex{Bonaventure}
\feast{0715}{S. Henrici Imperatoris Confessoris}{15 juillet}
\saintindex{}
\saintindex{}
\oratio{Deus, qui hodiérna die beátum Henrícum Confessórem tuum e terréni cúlmine impérii ad regnum ætérnum transtulísti: te súpplices exorámus; ut, sicut illum, grátiæ tuæ ubertáte prævéntum, illécebras sǽculi superáre fecísti, ita nos fácias, ejus imitatióne, mundi hujus blandiménta vitáre, et ad te puris méntibus perveníre. \perdominumla}{Ô Dieu, en ce jour, vous avez fait passer le bienheureux Henri, Votre Confesseur, du sommet de l’empire de la terre au royaume du ciel : nous Vous demandons en suppliant que, comme en le prévenant par l’abondance de Votre grâce, Vous l’avez fait triompher des attraits du siècle, Vous nous fassiez aussi, à son imitation, éviter les séductions du monde et parvenir jusqu’à Vous avec des cœurs purs. \perdominumfr}
\feast{0716}{}{16 juillet}
\saintindex{}
\saintindex{}
\oratio{Deus, qui beatíssimæ semper Vírginis et Genetrícis tuæ Maríæ singulári título Carméli órdinem decorásti: concéde propítius; ut, cujus hódie Commemoratiónem solémni celebrámus offício, ejus muníti præsídiis, ad gáudia sempitérna perveníre mereámur:
$Qui vivis}{Dieu, Vous avez orné l’Ordre du Carmel du titre particulier de la bienheureuse Marie toujours Vierge et Votre Mère : accordez-nous, dans Votre bonté, que, soutenus de la protection de celle dont nous honorons aujourd’hui solennellement la mémoire, nous méritions de parvenir aux joies éternelles.
$Qui vivis}
\feast{0716}{}{16 juillet}
\saintindex{}
\saintindex{}
\oratio{@Sancti/07-16}{}
\feast{0717}{S. Alexii Confessoris}{17 juillet}
\saintindex{Aléxii}
\saintindex{Alexis}
\feast{0718}{S. Camilli de Lellis Confessoris}{18 juillet}
\saintindex{}
\saintindex{}
\oratio{Deus, qui sanctum Camíllum, ad animárum in extrémo agóne luctántium subsídium, singulári caritátis prærogatíva decorásti: ejus, quǽsumus, méritis, spíritum nobis tuæ dilectiónis infúnde; ut in hora éxitus nostri hostem víncere, et ad cæléstem mereámur corónam perveníre. \perdominumla}{Dieu, Vous avez fait don à saint Camille d’une charité extraordinaire pour aider les âmes dans la lutte suprême de l’agonie : nous Vous supplions, par ses mérites, de répandre en nous l’esprit de Votre charité afin que nous puissions, à l’heure du trépas, vaincre l’ennemi et parvenir à la céleste couronne. \perdominumfr}
\feast{0718}{Ss. Symphorosæ et Septem Filiorum Martyrum}{18 juillet}
\saintindex{Symphorósæ et filiórum ejus}
\saintindex{Symphorose et ses fils}
\feast{0719}{S. Vincentii a Paulo Confessoris}{19 juillet}
\saintindex{}
\saintindex{}
\oratio{Deus, qui, ad evangelizándum paupéribus et ecclesiástici órdinis decórem promovéndum, beátum Vincéntium apostólica virtúte roborásti: præsta, quǽsumus; ut, cujus pia mérita venerámur, virtútum quoque instruámur exémplis. \perdominumla}{Dieu, Vous avez fortifié le bienheureux Vincent par un courage apostolique pour évangéliser les pauvres et augmenter la gloire de l’ordre ecclésiastique ; faites, s’il Vous plaît, qu’en honorant sa piété et ses mérites, l’exemple de ses vertus nous instruise. \perdominumfr}
\feast{0720}{S. Hieronymi Æmiliani Confessoris}{20 juillet}
\saintindex{}
\saintindex{}
\oratio{Deus, misericordiárum pater, per mérita et intercessiónem beáti Hierónymi, quem órphanis adjutórem et patrem esse voluísti: concéde; ut spíritum adoptiónis, quo fílii tui nominámur et sumus, fidéliter custodiámus. \perdominumla}{Dieu, Père des miséricordes, par les mérites et l’intercession du bienheureux Jérôme, que Vous avez donné pour soutien et pour père aux orphelins, faites-nous la grâce de conserver fidèlement cet esprit d’adoption en vertu duquel nous sommes appelés Vos fils et le devenons réellement. \perdominumfr}
\feast{0720}{S. Margaritæ, Virginis et Martyris}{20 juillet}
\saintindex{Margaríta}
\saintindex{Marguerite}
\feast{0721}{S. Praxedis Virginis}{21 juillet}
\saintindex{Praxédis}
\saintindex{Praxède}
\feast{0721}{S. Laurentii a Brundusio Confessoris et Ecclesiæ Doctoris}{21 juillet}
\saintindex{Lauréntio}
\saintindex{Laurent}
\oratio{Deus, qui ad árdua quæque pro nóminis tui glória et animárum salúte beáto Lauréntio, Confessóri tuo atque Doctóri, spíritum sapiéntiæ et fortitúdinis contulísti: da nobis in eódem spíritu et agénda cognóscere; et cógnita, ejus intercessióne, perfícere. \perdominumla}{Dieu, pour travailler avec ardeur à la gloire de Votre nom et au salut des âmes, Vous avez accordé au bienheureux Laurent, Votre Confesseur et Docteur, l’esprit de sagesse et de force : donnez-nous, dans ce même esprit, de connaître ce que nous devons faire et, une fois que nous le connaîtrons, la force de l’accomplir, avec le secours de son intercession. \perdominumfr}
\feast{0722}{S. Mariæ Magdalenæ Pœnitentis}{22 juillet}
\saintindex{}
\saintindex{}
\oratio{Beátæ Maríæ Magdalénæ, quǽsumus, Dómine, suffrágiis adjuvémur: cujus précibus exorátus, quatriduánum fratrem Lázarum vivum ab ínferis resuscitásti:
$Qui vivis}{Nous Vous prions, Seigneur, par les suffrages de la bienheureuse Marie-Madeleine, de venir à notre aide : Vous qui, fléchi par ses prières, avez ressuscité vivant des enfers son frère Lazare, mort depuis quatre jours.
$Qui vivis}
\feast{0723}{S. Apollinaris Episcopi et Martyris}{23 juillet}
\saintindex{}
\saintindex{}
\oratio{Deus, fidélium remunerátor animárum, qui hunc diem beáti Apollináris Sacerdótis tui martýrio consecrásti: tríbue nobis, quǽsumus, fámulis tuis; ut, cujus venerándam celebrámus festivitátem, précibus ejus indulgéntiam consequámur. \perdominumla}{Dieu, qui récompensez les âmes fidèles, Vous avez consacré ce jour par le martyre du bienheureux Apollinaire, Votre Prêtre : faites, s’il Vous plaît, que nous, qui sommes Vos serviteurs, nous obtenions Votre pardon au moyen des prières de celui dont nous célébrons la fête vénérable. \perdominumfr}
\feast{0723}{S. Liborii Episcopi et Confessoris}{23 juillet}
\saintindex{Libórii}
\saintindex{Liboire}
\oratio{@Commune/C4
(sed feria 7 nisi rubrica 1955 nisi rubrica 196)
@Commune/C4-1
}{}
\feast{0724}{In Vigilia S. Jacobi Ap.}{24 juillet}
\saintindex{Jacóbi}
\saintindex{Jacques}
\feast{0724}{S. Sharbel Makhluf, hermit}{24 juillet}
\saintindex{}
\saintindex{}
\feast{0724}{S. Christinæ Virginis et Martyris}{24 juillet}
\saintindex{Christína}
\saintindex{Christine}
\feast{0725}{S. Jacobi Apostoli}{25 juillet}
\saintindex{}
\saintindex{}
\oratio{Esto, Dómine, plebi tuæ sanctificátor et custos: ut, Apóstoli tui Jacóbi muníta præsídiis, et conversatióne tibi pláceat, et secúra mente desérviat. \perdominumla}{Seigneur, soyez le sanctificateur de Votre peuple et son gardien afin qu’aidé par l’assistance de votre Apôtre Jacques, il mène une vie qui Vous soit agréable et Vous serve avec tranquillité et confiance. \perdominumfr}
\feast{0725}{S. Christophori Martyris}{25 juillet}
\saintindex{Christóphori}
\saintindex{Christophe}
\feast{0726}{S. Annæ Matris B.M.V.}{26 juillet}
\saintindex{}
\saintindex{}
\oratio{Deus, qui beátæ Annæ grátiam conférre dignátus es, ut Genetrícis unigéniti Fílii tui mater éffici mererétur: concéde propítius; ut, cujus solémnia celebrámus, ejus apud te patrocíniis adjuvémur. \pereumdemla}{Dieu, vous avez daigné donner à sainte Anne la grâce de mériter d’être la mère de celle par laquelle votre Fils unique est né, faites que nous soyons aidés de son patronage en ce jour où nous célébrons sa solennité. \pereumdemfr}
\feast{0726}{S. Annae et Joachim Parentes B.M.V.}{26 juillet}
\saintindex{}
\saintindex{}
\feast{0727}{S. Pantaleonis Martyris}{27 juillet}
\saintindex{Pantaleóne}
\saintindex{Pantaléon}
\feast{0728}{Ss. Nazarii et Celsi Martyrum, Victoris I Papæ et Martyris ac Innocentii I Papæ et Confessoris}{28 juillet}
\saintindex{}
\saintindex{}
\oratio{Sanctórum tuórum nos, Dómine, Nazárii, Celsi, Victóris et Innocéntii conféssio beáta commúniat: et fragilitáti nostræ subsídium dignánter exóret. \perdominumla}{Que la bienheureuse profession de foi de Vos saints Nazaire, Celse, Victor et Innocent nous fortifie, Seigneur et qu’elle obtienne de Votre bonté des secours pour notre faiblesse. \perdominumfr}
\feast{0729}{S. Marthæ Virginis}{29 juillet}
\saintindex{Marthæ}
\saintindex{Marthe}
\feast{0729}{Ss. Felicis, Simplicii, Faustini et Beatricis Martyrum}{29 juillet}
\saintindex{}
\saintindex{}
\oratio{Præsta, quǽsumus, Dómine: ut, sicut pópulus christiánus Mártyrum tuórum Felícis, Simplícii, Faustíni et Beatrícis temporáli solemnitáte congáudet, ita perfruátur ætérna; et, quod votis celébrat, comprehéndat efféctu. \perdominumla}{Accordez-nous, s'il Vous plaît, Seigneur, que le peuple chrétien qui se réjouit ici-bas de la solennité temporelle de Vos Martyrs Félix, Simplice, Faustin et Béatrice, jouisse de même un jour de leur fête éternelle et qu'il possède éternellement ce qu'il célèbre par ses vœux. \perdominumfr}
\feast{0730}{S. Abdon et Sennen Martyrum}{30 juillet}
\saintindex{}
\saintindex{}
\oratio{Deus, qui sanctis tuis Abdon et Sennen ad hanc glóriam veniéndi copiósum munus grátiæ contulísti: da fámulis tuis suórum véniam peccatórum; ut, Sanctórum tuórum intercedéntibus méritis, ab ómnibus mereántur adversitátibus liberári. \perdominumla}{Dieu, Vous avez fait à Vos saints Abdon et Sennen le don insigne de la grâce d’arriver à cette gloire. Accordez à Vos serviteurs le pardon de leurs péchés afin que, aidés des mérites de Vos saints, nous puissions être délivrés de toute adversité. \perdominumfr}
\feast{0731}{S. Ignatii Confessoris}{31 juillet}
\saintindex{}
\saintindex{}
\oratio{Deus, qui ad majórem tui nóminis glóriam propagándam, novo per beátum Ignátium subsídio militántem Ecclésiam roborásti: concéde; ut, ejus auxílio et imitatióne certántes in terris, coronári cum ipso mereámur in cælis. \perdominumla}{Dieu, pour propager la plus grande gloire de Votre nom, vous avez, par le bienheureux Ignace, procuré à Votre Église militante de nouveaux renforts : accordez-nous, avec son secours et combattant à son exemple sur la terre, de mériter d’être couronnés avec lui dans le ciel. \perdominumfr}
\oratio{@Commune/C5-1::s/N\./Ignátii/}{}
\feast{0801}{S. Petri ad Vincula}{1 août}
\saintindex{}
\saintindex{}
\oratio{Deus, qui beátum Petrum Apóstolum, a vínculis absolútum, illǽsum abíre fecísti: nostrórum, quæsumus, absólve víncula peccatórum; et ómnia mala a nobis propitiátus exclúde. \perdominumla
_
@Sancti/02-22:Commemoratio4:s/\$Per Dominum//}{Ô Dieu, qui après avoir fait tomber les chaînes du bienheureux Pierre, Apôtre, l’avez fait sortir de prison, sans qu’il eût reçu aucun mal, nous Vous en prions, brisez les liens de nos péchés et, dans Votre bonté, éloignez de nous tous les maux. \perdominumfr
_
@Sancti/02-22:Commemoratio4:s/\$Per Dominum//}
\feast{0801}{Ss. Martyrum Machabæorum}{1 août}
\saintindex{}
\saintindex{}
\oratio{Fratérna nos, Dómine, Mártyrum tuórum coróna lætíficet: quæ et fídei nostræ prǽbeat increménta virtútum; et multíplici nos suffrágio consolétur. \perdominumla}{Faites, Seigneur, que la couronne fraternelle de Vos Martyrs soit une source de joie : qu’elle procure à notre foi une augmentation de force et qu’elle nous console par ces intercessions multiples. \perdominumfr}
\feast{0802}{S. Alfonsi Mariæ de Ligorio Episc. Conf. et Eccles. Doct.}{2 août}
\saintindex{Alfónsum Maríam}
\saintindex{Alphonse-Marie}
\oratio{Deus, qui per beátum Alfónsum Maríam Confessórem tuum atque Pontíficem, animárum zelo succénsum, Ecclésiam tuam nova prole fecundásti: quǽsumus; ut, ejus salutáribus mónitis edócti et exémplis roboráti, ad te perveníre felíciter valeámus. \perdominumla}{Dieu, Vous avez fécondé Votre Église d’une nouvelle famille par le ministère du bienheureux Alphonse-Marie, Votre Confesseur et Pontife, qui brûlait de zèle pour le salut des âmes : faites, nous Vous en prions, qu’instruits par ses leçons salutaires et fortifiés par ses exemples, nous puissions parvenir heureusement jusqu’à Vous. \perdominumfr}
\feast{0802}{S. Eusebius of Vercelli, bishop}{2 août}
\saintindex{}
\saintindex{}
\feast{0802}{S. Stephani Papæ Martyris}{2 août}
\saintindex{Stéphani}
\saintindex{Etienne}
\feast{0803}{De Inventione S. Stephani Protomartyris}{3 août}
\saintindex{}
\saintindex{}
\oratio{@Sancti/12-26::s/natalícia/Inventiónem/}{@Sancti/12-26::s/natalícia/Inventiónem/}
\feast{0804}{S. Dominici Confessoris}{4 août}
\saintindex{}
\saintindex{}
\oratio{Deus, qui Ecclésiam tuam beáti Domínici Confessóris tui illumináre dignátus es méritis et doctrínis: concéde; ut ejus intercessióne temporálibus non destituátur auxíliis, et spirituálibus semper profíciat increméntis. \perdominumla}{Dieu, Vous avez daigné éclairer votre Église par les mérites et les leçons du bienheureux Dominique, Votre Confesseur, faites que, par son intercession, elle ne soit pas privée des secours temporels et qu’elle fasse toujours de nouveaux progrès dans les voies spirituelles. \perdominumfr}
\feast{0805}{}{5 août}
\saintindex{}
\saintindex{}
\oratio{}{Accordez, Seigneur Dieu, à nous Vos serviteurs, de nous réjouir d’une constante santé de l’âme et du corps ; et, par la glorieuse intercession de la Bienheureuse Marie toujours Vierge, délivrez-nous des tristesses présentes et assurez-nous de jouir de la joie éternelle.  \perdominumfr}
\feast{0806}{}{6 août}
\saintindex{}
\saintindex{}
\oratio{Deus, qui fídei sacraménta in Unigéniti tui gloriósa Transfiguratióne patrum testimónio roborásti, et adoptiónem filiórum perféctam, voce delápsa in nube lúcida, mirabíliter præsignásti: concéde propítius; ut ipsíus Regis glóriæ nos coherédes effícias, et ejúsdem glóriæ tríbuas esse consórtes. \pereumdemla}{Dieu, vous avez affermi les mystères de la foi dans la glorieuse Transfiguration de Votre Fils unique, par le témoignage des anciens pères et, par la voix que Vous avez fait entendre dans la nuée lumineuse, Vous nous avez marqués de la grâce de la parfaite adoption : faites-nous, au moyen de Votre miséricorde, cohéritiers de ce même Roi de gloire et participants de cette même gloire. \pereumdemfr}
\feast{0806}{Ss. Xysti II Papæ, Felicissimi et Agapiti Martyrum}{6 août}
\saintindex{Xysti, Felicíssimi et Agapíti }
\saintindex{Sixte, Félicissime et Agapit}
\feast{0807}{S. Cajetani Confessoris}{7 août}
\saintindex{}
\saintindex{}
\oratio{Deus, qui beáto Cajetáno Confessóri tuo apostólicam vivéndi formam imitári tribuísti: da nobis, ejus intercessióne et exémplo, in te semper confídere et sola cæléstia desideráre. \perdominumla}{Dieu, Vous avez donné au bienheureux Gaétan, Votre Confesseur, d’imiter le genre de vie des Apôtres : accordez-nous, par son intercession et à son exemple, de mettre toujours en Vous notre confiance et de ne désirer que les biens du ciel. \perdominumfr}
\feast{0807}{S. Sixti II, Papæ, et Sociorum, Martyrum}{7 août}
\saintindex{}
\saintindex{}
\feast{0807}{S. Donati Episcopi et Martyris}{7 août}
\saintindex{}
\saintindex{}
\oratio{Deus, tuórum glória sacerdótum: præsta, quǽsumus, ut sancti Martyris tui et Epíscopi Donáti, cujus festa gérimus, sentiámus auxílium. \perdominumla}{Dieu, vous êtes la gloire de vos prêtres : faites, nous vous en prions, que nous ressentions le secours de votre saint Martyr et Évêque Donat, dont nous célébrons la fête. \perdominumfr}
\feast{0808}{Ss. Cyriaci, Largi et Smaragdi Martyrum}{8 août}
\saintindex{Cyríaci, Largi et Smarágdi}
\saintindex{Cyriaque, Large et Smaragde}
\feast{0809}{S. Joannis Mariæ Vianney Confessoris}{9 août}
\saintindex{}
\saintindex{}
\oratio{Omnípotens et miséricors Deus, qui sanctum Joánnem Maríam pastoráli stúdio et jugi oratiónis ac pœniténtiæ ardóre mirábilem effecísti: da, quǽsumus; ut, ejus exémplo et intercessióne, ánimas fratrum lucrári Christo, et cum eis ætérnam glóriam cónsequi valeámus. \pereumdemla}{Dieu tout-puissant et miséricordieux, Vous avez rendu saint Jean-Marie admirable par son zèle pastoral et son ardeur soutenue pour la prière et la pénitence : faites, nous Vous en prions, qu’à son exemple et par son intercession, nous puissions gagner au Christ les âmes de nos frères et obtenir avec eux la gloire éternelle. \pereumdemfr}
\feast{0809}{S. Romani Martyris}{9 août}
\saintindex{Románo}
\saintindex{Romain}
\feast{0809}{S. Terésiæ Benedictæ a Cruce (Edith Stein) Virginis et Martyris}{9 août}
\saintindex{Terésiæ}
\saintindex{Thérèse}
\feast{0809}{In Vigilia S. Laurentii Mart.}{9 août}
\saintindex{}
\saintindex{}
\oratio{Adésto, Dómine, supplicatiónibus nostris: et intercessióne beáti Lauréntii Mártyris tui, cujus prævenímus festivitátem; perpétuam nobis misericórdiam benígnus impénde. \perdominumla}{Exaucez nos supplications, Seigneur et, par l’intercession du bienheureux Laurent, Votre Martyr, dont nous devançons la fête, faites-nous ressentir avec bienveillance les effets de Votre miséricorde. \perdominumfr}
\feast{0810}{S. Laurentii Martyris}{10 août}
\saintindex{}
\saintindex{}
\oratio{Da nobis, quǽsumus, omnípotens Deus: vitiórum nostrórum flammas exstínguere; qui beáto Lauréntio tribuísti tormentórum suórum incéndia superáre. \perdominumla}{Nous Vous prions, Seigneur, d’éteindre en nous l’ardeur de nos vices, Vous qui avez donné au bienheureux Laurent la force de surmonter les flammes de ses tourments. \perdominumfr}
\feast{0811}{Ss. Tiburtii et Susannæ Virginis, Martyrum}{11 août}
\saintindex{}
\saintindex{}
\oratio{Sanctórum Mártyrum tuórum Tibúrtii et Susánnæ nos, Dómine, fóveant continuáta præsídia: quia non désinis propítius intuéri; quos tálibus auxíliis concésseris adjuvári. \perdominumla}{Faites, Seigneur, que nous soyons toujours soutenus par la protection de Vos saints Martyrs Tiburce et Suzanne puisque vous ne pouvez manquer d’accueillir favorablement ceux qui par Votre grâce, jouissent d’un tel appui. \perdominumfr}
\feast{0811}{Secunda die infra Octavam S. Laurentii Martyris}{11 août}
\saintindex{}
\saintindex{}
\feast{0812}{S. Claræ Virginis}{12 août}
\saintindex{Claræ}
\saintindex{Claire}
\feast{0812}{Tertia die infra Octavam S. Laurentii Martyris}{12 août}
\saintindex{}
\saintindex{}
\feast{0813}{Ss. Hippolyti et Cassiani Martyrum}{13 août}
\saintindex{}
\saintindex{}
\oratio{Da, quǽsumus, omnípotens Deus: ut beatórum Mártyrum tuórum Hippólyti et Cassiáni veneránda solémnitas, et devotiónem nobis áugeat, et salútem. \perdominumla}{Faites, s’il Vous plaît, ô Dieu tout puissant, que la solennité vénérable de Vos bienheureux Martyrs Hippolyte et Cassien, nous apporte un accroissement de dévotion et de salut. \perdominumfr}
\feast{0813}{Ss Pontiáni Papæ et Hippólyti Mártyrum}{13 août}
\saintindex{}
\saintindex{}
\oratio{@Commune/C4::s/beáti N. Confessóris tui atque Pontíficis/ut beátorum Mártyrum tuórum Pontiáni et Hippólyti/
}{}
\feast{0813}{Quarta die infra Octavam S. Laurentii Martyris}{13 août}
\saintindex{}
\saintindex{}
\feast{0814}{}{14 août}
\saintindex{}
\saintindex{}
\oratio{Deus, qui virginálem aulam beátæ Maríæ, in qua habitáres, elígere dignátus es: da, quǽsumus; ut, sua nos defensióne munítos, jucúndos fácias suæ interésse festivitáti:
$Qui vivis}{Dieu, Vous avez daigné choisir le palais virginal de la bienheureuse Vierge Marie pour y habiter : faites, nous Vous en prions, qu’à l’abri de sa protection, nous puissions avec joie prendre part à sa fête.
$Qui vivis}
\feast{0814}{S. Eusebii Confessoris}{14 août}
\saintindex{}
\saintindex{}
\oratio{@Commune/C5:Oratio:s/N\./Eusébii/ s/étiam actiónes imitémur/per ejus ad te exémpla gradiámur/
}{Dieu, Vous nous réjouissez par la solennité annuelle du bienheureux Eusèbe, Votre Confesseur : faites, dans Votre bonté, qu’honorant sa naissance au ciel, nous marchions vers Vous en suivant ses exemples. \perdominumfr.
}
\feast{0814}{Quinta die infra Octavam S. Laurentii Martyris}{14 août}
\saintindex{}
\saintindex{}
\feast{0815}{}{15 août}
\saintindex{}
\saintindex{}
\oratio{Omnípotens sempitérne Deus, qui Immaculátam Vírginem Maríam, Fílii tui Genetrícem, córpore et ánima ad cæléstem glóriam assumpsísti: concéde, quǽsumus; ut ad supérna semper inténti, ipsíus glóriæ mereámur esse consórtes. \pereumdemla}{Dieu éternel et tout-puissant, Vous avez élevé, en son corps et en son âme, à la gloire du ciel, Marie, la Vierge immaculée, mère de Votre Fils : faites, nous Vous en prions, que, sans cesse tendus vers les choses d’en-haut, nous méritions d’avoir part à sa gloire. \pereumdemfr}
\feast{0815}{}{15 août}
\saintindex{}
\saintindex{}
\oratio{Famulórum tuórum, quǽsumus, Dómine, delíctis ignósce: ut qui tibi placére de áctibus nostris non valémus: Genitrícis Fílii tui Dómini nostri intercessióne salvémur:
$Qui tecum}{Nous supplions, Seigneur, de pardonner les fautes de vos serviteurs : et dans l’impuissance où nous sommes de vous plaire par nos propres mérites, accordez-nous d’être sauvés par l’intercession de la Mère de votre Fils, notre Seigneur, qui, étant Dieu, vit et règne.
$Qui tecum}
\feast{0816}{S. Joachim Confessoris, Patris B. M. V.}{16 août}
\saintindex{}
\saintindex{}
\oratio{Deus, qui præ ómnibus Sanctis tuis beátum Jóachim Genetrícis Fílii tui patrem esse voluísti: concéde, quǽsumus; ut, cujus festa venerámur, ejus quoque perpétuo patrocínia sentiámus. \pereumdemla}{Dieu, de préférence à tous Vos saints, vous avez choisi le bienheureux Joachim pour qu’il fût le père de la Mère de Votre Fils : accordez-nous, s’il Vous plaît, la grâce d’être constamment protégés par celui dont nous célébrons la fête. \pereumdemfr}
\feast{0816}{}{16 août}
\saintindex{}
\saintindex{}
\feast{0816}{Septima die infra Octavam S. Laurentii Martyris}{16 août}
\saintindex{}
\saintindex{}
\feast{0817}{S. Hyacinthi Confessoris}{17 août}
\saintindex{Hyacínthi}
\saintindex{Hyacinthe}
\feast{0817}{}{17 août}
\saintindex{}
\saintindex{}
\feast{0817}{In Octava S. Laurentii Martyris}{17 août}
\saintindex{}
\saintindex{}
\oratio{Excita, Domine, in Ecclésia tua Spíritum, cui beátus Lauréntius Levíta servívit: ut eódem nos repléti, studeámus amáre quod amávit, et ópere exercére quod dócuit. \perdominumejusdemla}{Seigneur, faites lever dans votre Église l’Esprit auquel obéit le bienheureux Lévite Laurent : afin que, remplis de ce même Esprit, nous nous appliquions à aimer ce qu’il a aimé et à pratiquer ce qu’il a enseigné. \perdominumejusdemfr}
\feast{0818}{}{18 août}
\saintindex{}
\saintindex{}
\feast{0818}{S. Agapiti Martyris}{18 août}
\saintindex{}
\saintindex{}
\oratio{Lætétur Ecclésia tua, Deus, beáti Agapíti Mártyris tui confísa suffrágiis: atque, ejus précibus gloriósis, et devóta permáneat et secúra consístat. \perdominumla}{Que votre Église, Seigneur, se réjouisse de la confiance qu’elle a placée dans l’intercession du bienheureux Agapit, Votre Martyr : afin que, par ses glorieuses prières, elle puisse persévérer dans Votre service et jouir d’une paisible sécurité. \perdominumfr}
\feast{0819}{S. Joannis Eudes Confessoris}{19 août}
\saintindex{}
\saintindex{}
\oratio{Deus, qui beátum Joánnem, Confessórem tuum, ad cultum sacrórum Córdium Jesu et Maríæ rite promovéndum, mirabíliter inflammásti, et per eum novas in Ecclésia tua famílias congregáre voluísti: præsta, quǽsumus; ut, cujus pia mérita venerámur, virtútum quoque instruámur exémplis. \pereumdemla}{Dieu, Vous avez admirablement enflammé le bienheureux Jean, Votre Confesseur, de zèle à promouvoir dignement le culte des Sacrés-Cœurs de Jésus et de Marie et Vous avez voulu rassembler par lui dans Votre Église de nouvelles familles religieuses : accordez-nous, nous Vous en prions, de retirer un enseignement de ses exemples après avoir rendu hommage à ses pieux mérites. \pereumdemfr}
\feast{0819}{}{19 août}
\saintindex{}
\saintindex{}
\feast{0820}{S. Bernardi Abbatis et Ecclesiæ Doctoris}{20 août}
\saintindex{Bernárdum}
\saintindex{Bernard}
\feast{0820}{}{20 août}
\saintindex{}
\saintindex{}
\feast{0821}{S. Joannæ Franciscæ Frémiot de Chantal Viduæ}{21 août}
\saintindex{}
\saintindex{}
\oratio{Omnípotens et miséricors Deus, qui beátam Joánnam Francíscam, tuo amóre succénsam, admirábili spíritus fortitúdine per omnes vitæ sémitas in via perfectiónis donásti, quique per illam illustráre Ecclésiam tuam nova prole voluísti: ejus méritis et précibus concéde; ut, qui infirmitátis nostræ cónscii de tua virtúte confídimus, cæléstis grátiæ auxílio cuncta nobis adversántia vincámus. \perdominumla}{Dieu tout-puissant et miséricordieux, après avoir embrasé de Votre amour la bienheureuse Jeanne Françoise, Vous lui avez donné la force d’âme admirable qui la fit avancer en perfection dans toutes les situations de la vie et Vous avez voulu orner par elle Votre Église d’une nouvelle famille religieuse : faites, en considération de ses mérites et de ses prières, que, conscients de notre faiblesse, mais confiants en Votre secours, nous puissions, à l’aide de la grâce céleste, surmonter tout ce qui nous est contraire. \perdominumfr}
\feast{0821}{}{21 août}
\saintindex{}
\saintindex{}
\feast{0822}{}{22 août}
\saintindex{}
\saintindex{}
\oratio{Omnípotens sempitérne Deus, qui in Corde beátæ Maríæ Vírginis dignum Spíritus Sancti habitáculum præparásti: concéde propítius; ut ejúsdem immaculáti Cordis festivitátem devóta mente recoléntes, secúndum Cor tuum vívere valeámus. \perdominumejusdemla}{Dieu éternel et tout puissant, qui avez préparé dans le Cœur de la bienheureuse Vierge Marie une demeure digne du Saint-Esprit, faites, dans Votre bonté, qu’en célébrant de toute notre âme cette fête en l’honneur de son cœur immaculé, nous arrivions à vivre selon Votre cœur. \perdominumejusdemfr}
\feast{0822}{Ss. Timothei, Hippolyti et Symphoriani Martyres}{22 août}
\saintindex{}
\saintindex{}
\oratio{Auxílium tuum nobis, Dómine, quǽsumus, placátus impende: et, intercedéntibus beátis Martýribus tuis Timótheo, Hippólyto et Symphoriáno, déxteram super nos tuæ propitiatiónis exténde. \perdominumla}{Dans votre bonté, Seigneur, accordez-nous votre secours, et grâce à l’intercession de vos bienheureux Martyrs Timothée, Hippolyte et Symphorien, étendez sur nous votre bras protecteur. \perdominumfr}
\feast{0822}{}{22 août}
\saintindex{}
\saintindex{}
\feast{0823}{S. Philippi Benitii Confessoris}{23 août}
\saintindex{}
\saintindex{}
\oratio{Deus, qui per beátum Philíppum Confessórem tuum, exímium nobis humilitátis exémplum tribuísti: da fámulis tuis próspera mundi ex ejus imitatióne despícere, et cæléstia semper inquírere. \perdominumla}{Dieu, Vous nous avez donné un excellent modèle d’humilité en la personne de Votre Confesseur, le bienheureux Philippe : accordez à Vos serviteurs de mépriser, à son exemple, les biens de ce monde et de chercher toujours les biens du ciel. \perdominumfr}
\feast{0823}{In Vigilia S.Bartholomæi Apostoli}{23 août}
\saintindex{Bartholomǽi}
\saintindex{Barthélemy}
\feast{0824}{S. Bartholomæi Apostoli}{24 août}
\saintindex{}
\saintindex{}
\oratio{Omnípotens sempitérne Deus, qui hujus diéi venerándam sanctámque lætítiam in beáti Apóstoli tui Bartholomǽi festivitáte tribuísti: da Ecclésiæ tuæ, quǽsumus; et amáre quod crédidit, et prædicáre quod dócuit. \perdominumla}{Dieu tout-puissant et éternel, de qui nous vient la religieuse et sainte joie que nous éprouvons à célébrer aujourd’hui la fête de Votre bienheureux Apôtre Barthélemy, accordez à Votre Église, nous Vous en prions, la grâce d’aimer ce qu’il a cru et de prêcher ce qu’il a enseigné. \perdominumfr}
\feast{0825}{S. Ludovici Regis Franciæ Confessoris}{25 août}
\saintindex{}
\saintindex{}
\oratio{Deus, qui beátum Ludovícum Confessórem tuum de terréno regno ad cæléstis regni glóriam transtulísti: ejus, quǽsumus, méritis et intercessióne; Regis regum Jesu Christi, Fílii tui, fácias nos esse consórtes:
$Qui tecum}{Dieu, d’une royauté terrestre, Vous avez élevé saint Louis, Votre Confesseur, à la gloire du céleste royaume : daignez, nous Vous en prions, nous accorder, en considération de ses mérites et de son intercession, la grâce d’être associés à la gloire du Roi des rois, Jésus-Christ Votre Fils,
$Qui tecum}
\feast{0826}{S. Zephyrini Papæ et Martyris}{26 août}
\saintindex{Zephyrínum}
\saintindex{Zephyrin}
\oratio{Præsta, quǽsumus, omnípotens Deus: ut beáti Zephyríni, Mártyris tui atque Pontíficis, cujus gaudémus méritis, instruámur exémplis. \perdominumla}{Dieu tout puissant, accordez-nous, nous Vous en prions : puisque nous nous réjouissons des mérites du bienheureux Zéphyrin, Votre Martyr et Pontife, que nous soyons instruits de ses exemples. \perdominumfr}
\feast{0827}{S. Josephi Calasanctii Confessoris}{27 août}
\saintindex{}
\saintindex{}
\oratio{Deus, qui per sanctum Joséphum Confessórem tuum, ad erudiéndam spíritu intellegéntiæ ac pietátis juventútem, novum Ecclésiæ tuæ subsídium providére dignátus es: præsta, quǽsumus; nos, ejus exémplo et intercessióne, ita fácere et docére, ut prǽmia consequámur ætérna. \perdominumla}{Dieu, par saint Joseph, Votre Confesseur, Vous avez daigné pourvoir Votre Église d’un nouveau secours pour former la jeunesse à la science et à la piété : faites, nous Vous en prions, qu’aidés de son exemple et de son intercession, nous puissions agir et enseigner de manière à mériter les éternelles récompenses. \perdominumfr}
\feast{0828}{S. Augustini Episcopi et Confessoris et Ecclesiæ Doctoris}{28 août}
\saintindex{Augustíno}
\saintindex{Augustin}
\oratio{Adésto supplicatiónibus nostris, omnípotens Deus: et, quibus fidúciam sperándæ pietátis indúlges, intercedénte beáto Augustíno Confessóre tuo atque Pontífice, consuétæ misericórdiæ tríbue benígnus efféctum. \perdominumla}{Recevez favorablement nos supplications, Dieu tout-puissant et, puisque Vous voulez bien nous permettre d’espérer en Votre bonté, daignez, grâce à l’intercession du bienheureux Augustin, Votre Confesseur et Pontife, nous accorder les effets de Votre miséricorde habituelle. \perdominumfr}
\feast{0828}{S. Hermetis Martyris}{28 août}
\saintindex{}
\saintindex{}
\oratio{Deus, qui beátum Hermétem Mártyrem tuum virtúte constántiæ in passióne roborásti: ex ejus nobis imitatióne tríbue; pro amóre tuo próspera mundi despícere, et nulla ejus advérsa formidáre. \perdominumla}{Dieu, Vous avez donné à Votre Martyr le bienheureux Hermès, le courage et la constance au milieu des supplices : faites qu’à son exemple et pour Votre amour, nous méprisions les faveurs du monde et ne redoutions aucune adversité. \perdominumfr}
\feast{0829}{In Decollatione S. Joannis Baptistæ}{29 août}
\saintindex{}
\saintindex{}
\oratio{Sancti Joánnis Baptístæ Præcursóris et Mártyris tui, quǽsumus, Dómine, veneránda festívitas: salutáris auxílii nobis præstet efféctum:
$Qui vivis}{Nous Vous en prions, Seigneur, faites que la fête solennelle de Votre saint Précurseur et Martyr Jean-Baptiste, nous procure des grâces efficaces de salut.
$Qui vivis}
\feast{0829}{S. Sabinæ Martyris}{29 août}
\saintindex{Sabínæ}
\saintindex{Sabine}
\oratio{}{Dieu, entre autres merveilles de Votre puissance, Vous avez fait remporter la victoire du martyre même par le sexe le plus faible : faites, dans Votre bonté, qu’honorant la naissance au ciel de la bienheureuse Sabine, Votre Martyre, nous tendions vers Vous par l’imitation de ses exemples. \perdominumfr.}
\feast{0830}{S. Rosæ a Sancta Maria Limanæ Virginis}{30 août}
\saintindex{}
\saintindex{}
\oratio{Bonórum ómnium largítor, omnípotens Deus, qui beátam Rosam, cæléstis grátiæ rore prævéntam, virginitátis et patiéntiæ decóre Indis floréscere voluísti: da nobis fámulis tuis; ut, in odórem suavitátis ejus curréntes, Christi bonus odor éffici mereámur:
$Qui tecum}{Dieu tout-puissant, dispensateur de tous les biens, Vous avez prévenu de la rosée céleste de Votre grâce la bienheureuse Rose et Vous l’avez fait briller dans les Indes de tout l’éclat de la pureté et de la patience : accordez à Vos serviteurs de courir à l’odeur de ses parfums, afin qu’ils méritent de devenir eux-mêmes la bonne odeur de Jésus-Christ.
$Qui tecum}
\feast{0830}{Ss. Felicis et Adaucti Martyrum}{30 août}
\saintindex{}
\saintindex{}
\oratio{Majestátem tuam, Dómine, súpplices exorámus: ut, sicut nos júgiter Sanctórum tuórum commemoratióne lætíficas; ita semper supplicatióne deféndas. \perdominumla}{Nous adressons, Seigneur, de suppliantes prières à Votre majesté, à l’effet d’obtenir que, nous donnant dans la fête de Vos saints un continuel sujet de joie, Vous nous fassiez aussi trouver dans leurs prières une perpétuelle assistance. \perdominumfr}
\feast{0831}{S. Raymundi Nonnati Confessoris}{31 août}
\saintindex{}
\saintindex{}
\oratio{Deus, qui in liberándis fidélibus tuis ab impiórum captivitáte beátum Raymúndum Confessórem tuum mirábilem effecísti: ejus nobis intercessióne concéde; ut, a peccatórum vínculis absolúti, quæ tibi sunt plácita, líberis méntibus exsequámur. \perdominumla}{Dieu, Vous avez rendu le bienheureux Raymond, Votre Confesseur, admirable par son dévouement pour délivrer Vos fidèles de la captivité des impies : accordez-nous, par son intercession, d’être délivrés des liens du péché et d’accomplir d’une âme libre ce qui Vous est agréable. \perdominumfr}
\feast{0901}{S. Ægidii Abbatis}{1 septembre}
\saintindex{Ægídii}
\saintindex{Gilles}
\feast{0901}{Ss. Duodecim Fratrum Mártyrum}{1 septembre}
\saintindex{}
\saintindex{}
\oratio{Fratérna nos, Dómine, Mártyrum tuórum coróna lætíficet: quæ et fídei nostræ prǽbeat increménta virtútum, et multíplici nos suffrágio consolétur. \perdominumla}{Faites, Seigneur, que la couronne fraternelle de Vos Martyrs soit une source de joie : qu’elle procure à notre foi une augmentation de force et qu’elle nous console par ces intercessions multiples. \perdominumfr}
\feast{0902}{S. Stephani Hungariæ Regis Confessoris}{2 septembre}
\saintindex{}
\saintindex{}
\oratio{Concéde, quǽsumus, Ecclésiæ tuæ, omnípotens Deus: ut beátum Stéphanum Confessórem tuum, quem regnántem in terris propagatórem hábuit, propugnatórem habére mereátur gloriósum in cælis. \perdominumla}{Nous vous en prions, ô Dieu tout-puissant, accordez à votre Église, qu’après avoir été propagée par le bienheureux Étienne, quand il régnait sur la terre, elle mérite de l’avoir pour défenseur, maintenant qu’il est glorieux dans le ciel. \perdominumfr}
\feast{0903}{S. Pii X Papæ Confessoris}{3 septembre}
\saintindex{}
\saintindex{}
\oratio{Deus, qui ad tuéndam cathólicam fidem, et univérsa in Christo instauránda sanctum Pium, Summum Pontíficem, cælésti sapiéntia et apostólica fortitúdine replevísti: concéde propítius; ut, ejus institúta et exémpla sectántes, prǽmia consequámur ætérna. \pereumdemla}{Ô Dieu qui, pour protéger la foi catholique et instaurer toute chose dans le Christ, avez rempli le Souverain Pontife saint Pie X de sagesse céleste et de force apostolique : accordez-nous favorablement de suivre son enseignement et ses exemples afin de parvenir aux biens éternels. \pereumdemfr}
\feast{0905}{S. Laurentii Justiniani Episcopi et Confessoris}{5 septembre}
\saintindex{Lauréntii}
\saintindex{Laurent}
\feast{0908}{}{8 septembre}
\saintindex{}
\saintindex{}
\oratio{Fámulis tuis, quǽsumus, Dómine, cæléstis grátiæ munus impertíre: ut, quibus beátæ Vírginis partus éxstitit salútis exórdium; Nativitátis ejus votíva solémnitas pacis tríbuat increméntum. \perdominumla}{Seigneur, nous vous prions d’accorder à Vos serviteurs le don de la grâce céleste ; et, comme l’enfantement de la bienheureuse Vierge a été le principe de leur salut, qu’ainsi la pieuse solennité de sa Nativité leur procure un accroissement de paix. \perdominumfr}
\feast{0908}{S. Hadriani Martyris}{8 septembre}
\saintindex{Hadriáni}
\saintindex{Adrien}
\oratio{}{Accordez, Dieu tout-puissant, à nous qui célébrons la naissance au ciel du bienheureux Adrien, Votre Martyr, la grâce d’être, par son intercession, fortifiés dans l’amour de Votre nom.}
\feast{0909}{S. Gorgonii Martyris}{9 septembre}
\saintindex{}
\saintindex{}
\oratio{Sanctus tuus, Dómine, Gorgónius sua nos intercessióne lætíficet: et pia fáciat solemnitáte gaudére. \perdominumla}{Seigneur, faites que saint Gorgon, Votre martyr, nous console par son intercession et nous fasse goûter la joie de cette pieuse solennité. \perdominumfr}
\feast{0909}{}{9 septembre}
\saintindex{}
\saintindex{}
\feast{0909}{}{9 septembre}
\saintindex{}
\saintindex{}
\feast{0909}{S. Petri Claver Confessoris}{9 septembre}
\saintindex{}
\saintindex{}
\oratio{Deus, qui, abréptos in servitútem Nigrítas ad agnitiónem tui nóminis vocatúrus, beátum Petrum mira in eis juvándis caritáte et patiéntia roborásti: ejus nobis intercessióne concéde; ut, quæ Jesu Christi sunt quæréntes, próximos ópere et veritáte diligámus. \pereumdemla}{}
\feast{0910}{S. Nicolai de Tolentino Confessoris}{10 septembre}
\saintindex{Nicolái}
\saintindex{Nicolas}
\oratio{}{Agréez, Seigneur, nos supplications offertes en la solennité du bienheureux Nicolas, Votre confesseur, en sorte que, n’ayant aucune confiance en notre jutice nous soyons aidés par les prières de celui qui a su Vous plaire. \perdominumfr}
\feast{0910}{}{10 septembre}
\saintindex{}
\saintindex{}
\feast{0911}{Ss. Proti et Hyacinthi Martyrum}{11 septembre}
\saintindex{}
\saintindex{}
\oratio{Beatórum Mártyrum tuórum Proti et Hyacínthi nos, Dómine, fóveat pretiósa conféssio: et pia júgiter intercéssio tueátur. \perdominumla}{Faites, Seigneur, que le magnifique témoignage de Vos bienheureux Martyrs Prote et Hyacinthe nous encourage et nous anime et que leur pieuse intercession nous protège toujours. \perdominumfr}
\feast{0911}{}{11 septembre}
\saintindex{}
\saintindex{}
\feast{0912}{}{12 septembre}
\saintindex{}
\saintindex{}
\oratio{Concéde, quǽsumus, omnípotens Deus: ut fidéles tui, qui sub sanctíssimæ Vírginis Maríæ Nómine et protectióne lætántur; ejus pia intercessióne, a cunctis malis liberéntur in terris, et ad gáudia ætérna perveníre mereántur in cælis. \perdominumla}{Nous Vous en prions, ô Dieu tout-puissant, accordez à Vos fidèles, qui mettent leur joie dans le nom et la protection de la très sainte Vierge Marie, d’être, par un effet de sa maternelle intercession, préservés de tous les maux sur la terre et d’arriver dans le ciel aux joies de l’éternité. \perdominumfr}
\feast{0912}{}{12 septembre}
\saintindex{}
\saintindex{}
\feast{0913}{}{13 septembre}
\saintindex{}
\saintindex{}
\feast{0914}{}{14 septembre}
\saintindex{}
\saintindex{}
\oratio{Deus, qui nos hodiérna die Exaltatiónis sanctæ Crucis ánnua solemnitáte lætíficas: præsta, quǽsumus; ut, cujus mystérium in terra cognóvimus, ejus redemptiónis prǽmia in cælo mereámur. \pereumdemla}{Ô Dieu, qui nous donnez aujourd’hui un sujet de joie dans la fête annuelle de l’Exaltation de la sainte Croix, faites, nous Vous en prions, que nous méritions de recueillir dans le ciel les récompenses acquises au moyen de la rédemption de Celui dont nous avons connu le mystère ici-bas. \pereumdemfr}
\feast{0914}{}{14 septembre}
\saintindex{}
\saintindex{}
\feast{0915}{}{15 septembre}
\saintindex{}
\saintindex{}
\oratio{Deus, in cujus passióne, secúndum Simeónis prophetíam, dulcíssimam ánimam gloriósæ Vírginis et Matris Maríæ dolóris gládius pertransívit: concéde propítius; ut, qui dolóres ejus venerándo recólimus, passiónis tuæ efféctum felícem consequámur:
$Qui vivis}{Ô Dieu, dans la passion Duquel, suivant la prophétie de Siméon, un glaive de douleur a percé le cœur très doux de la glorieuse Vierge Marie, Votre Mère, faites, dans Votre miséricorde, que célébrant avec respect le souvenir de ses douleurs, nous recueillions les heureux fruits de Votre passion.
$Qui vivis}
\feast{0915}{S. Nicomedis Martyris}{15 septembre}
\saintindex{}
\saintindex{}
\oratio{v. Adésto, Dómine, pópulo tuo: ut, beáti Nicomédis Mártyris tui mérita præclára suscípiens, ad impetrándam misericórdiam tuam semper ejus patrocíniis adjuvétur. \perdominumla}{v. Montrez-Vous favorable à Votre peuple, Seigneur, afin que, célébrant les mérites si glorieux de Votre bienheureux Martyr Nicomède, il soit toujours aidé de ses prières pour obtenir Vos miséricordes. \perdominumfr}
\feast{0915}{}{15 septembre}
\saintindex{}
\saintindex{}
\feast{0916}{Ss. Cornelii Papæ et Cypriani Episcopi, Martyrum}{16 septembre}
\saintindex{Cornélii et Cypriáni}
\saintindex{Corneille et Cyrprien}
\oratio{}{Nous Vous en prions, Seigneur, faites qu’en célébrant les fêtes de Vos bienheureux Martyrs et Pontifes Corneille et Cyrprien, nous obtenions leur protection et que leur sainte prière nous serve de recommandation auprès de Vous. \perdominumfr}
\feast{0916}{Ss. Euphemiæ, Luciæ et Geminiani Martyrum}{16 septembre}
\saintindex{}
\saintindex{}
\oratio{Præsta, Dómine, précibus nostris cum exsultatióne provéntum: ut sanctórum Mártyrum Euphémiæ, Lúciæ et Geminiáni, quorum diem passiónis ánnua devotióne recólimus, étiam fídei constántiam subsequámur. \perdominumla}{v. Touché de nos prières, faites, Seigneur, que cette Fête nous profite ; en sorte que, célébrant avec dévotion, chaque année, le jour où Vos saints Martyrs Euphémie, Lucie et Oéminien ont souffert, nous les imitions dans la constance de la foi. \perdominumfr}
\feast{0917}{Impressionis Stigmatum S. Francisci}{17 septembre}
\saintindex{}
\saintindex{}
\oratio{Dómine Jesu Christe, qui, frigescénte mundo, ad inflammándum corda nostra tui amóris igne, in carne beatíssimi Francísci passiónis tuæ sacra Stígmata renovásti: concéde propítius; ut ejus méritis et précibus crucem júgiter ferámus, et dignos fructus pæniténtiæ faciámus:
$Qui vivis}{Seigneur Jésus-Christ, qui, lorsque le monde se refroidissait, avez voulu, pour enflammer nos cœurs du feu de Votre amour, renouveler les sacrés stigmates de Votre passion dans la chair du bienheureux François, accordez-nous, s’il Vous plaît, que, par ses mérites et ses prières, nous portions continuellement la croix et que nous fassions de dignes fruits de pénitence.
$Qui vivis}
\feast{0918}{S. Josephi de Cupertino Confessoris}{18 septembre}
\saintindex{}
\saintindex{}
\oratio{Deus, qui ad unigénitum Fílium tuum exaltátum a terra ómnia tráhere disposuísti: pérfice propítius; ut, méritis et exémplo seráphici Confessóris tui Joséphi, supra terrénas omnes cupiditátes eleváti, ad eum perveníre mereámur:
$Qui tecum}{Ô Dieu, qui, après que Votre Fils unique eut été élevé de terre, avez voulu attirer tout à Lui, faites, dans Votre miséricorde, que, nous élevant au-dessus de tous les désirs terrestres, à l’exemple et par les mérites de Votre séraphique Confesseur Joseph, nous méritions d’arriver auprès de
$Qui tecum}
\feast{0919}{S. Januarii Episcopi et Sociorum Martyrum}{19 septembre}
\saintindex{Januárii et Sociórum ejus}
\saintindex{Janvier et ses compagnons}
\feast{0920}{Ss. Eustachii et Sociorum Martyrum}{20 septembre}
\saintindex{Eustáchii et Sociórum ejus}
\saintindex{Janvier et ses compagnons}
\oratio{}{Ô Dieu qui nous faites la grâce de célébrer la naissance au ciel de Vos saints Martyrs, Eustache et ses compagnons, accordez-nous de jouir de leur société, dans l'éternelle béatitude. \perdominumfr.}
\feast{0920}{S. Andrew Kim Tægon, priest, and Paul Chong Hasang and companions, martyrs}{20 septembre}
\saintindex{}
\saintindex{}
\feast{0920}{In Vigilia S. Matthæi Ap.}{20 septembre}
\saintindex{Matthǽi}
\saintindex{Matthieu}
\oratio{@Commune/C1v:Oratio pro Evangelistae}{}
\feast{0921}{S. Matthæi Apostoli et Evangelistæ}{21 septembre}
\saintindex{}
\saintindex{}
\oratio{Beáti Apóstoli et Evangelístæ Matthǽi, Dómine, précibus adjuvémur: ut, quod possibílitas nostra non óbtinet, ejus nobis intercessióne donétur. \perdominumla}{Que les prières du bienheureux Apôtre et Évangéliste Matthieu nous viennent en aide, Seigneur, afin que les grâces que notre insuffisance ne peut obtenir nous soient accordées par son intercession. \perdominumfr}
\feast{0922}{S. Thomæ de Villanova Episcopi et Confessoris}{22 septembre}
\saintindex{}
\saintindex{}
\oratio{Deus, qui beátum Thomam Pontíficem insígnis in páuperes misericórdiæ virtúte decorásti: quǽsumus; ut, ejus intercessióne, in omnes, qui te deprecántur, divítias misericórdiæ tuæ benígnus effúndas. \perdominumla}{Ô Dieu, qui avez enrichi et illustré le bienheureux Pontife Thomas d’une insigne compassion envers les pauvres, faites, nous Vous en prions, que son intercession obtienne de Votre bonté, pour tous ceux qui Vous implorent, l’effusion des trésors de Votre miséricorde. \perdominumfr}
\feast{0922}{Ss. Mauritii et Sociorum Martyrum}{22 septembre}
\saintindex{}
\saintindex{}
\oratio{Annue, quǽsumus, omnípotens Deus: ut sanctórum Mártyrum tuórum Maurítii et Sociórum ejus nos lætíficet festíva solémnitas; ut, quorum suffrágiis nítimur, eórum natalíciis gloriémur. \perdominumla}{Qu’il Vous plaise, Dieu tout-puissant, que la fête solennelle de Vos saints Martyrs, Maurice et ses compagnons, nous procure la joie, afin qu’ayant l’appui de leurs prières, nous participions à la gloire de leur naissance au ciel. \perdominumfr}
\feast{0923}{S. Lini Papæ et Martyris}{23 septembre}
\saintindex{Lini}
\saintindex{Lin}
\feast{0923}{S. Theclæ Virginis et Martyris}{23 septembre}
\saintindex{Theclæ}
\saintindex{Thècle}
\oratio{Da, quǽsumus, omnípotens Deus: ut, qui beátæ Theclæ Vírginis et Mártyris tuæ natalícia cólimus; et ánnua solemnitáte lætémur, et tantæ fídei proficiámus exémplo. \perdominumla}{Accordez-nous, s’il Vous plaît, Dieu tout-puissant, que célébrant la naissance au ciel de la bienheureuse Thècle, Votre Vierge et Martyre, nous retirions une sainte joie de cette fête annuelle et profitions de l’exemple que nous donne sa foi si grande. \perdominumfr}
\feast{0923}{S. Pio of Pietrelcina (Padre Pio), priest}{23 septembre}
\saintindex{}
\saintindex{}
\feast{0924}{}{24 septembre}
\saintindex{}
\saintindex{}
\oratio{Deus, qui per gloriosíssimam Fílii tui Matrem, ad liberándos Christi fidéles a potestáte paganórum nova Ecclésiam tuam prole amplificáre dignátus es: præsta, quǽsumus; ut, quam pie venerámur tanti operis institutrícem, ejus páriter méritis et intercessióne, a peccátis ómnibus et captivitáte dǽmonis liberémur. \pereumdemla}{Ô Dieu, qui, par la très glorieuse Mère de Votre Fils, avez daigné enrichir Votre Église d’une nouvelle famille destinée à délivrer les fidèles du Christ de la puissance des païens, faites, nous Vous prions, que, vénérant avec piété l’inspiratrice d’une si grande œuvre, nous soyons, grâce à ses mérites et son intercession, délivrés de nos péchés et de la captivité du démon. \pereumdemfr}
\feast{0926}{Ss. Cypriani et Justinæ Virginis, Martyrum}{26 septembre}
\saintindex{}
\saintindex{}
\oratio{Beatórum Mártyrum Cypriáni et Justínæ nos, Dómine, fóveant continuáta præsídia: quia non désinis propítius intuéri, quos tálibus auxíliis concésseris adjuvári. \perdominumla}{Faites, Seigneur, que Vos bienheureux Martyrs Cyprien et Justine, nous couvrent de leur protection continuelle, puisque Vous ne cessez de regarder favorablement ceux à qui Vous accordez l’assistance de tels secours. \perdominumfr}
\feast{0926}{Ss. Isaaci Jogues, Joannis de Brébeuf et Sociorum Martyrum}{26 septembre}
\saintindex{}
\saintindex{}
\oratio{Deus, qui primítias fídei in boreálibus Américæ regiónibus sanctórum Mártyrum tuórum Isaáci, Joánnis, eorúmque Sociórum prædicatióne et sánguine consecrásti: concéde propítius; ut, eórum intercessióne, flórida christianórum seges ubíque in dies augeátur.  \perdominumla}{Ô Dieu, qui avez consacré les prémices de la foi dans les régions d'Amérique du Nord par la prédication et le sang de vos saints martyrs Jean, Isaac et de leurs compagnons, accordez-nous, nous vous en prions, que par leur intercession, grandisse partout de jour en jour une abondante moisson de chrétiens. \perdominumfr
}
\feast{0927}{Ss. Cosmæ et Damiani Martyrum}{27 septembre}
\saintindex{}
\saintindex{}
\oratio{Præsta, quǽsumus, omnípotens Deus: ut, qui sanctórum Mártyrum tuórum Cosmæ et Damiáni natalícia cólimus, a cunctis malis imminéntibus, eórum intercessiónibus, liberémur. \perdominumla}{Accordez-nous, s’il Vous plaît, ô Dieu tout-puissant, que, célébrant la naissance au ciel de Vos bienheureux Martyrs Côme et Damien, nous soyons délivrés, grâce à leur intercession, de tous les maux qui nous menacent. \perdominumfr}
\feast{0928}{S. Wenceslai Ducis et Martyris}{28 septembre}
\saintindex{}
\saintindex{}
\oratio{Deus, qui beátum Wencesláum per martýrii palmam a terréno principátu ad cæléstem glóriam transtulísti: ejus précibus nos ab omni adversitáte custódi; et ejúsdem tríbue gaudére consórtio. \perdominumla}{Ô Dieu, qui, par le triomphe du martyre, avez fait passer le bienheureux Wenceslas d’une principauté terrestre à la gloire du ciel, accordez-nous, grâce à l’intercession de ses prières, d’être préservés de toute adversité et de partager son sort glorieux. \perdominumfr}
\feast{0929}{In Dedicatione S. Michaëlis Archangelis}{29 septembre}
\saintindex{}
\saintindex{}
\feast{0930}{S. Hieronymi Presbyteris Confessoris et Ecclesiæ Doctoris}{30 septembre}
\saintindex{Hierónymum}
\saintindex{Jérôme}
\oratio{Deus, qui Ecclésiæ tuæ in exponéndis sacris Scriptúris beátum Hierónymum, Confessórem tuum, Doctórem máximum providére dignátus es: præsta, quǽsumus; ut, ejus suffragántibus méritis, quod ore simul et ópere dócuit, te adjuvánte, exercére valeámus. \perdominumla}{Ô Dieu, qui avez ménagé à votre Église pour expliquer les saintes Écritures, un éminent Docteur dans la personne du bienheureux Jérôme, Votre Confesseur, faites, nous Vous en prions, que, secondés par ses mérites, nous puissions, Votre grâce aidant, pratiquer ce qu’il a enseigné tout à la fois au moyen de la parole et de l’action. \perdominumfr}
\feast{1001}{S. Remigii Episcopi et Confessoris}{1 octobre}
\saintindex{Remígii}
\saintindex{Remi}
\feast{1002}{Ss. Angelorum Custodum}{2 octobre}
\saintindex{}
\saintindex{}
\oratio{Deus, qui ineffábili providéntia sanctos Angelos tuos ad nostram custódiam míttere dignáris: largíre supplícibus tuis; et eórum semper protectióne deféndi, et ætérna societáte gaudére. \perdominumla}{Ô Dieu, qui, en votre providence ineffable, daignez envoyer vos saints Anges pour nous garder, accordez à ceux qui vous en supplient, d’être sans cesse sous leur défense et leur protection, et de jouir éternellement de leur société. \perdominumfr}
\feast{1003}{S. Theresiæ a Jesu Infante Virginis}{3 octobre}
\saintindex{}
\saintindex{}
\oratio{Dómine, qui dixísti: Nisi efficiámini sicut párvuli, non intrábitis in regnum cælórum: da nobis, quǽsumus: ita sanctæ Terésiæ Vírginis in humilitáte et simplicitáte cordis vestígia sectári, ut prǽmia consequámur ætérna.
$Qui vivis}{Seigneur qui avez dit : Si vous ne devenez semblables à de petits enfants, vous n’entrerez pas dans le Royaume des Cieux, donnez-nous, nous vous en supplions, de suivre les traces de votre sainte Vierge Thérèse dans la voie de l’humilité et de la simplicité du cœur, en sorte que nous méritions de partager sa récompense éternelle.
$Qui vivis}
\feast{1004}{S. Francisci Confessoris}{4 octobre}
\saintindex{}
\saintindex{}
\oratio{Deus, qui Ecclésiam tuam beáti Francísci méritis fœtu novæ prolis amplíficas: tríbue nobis; ex ejus imitatióne, terréna despícere, et cæléstium donórum semper participatióne gaudére. \perdominumla}{Ô Dieu, qui, par les mérites du bienheureux François, avez enrichi votre Église, en lui donnant une nouvelle famille, faites-nous la grâce de l’imiter en méprisant les biens de la terre, et d’avoir la joie de participer toujours aux dons célestes. \perdominumfr}
\feast{1005}{Ss. Placidi et Sociorum Martyrum}{5 octobre}
\saintindex{Plácidi et Sociórum ejus}
\saintindex{Placide et ses Compagnons}
\oratio{@Commune/C3a}{Ô Dieu qui nous faites la grâce d’honorer la naissance au ciel de vos Saints Martyrs Placide et ses Compagnons, accordez-nous de jouir de leur société dans l’éternité bienheureuse. \perdominumfr.}
\feast{1006}{S. Brunonis Confessoris}{6 octobre}
\saintindex{}
\saintindex{}
\oratio{Sancti Brunónis Confessóris tui, quǽsumus, Dómine, intercessiónibus adjuvémur: ut, qui majestátem tuam gráviter delinquéndo offéndimus, ejus méritis et précibus, nostrórum delictórum véniam consequámur. \perdominumla}{Que l’intercession de votre Confesseur saint Bruno nous soit en aide, nous vous en prions, Seigneur, afin que ses mérites et ses prières nous obtiennent le pardon des péchés, par lesquels nous avons gravement offensé votre majesté. \perdominumfr}
\feast{1007}{}{7 octobre}
\saintindex{}
\saintindex{}
\oratio{Deus, cujus Unigénitus per vitam, mortem et resurrectiónem suam nobis salútis ætérnæ prǽmia comparávit: concéde, quǽsumus; ut hæc mystéria sacratíssimo beátæ Maríæ Vírginis Rosário recoléntes, et imitémur quod cóntinent, et quod promíttunt, assequámur. \pereumdemla}{Ô Dieu, dont le Fils unique nous a obtenu le prix du salut éternel par sa vie, sa mort et sa résurrection ; faites, nous vous en prions, qu’honorant ces mystères au moyen du très saint Rosaire de la bienheureuse Vierge Marie, nous imitions ce qu’ils contiennent et obtenions ce qu’ils promettent. \pereumdemfr}
\feast{1007}{Ss. Sergii, Bacchi, Marcelli et Apuleji Martyrum}{7 octobre}
\saintindex{}
\saintindex{}
\oratio{Sanctórum Mártyrum tuórum nos, Dómine, Sergii, Bacchi, Marcelli et Apuleji beáta merita prosequántur: et tuo semper fáciant amóre fervéntes. \perdominumla}{Faites, Seigneur, que les mérites de vos saints Martyrs Serge, Bacchus, Marcel et Apulée, nous soient acquis et nous rendent toujours fervents dans votre amour. \perdominumfr
}
\feast{1007}{S. Marci Papæ Confessoris}{7 octobre}
\saintindex{Marco}
\saintindex{Marc}
\oratio{Exáudi, Dómine, preces nostras: et, interveniénte beáto Marco, Confessóre tuo atque Pontífice, indulgéntiam nobis tríbue placátus et pacem. \perdominumla}{Exaucez, Seigneur, nos prières, et par l'intercession du bienheureux Marc, votre Confesseur et Pontife, daignez nous accorder l'indulgence et la paix. \perdominumfr}
\feast{1008}{S. Birgittæ Viduæ}{8 octobre}
\saintindex{}
\saintindex{}
\oratio{Dómine, Deus noster, qui beátæ Birgíttæ per Fílium tuum unigénitum secréta cæléstia revelásti: ipsíus pia intercessióne da nobis fámulis tuis; in revelatióne sempitérnæ glóriæ tuæ gaudére lætántes. \pereumdemla}{Seigneur notre Dieu, qui avez révélé par votre Fils unique, à la bienheureuse Brigitte, les secrets célestes, faites que, par sa pieuse intercession, nous qui sommes vos serviteurs, nous jouissions dans l’éternelle félicité de la manifestation de votre gloire. \pereumdemfr}
\feast{1009}{S. Joannis Leonardi Confessoris}{9 octobre}
\saintindex{}
\saintindex{}
\oratio{Deus, qui beátum Joánnem Confessórem tuum ad fidem in géntibus propagándam mirabíliter excitáre dignátus es, ac per eum in erudiéndis fidélibus novam in Ecclésia tua famíliam congregásti: da nobis fámulis tuis; ita ejus institútis profícere, ut prǽmia consequámur ætérna. \perdominumla}{Ô Dieu, afin de répandre la foi parmi les païens, vous avez suscité d’une merveilleuse façon le bienheureux Jean votre Confesseur, et vous lui avez fait fonder dans votre Église une nouvelle famille pour instruire les fidèles : donnez-nous, à nous vos serviteurs, de profiter si bien de ses enseignements que nous obtenions la vie éternelle. \perdominumfr}
\feast{1009}{S. Dionysii Epíscopi, Rustici et Eleutherii Martyrum}{9 octobre}
\saintindex{}
\saintindex{}
\oratio{Deus, qui hodiérna die beátum Dionysium, Mártyrem tuum atque Pontificem, virtúte constantiæ in passióne roborásti, quique illi, ad prædicándum géntibus glóriam tuam, Rusticum et Eleuthérium sociare dignátus es: tríbue nobis, quǽsumus: eórum imitatióne, pro amóre tuo prospera mundi despícere, et nulla ejus adversa formidare. \perdominumla}{Ô Dieu, qui, en ce jour, avez fortifié le bienheureux Denis, votre Martyr et Pontife, lui donnant la constance dans l’épreuve du martyre, et qui avez daigné lui associer Rustique et Éleuthère, pour annoncer votre gloire aux Gentils ; faites-nous, s’il vous plaît, la grâce de mépriser, à leur exemple et pour l’amour de vous, les prospérités du monde et de ne craindre aucune de ses adversités. \perdominumfr}
\feast{1010}{S. Francisci Borgiæ Confessoris}{10 octobre}
\saintindex{}
\saintindex{}
\oratio{Dómine Jesu Christe, veræ humilitátis et exémplar et prǽmium: quǽsumus; ut, sicut beátum Francíscum in terréni honóris contémptu imitatórem tui gloriósum effecísti, ita nos ejúsdem imitatiónis et glóriæ tríbuas esse consórtes:
$Qui vivis}{Seigneur Jésus-Christ, vous qui êtes le modèle et la récompense de la véritable humilité, et qui avez fait du bienheureux François votre glorieux imitateur dans le mépris des honneurs terrestres, accordez-nous la grâce de l’imiter et de partager sa gloire.
$Qui vivis}
\feast{1011}{}{11 octobre}
\saintindex{}
\saintindex{}
\oratio{Deus, qui de beátæ Maríæ Vírginis útero Verbum tuum, Angelo nuntiánte, carnem suscípere voluísti: præsta supplícibus tuis: ut qui vere eam Genetrícem Dei crédimus, ejus apud te intercessiónibus adjuvémur. \pereumdemla}{Ô Dieu, vous avez voulu qu’à l’annonce de l’Ange, le Verbe Votre Fils, prenne chair dans le sein de la bienheureuse Vierge Marie : accordez à ceux qui Vous supplient et qui croient qu’elle est vraiment la Mère de Dieu, qu’elle nous secoure par son intercession. \pereumdemfr}
\feast{1013}{S. Eduardi Regis Confessoris}{13 octobre}
\saintindex{}
\saintindex{}
\oratio{Deus, qui beátum regem Eduárdum Confessórem tuum æternitátis glória coronásti: fac nos, quǽsumus; ita eum venerári in terris, ut cum eo regnáre possímus in cælis. \perdominumla}{Ô Dieu, qui avez couronné de la gloire éternelle le roi Édouard, votre bienheureux Confesseur ; faites, s’il vous plaît, que nous l’honorions sur la terre de manière à mériter de régner avec lui dans les cieux. \perdominumfr}
\feast{1014}{S. Callisti Papæ et Martyris}{14 octobre}
\saintindex{}
\saintindex{}
\oratio{Deus, qui nos cónspicis ex nostra infirmitáte defícere: ad amórem tuum nos misericórditer per Sanctórum tuórum exémpla restáura. \perdominumla}{Ô Dieu, qui nous voyez défaillir à cause de notre faiblesse, raffermissez-nous miséricordieusement dans votre amour au moyen des exemples de vos Saints. \perdominumfr}
\feast{1015}{S. Teresiæ Virginis}{15 octobre}
\saintindex{Terésiæ}
\saintindex{Thérèse}
\oratio{Exáudi nos, Deus, salutáris noster: ut sicut de beátæ Terésiæ Vírginis tuæ festivitáte gaudémus; ita cæléstis ejus doctrínæ pábulo nutriámur, ita piæ devotiónis erudiámur afféctu.  \perdominumla}{Exaucez-nous, ô Dieu, qui êtes notre salut, et faites que, célébrant avec joie la fête de la bienheureuse Thérèse, Votre Vierge, nous soyons nourris du pain de sa céleste doctrine et formés aux sentiments d’une piété fervente. \perdominumfr}
\feast{1016}{S. Hedwigis Viduæ}{16 octobre}
\saintindex{}
\saintindex{}
\oratio{Deus, qui beátam Hedwígem a sǽculi pompa ad húmilem tuæ crucis sequélam toto corde transíre docuísti: concéde; ut ejus méritis et exémplo discámus peritúras mundi calcáre delícias, et in ampléxu tuæ crucis ómnia nobis adversántia superáre:
$Qui vivis}{Ô Dieu, de qui la bienheureuse Hedwige apprit à passer généreusement du sein des pompes du siècle en l’humble voie de votre croix ; faites que, par ses mérites et à son exemple, nous apprenions à fouler aux pieds les délices périssables du monde et à surmonter, en embrassant votre croix, tout ce qui nous est contraire.
$Qui vivis}
\oratio{Tuórum corda fidélium, Deus miserátor, illústra: et beátæ Hedwígis précibus gloriósis, fac nos próspera mundi despícere, et cœlésti semper consolatióne gaudére. \perdominumla}{}
\feast{1017}{S. Margaritæ Mariæ Alacoque Virginis}{17 octobre}
\saintindex{}
\saintindex{}
\oratio{Dómine Jesu Christe, qui investigábiles divítias Cordis tui beátæ Margarítæ Maríæ Vírgini mirabíliter revelásti; da nobis ejus méritis et imitatióne; ut te in ómnibus et super ómnia diligéntes, jugem in eódem Corde tuo mansiónem habére mereámur:
$Qui vivis.}{Seigneur Jésus-Christ, vous avez révélé miraculeusement à la bienheureuse vierge Marguerite Marie les richesses inépuisables de votre Cœur : que, par ses mérites et à son imitation, nous obtenions de vous aimer en tout et par dessus-tout, et d’établir en votre Cœur notre demeure permanente.
$Qui vivis.}
\feast{1018}{S. Lucæ Evangelistæ}{18 octobre}
\saintindex{}
\saintindex{}
\oratio{Intervéniat pro nobis, quǽsumus, Dómine, sanctus tuus Lucas Evangelísta: qui crucis mortificatiónem júgiter in suo córpore, pro tui nóminis honóre, portávit. \perdominumla}{Nous vous en prions, Seigneur, que votre saint Évangéliste Luc intercède pour nous, lui qui n’a jamais cessé de porter dans son corps la mortification de la croix, pour l’honneur de votre nom. \perdominumfr}
\feast{1019}{S. Petri de Alcantara Confessoris}{19 octobre}
\saintindex{}
\saintindex{}
\oratio{Deus, qui beátum Petrum Confessórem tuum admirábilis pœniténtiæ et altíssimæ contemplatiónis múnere illustráre dignátus es: da nobis, quǽsumus; ut, ejus suffragántibus méritis, carne mortificáti, facílius cæléstia capiámus.  \perdominumla}{Ô Dieu, qui avez daigné faire briller dans votre Confesseur, le bienheureux Pierre, les dons d’une admirable pénitence et d’une sublime contemplation, faites, s’il vous plaît, qu’aidés de ses mérites et mortifiant notre chair, nous obtenions plus facilement les biens célestes. \perdominumfr}
\feast{1019}{Ss. Isaaci Jogues, Joannis de Brebeuf et Sociorum Martyrum}{19 octobre}
\saintindex{}
\saintindex{}
\feast{1020}{S. Joannis Cantii Confessoris}{20 octobre}
\saintindex{}
\saintindex{}
\oratio{Da, quǽsumus, omnípotens Deus: ut, sancti Joánnis Confessóris exémplo in sciéntia Sanctórum proficiéntes, atque áliis misericórdiam exhibéntes; ejus méritis, indulgéntiam apud te consequámur. \perdominumla}{Nous vous en prions, Dieu tout-puissant, faites que, progressant dans la science des Saints et montrant de la compassion envers nos frères, à l’exemple du saint Confesseur Jean, nous puissions, grâce à ses mérites, trouver indulgence auprès de vous. \perdominumfr}
\feast{1021}{S. Hilarionis Abbatis}{21 octobre}
\saintindex{Hilariónis}
\saintindex{Hilarion}
\feast{1021}{Ss. Ursulæ et Sociarum Virg. et Mart.}{21 octobre}
\saintindex{Ursulæ et Sociárum}
\saintindex{Ursule et ses Compagnes}
\feast{1023}{S. Antonii Mariæ Claret Episcopi et Confessoris}{23 octobre}
\saintindex{}
\saintindex{}
\oratio{Deus, qui beátum Antónium Maríam Confessórem tuum atque Pontíficem, apostólicis virtútibus sublimásti, et per eum novas in Ecclésia clericórum ac vírginum famílias collegísti: concéde, quǽsumus; ut, ejus dirigéntibus mónitis ac suffragántibus méritis, animárum salútem quǽrere júgiter studeámus. \perdominumla}{Ô Dieu, qui avez élevé le bienheureux Antoine-Marie, votre Confesseur et Evêque, par les vertus apostoliques, et qui avez institué par lui dans votre Eglise de nouvelles familles de clercs et de vierges : faites que, dirigés par son enseignement et soutenus par ses mérites, nous puissions travailler sans cesse à chercher le salut des âmes. \perdominumfr}
\feast{1024}{S. Raphaëlis Archangeli}{24 octobre}
\saintindex{}
\saintindex{}
\oratio{Deus, qui beátum Raphaélem Archángelum Tobíæ fámulo tuo cómitem dedísti in via: concéde nobis fámulis tuis; ut ejúsdem semper protegámur custódia et muniámur auxílio. \perdominumla}{Ô Dieu, qui avez donné le bienheureux Archange Raphaël comme compagnon de route à votre serviteur Tobie, accordez-nous, à nous vos serviteurs, la grâce d’être toujours protégés et secourus par ce même Archange. \perdominumfr}
\feast{1025}{Ss. Chrysanthi et Dariæ Martyrum}{25 octobre}
\saintindex{}
\saintindex{}
\oratio{Beatórum Mártyrum tuórum, Dómine, Chrysánthi et Daríæ, quǽsumus, adsit nobis orátio: ut, quos venerámur obséquio, eórum pium júgiter experiámur auxílium. \perdominumla}{Que nous assiste la prière de vos saints Martyrs Chrysanthe et Darie, Seigneur, nous vous en prions : afin que nous éprouvions sans cesse le secours dévoué de ceux que nous vénérons par notre hommage \perdominumfr}
\feast{1025}{S. Isidori Agricolæ Confessoris}{25 octobre}
\saintindex{}
\saintindex{}
\oratio{Da nobis, quǽsumus, miséricors Deus, beáto Isidóro Agrícola Confessóre tuo intercedénte, supérbe non sápere: sed ejus méritis et exémplis plácita tibi semper humilitáte deservíre. \perdominumla}{}
\feast{1026}{S. Evaristi Papæ et Martyris}{26 octobre}
\saintindex{Evarísti}
\saintindex{Évariste}
\feast{1027}{In Vigilia Ss. Simonis et Judæ Ap.}{27 octobre}
\saintindex{}
\saintindex{}
\oratio{Concede, quǽsumus, omnípotens Deus: ut, sicut Apostolórum tuórum Simónis et Judæ gloriósa natalícia prævenimus; sic ad tua benefícia promerenda, majestátem tuam pro nobis ipsi præveniant. \perdominumla}{Accordez-nous, nous vous en prions, Dieu tout-puissant, alors que nous devançons le jour glorieux de la naissance au ciel de vos Apôtres Simon et Jude, qu’ainsi ils nous préviennent eux-mêmes auprès de votre majesté pour nous obtenir vos bienfaits. \perdominumfr
}
\feast{1028}{Ss. Simonis et Judæ Apostolorum}{28 octobre}
\saintindex{}
\saintindex{}
\oratio{Deus, qui nos per beátos Apóstolos tuos Simónem et Judam ad agnitiónem tui nóminis veníre tribuísti: da nobis eórum glóriam sempitérnam et proficiéndo celebráre, et celebrándo profícere. \perdominumla}{Ô Dieu, vous nous avez accordé la grâce de parvenir à la connaissance de votre nom par vos bienheureux Apôtres Simon et Jude : faites qu’en progressant nous célébrions leur gloire éternelle et en la célébrant nous progressions. \perdominumfr}
\feast{1031}{In Vigilia Omnium Sanctorum}{31 octobre}
\saintindex{}
\saintindex{}
\oratio{Dómine, Deus noster, multíplica super nos grátiam tuam: et, quorum prævenimus gloriósa solemnia, tríbue subsequi in sancta professióne lætítiam. \perdominumla}{Seigneur notre Dieu, multipliez envers nous l’effusion de votre grâce, et faites que, par une vie sainte, nous méritions de suivre dans la félicité éternelle ceux dont nous anticipons la fête solennelle. \perdominumfr
}
\feast{1031}{}{31 octobre}
\saintindex{}
\saintindex{}
\feast{1101}{Omnium Sanctorum}{1 novembre}
\saintindex{}
\saintindex{}
\oratio{Omnípotens sempitérne Deus, qui nos ómnium Sanctórum tuórum mérita sub una tribuísti celebritáte venerári: quǽsumus; ut desiderátam nobis tuæ propitiatiónis abundántiam, multiplicátis intercessóribus, largiáris. \perdominumla}{Dieu tout-puissant et éternel, qui nous avez accordé de célébrer dans une même solennité les mérites de tous vos Saints ; faites, nous vous en prions, que nos intercesseurs étant multipliés, une abondante effusion de vos miséricordes, objet de nos désirs, nous vienne de votre munificence. \perdominumfr}
\feast{1102}{}{2 novembre}
\saintindex{}
\saintindex{}
\oratio{&Dominus_vobiscum
@Commune/C9:Oratio_Fid}{#Oratio
@Commune/C9:Oratio_a_porta:s/&psalm.*\)//s
$Oremus
v. Faites miséricorde, nous vous le demandons, Seigneur, aux âmes de tous vos serviteurs et de toutes vos servantes, pour lesquelles nous prions et supplions votre majesté, afin que par ces offices de pieuse supplication, elles méritent d’arriver au repos éternel. Par Notre-Seigneur. \perdominumfr}
\oratio{#Oratio
@Commune/C9:Oratio 2}{#Oratio
@Commune/C9:Oratio_a_porta:s/&psalm.*\)//s
_
$Oremus
v. Nous, suppliants, Seigneur, répandons nos prières pour les âmes de tous vos serviteurs et servantes, afin que vous leur pardonniez miséricordieusement tout ce qu'ils ont commis par leur fragilité humaine, et que vous leur accordiez gracieusement de les délivrer des douleurs du purgatoire.  \perdominumfr}
\oratio{#Oratio
@Commune/C9:Oratio_a_porta:s/&psalm.*\)//s
$Oremus
v. Propitiáre, quǽsumus, Dómine, animábus ómnium famulórum famularúmque tuárum, pro quibus majestátem tuam supplíciter exorámus: ut, per hæc piæ deprecatiónis offícia, perveníre mereántur ad réquiem sempitérnam. \perdominumla}{}
\oratio{#Oratio
@Commune/C9:Oratio_a_porta:s/&psalm.*\)//s
_
$Oremus
v. Súpplices, Dómine, pro animábus ómnium famulórum famularúmque tuárum preces effúndimus: ut, quidquid conversatióne contraxérunt humána, clemens indúlgeas, et pœnis eórum finem benígnus impónas. \perdominumla}{}
\feast{1102}{Secunda die infra Octavam Omnium Sanctorum}{2 novembre}
\saintindex{}
\saintindex{}
\feast{1103}{S. Martin de Porres, religious}{3 novembre}
\saintindex{}
\saintindex{}
\feast{1103}{Tertia die infra Octavam Omnium Sanctorum}{3 novembre}
\saintindex{}
\saintindex{}
\feast{1103}{}{3 novembre}
\saintindex{}
\saintindex{}
\feast{1104}{S. Caroli Episcopi et Confessoris}{4 novembre}
\saintindex{}
\saintindex{}
\oratio{Ecclésiam tuam, Dómine, sancti Cároli Confessóris tui atque Pontíficis contínua protectióne custódi: ut, sicut illum pastorális sollicitúdo gloriósum réddidit; ita nos ejus intercéssio in tuo semper fáciat amóre fervéntes. \perdominumla}{Daignez, Seigneur, garder continuellement votre Église sous la protection de saint Charles, votre Pontife et Confesseur ; et comme sa sollicitude pastorale l’a rendu glorieux, que son intercession nous obtienne d’être toujours fervents dans votre amour. \perdominumfr}
\feast{1104}{Ss. Mart. Vitalis et Agricolæ}{4 novembre}
\saintindex{}
\saintindex{}
\oratio{Præsta, quǽsumus, omnípotens Deus: ut, qui sanctórum Mártyrum tuórum Vitális et Agrícolæ solémnia cólimus, eórum apud te intercessiónibus adjuvémur. \perdominumla}{Faites, s’il vous plaît, ô Dieu tout-puissant, que vos saints Martyrs Vital et Agricola, dont nous célébrons la fête, nous assistent par leur intercession auprès de vous. \perdominumfr}
\feast{1104}{Quarta die infra Octavam Omnium Sanctorum}{4 novembre}
\saintindex{}
\saintindex{}
\feast{1105}{Quinta die infra Octavam Omnium Sanctorum}{5 novembre}
\saintindex{}
\saintindex{}
\feast{1106}{Sexta die infra Octavam Omnium Sanctorum}{6 novembre}
\saintindex{}
\saintindex{}
\feast{1107}{Septima die infra Octavam Omnium Sanctorum}{7 novembre}
\saintindex{}
\saintindex{}
\feast{1108}{In Octava Omnium Sanctorum}{8 novembre}
\saintindex{}
\saintindex{}
\feast{1108}{Ss. Quatuor Coronatorum Martyrum}{8 novembre}
\saintindex{}
\saintindex{}
\oratio{Præsta, quǽsumus, omnípotens Deus: ut, qui gloriósos Mártyres fortes in sua confessióne cognóvimus, pios apud te in nostra intercessióne sentiámus. \perdominumla}{Faites, nous vous en prions, Dieu tout-puissant, que connaissant le courage que les glorieux Martyrs ont déployé dans la confession de leur foi, nous ressentions les effets de leur charitable intercession auprès de vous. \perdominumfr}
\feast{1109}{}{9 novembre}
\saintindex{}
\saintindex{}
\feast{1109}{S. Theodori Martyris}{9 novembre}
\saintindex{}
\saintindex{}
\oratio{Deus, qui nos beáti Theodóri Mártyris tui confessióne gloriósa circúmdas et prótegis: præsta nobis ex ejus imitatióne profícere, et oratióne fulcíri. \perdominumla}{Ô Dieu, qui nous donnez la glorieuse profession de foi de votre bienheureux Martyr Théodore, comme appui et protection, accordez-nous la grâce de profiter de ses exemples, et d’être soutenus de ses prières. \perdominumfr}
\feast{1110}{S. Andreæ Avellini Confessoris}{10 novembre}
\saintindex{}
\saintindex{}
\oratio{Deus, qui in corde beáti Andréæ Confessóris tui, per árduum cotídie in virtútibus proficiéndi votum, admirábiles ad te ascensiónes disposuísti: concéde nobis, ipsíus méritis et intercessióne, ita ejúsdem grátiæ partícipes fíeri; ut, perfectióra semper exsequéntes, ad glóriæ tuæ fastígium felíciter perducámur. \perdominumla}{Ô Dieu, qui, par le vœu héroïque de faire chaque jour des progrès dans la vertu, avez disposé dans le cœur de votre Confesseur, le bienheureux André, des degrés admirables pour s’élever à vous ; accordez-nous, par ses mérites et son intercession, de participer à cette même grâce, de sorte que, tendant toujours au plus parfait, nous parvenions heureusement au faîte de votre gloire. \perdominumfr}
\feast{1110}{Ss. Tryphonis, Respicii, et Nymphæ Virg., Martyrum}{10 novembre}
\saintindex{}
\saintindex{}
\oratio{Fac nos, quǽsumus, Dómine, sanctórum Mártyrum tuórum Tryphónis, Respícii et Nymphæ semper festa sectári: quorum suffrágiis, protectiónis tuæ dona sentiámus. \perdominumla}{Faites, nous vous en prions, Seigneur, que nous honorions toujours la fête de vos saints Martyrs, Tryphon, Respice et Nymphe, afin que, favorisés de leurs suffrages, nous éprouvions les bienfaits de votre protection. \perdominumfr}
\feast{1111}{S. Martini Episcopi et Confessoris}{11 novembre}
\saintindex{}
\saintindex{}
\oratio{Deus, qui cónspicis, quia ex nulla nostra virtúte subsístimus: concéde propítius; ut, intercessióne beáti Martíni Confessóris tui atque Pontíficis, contra ómnia advérsa muniámur. \perdominumla}{Ô Dieu, qui voyez que nos forces ne suffisent nullement à nous soutenir, accordez-nous miséricordieusement que, par l'intercession du bienheureux Martin, Votre confesseur et pontife, nous soyons protégés de tout ce qui nous est contraire. \perdominumfr}
\feast{1111}{S. Mennæ Martyris}{11 novembre}
\saintindex{Mennæ}
\saintindex{Menne}
\oratio{@Commune/C2a}{Accordez-nous, s'il Vous plaît, Dieu tout-puissant, que, célébrant la naissance au ciel du bienheureux Menne, Votre martyr, Nous soyons, par son intercession, fortifiés dans l'amour de Votre nom. \perdominumfr}
\feast{1112}{S. Martini Papæ et Martyris}{12 novembre}
\saintindex{Martíni}
\saintindex{}
\feast{1113}{S. Didaci Confessoris}{13 novembre}
\saintindex{}
\saintindex{}
\oratio{Omnípotens sempitérne Deus, qui dispositióne mirábili infírma mundi éligis, ut fórtia quæque confúndas: concéde propítius humilitáti nostræ; ut, piis beáti Dídaci Confessóris tui précibus, ad perénnem in cælis glóriam sublimári mereámur. \perdominumla}{Dieu tout-puissant et éternel, qui, par une providence admirable, choisissez ce qu’il y a de plus faible dans le monde pour confondre ce qu’il y a de plus fort ; soyez propice à notre humilité, et accordez-nous grâce aux pieuses prières de votre bienheureux Confesseur Didace, d’être élevés dans les cieux à la gloire éternelle. \perdominumfr}
\feast{1113}{S. Franciscæ Xaveriæ Cabrini Virginis}{13 novembre}
\saintindex{}
\saintindex{}
\oratio{Dómine Jesu Christe, qui sanctam Vírginem Francíscam Xavériam, Sacratíssimi Cordis tui igne succénsam, per amplíssimas mundi plagas ad ánimas tibi lucrándas deduxísti, et per eam novam in Ecclésia tua Vírginum famíliam suscitásti; quǽsumus; ut ipsa intercedénte, ejúsdem Cordis tui virtútibus induámur, atque ad ætérnum beatitúdinis portum perveníre mereámur. 
$Qui vivis}{}
\feast{1114}{S. Josaphat Episcopi et Martyris}{14 novembre}
\saintindex{}
\saintindex{}
\oratio{Excita, quǽsumus, Dómine, in Ecclésia tua Spíritum, quo replétus beátus Jósaphat Martyr et Póntifex tuus ánimam suam pro óvibus pósuit: ut, eo intercedénte, nos quoque eódem Spíritu moti ac roboráti, ánimam nostram pro frátribus pónere non vereámur. \perdominumejusdemla}{Nous vous en prions, Seigneur, suscitez dans votre Église l’Esprit qui remplissait votre bienheureux Martyr et Pontife Josaphat, et qui le porta à donner sa vie pour ses brebis ; afin qu’étant, par son intercession, animés et fortifiés, nous aussi, de ce même Esprit, nous ne craignions point de sacrifier notre vie pour nos frères. \perdominumejusdemfr}
\feast{1115}{S. Alberti Magni Episcopi Confessoris et Ecclesiæ Doctoris}{15 novembre}
\saintindex{Albérte}
\saintindex{Albert}
\oratio{Deus, qui beátum Albértum Pontíficem tuum atque Doctórem in humána sapiéntia divínæ fídei subiciénda magnum effecísti: da nobis, quǽsumus; ita ejus magistérii inhærére vestígiis, ut luce perfécta fruámur in cælis. \perdominumla}{Ô Dieu, par qui le bienheureux Albert, Votre pontife et docteur, est devenu grand en soumettant l'humble sagesse humaine à la foi divine, donnez-nous, s'il Vous plaît, la grâce de suivre si fidèlement les traces de son enseignement que nous jouissions de la pleine lumière au ciel. \perdominumfr}
\feast{1116}{S. Gertrudis Virginis }{16 novembre}
\saintindex{}
\saintindex{}
\oratio{Deus, qui in corde beátæ Gertrúdis Vírginis jucúndam tibi mansiónem præparásti: ipsíus méritis et intercessióne; cordis nostri máculas cleménter abstérge, et ejúsdem tríbue gaudére consórtio. \perdominumla}{Ô Dieu, qui dans le cœur de la bienheureuse Gertrude votre Vierge, vous êtes préparé une douce demeure, daignez dans votre clémence, par ses mérites et son intercession, effacer les fautes de notre cœur et accordez-nous de jouir de sa compagnie. \perdominumfr}
\feast{1117}{S. Gregorii Thaumaturgi Episcopi et Confessoris}{17 novembre}
\saintindex{Gregórii}
\saintindex{Grégoire}
\oratio{}{Accordez à notre prière, Seigneur tout-puissant, que la vénérable solennité du bienheureux Grégoire votre Confesseur et Pontife augmente en nous dévotion et santé. \perdominumfr}
\feast{1118}{}{18 novembre}
\saintindex{}
\saintindex{}
\feast{1118}{}{18 novembre}
\saintindex{}
\saintindex{}
\feast{1118}{}{18 novembre}
\saintindex{}
\saintindex{}
\feast{1119}{S. Elisabeth Viduæ}{19 novembre}
\saintindex{}
\saintindex{}
\oratio{Tuórum corda fidélium, Deus miserátor, illústra: et, beátæ Elísabeth précibus gloriósis; fac nos próspera mundi despícere, et cælésti semper consolatióne gaudére. \perdominumla}{De vos fidèles, ô Dieu miséricordieux, illuminez les cœurs et par les glorieuses prières de la bienheureuse Elisabeth, faites-nous mépriser les prospérités de ce monde et jouir toujours des célestes consolations. \perdominumfr}
\feast{1119}{S. Pontiani Papæ et Martyris}{19 novembre}
\saintindex{Pontiáni}
\saintindex{Pontien}
\oratio{}{Voyez notre infirmité, Dieu tout-puissant, et puisque le poids de notre activité personnelle nous alourdit, que la glorieuse intercession du bienheureux Pontien, votre Martyr et Pontife, nous protège. \perdominumfr}
\feast{1120}{S. Felicis de Valois Confessoris}{20 novembre}
\saintindex{}
\saintindex{}
\oratio{Deus, qui beátum Felícem Confessórem tuum ex erémo ad munus rediméndi captívos cǽlitus vocáre dignátus es: præsta, quǽsumus; ut per grátiam tuam ex peccatórum nostrórum captivitáte, ejus intercessióne, liberáti, ad cæléstem pátriam perducámur. \perdominumla}{Ô Dieu qui avez daigné appeler par une voix du ciel le bienheureux Félix, votre Confesseur, de la solitude du désert à l’œuvre du rachat des captifs, faites, nous vous en prions, que délivrés par votre grâce et son intercession, de l’esclavage de nos péchés, nous soyons conduits vers la céleste patrie. \perdominumfr}
\feast{1121}{}{21 novembre}
\saintindex{}
\saintindex{}
\oratio{Deus, qui beátam Maríam semper Vírginem, Spíritus Sancti habitáculum, hodiérna die in templo præsentári voluísti: præsta, quǽsumus; ut, ejus intercessióne, in templo glóriæ tuæ præsentári mereámur. \perdominumejusdemla}{Ô Dieu qui avez voulu que la bienheureuse Marie toujours Vierge, demeure de l’Esprit-Saint, soit présentée aujourd’hui au temple, faites, nous vous le demandons, que, par son intercession, nous méritions d’être présentés dans le temple de votre gloire. \perdominumejusdemfr}
\feast{1122}{S. Cæciliæ Virginis et Martyris}{22 novembre}
\saintindex{}
\saintindex{}
\oratio{Deus, qui nos ánnua beátæ Cæcíliæ Vírginis et Mártyris tuæ solemnitáte lætíficas: da, ut, quam venerámur offício, étiam piæ conversatiónis sequámur exémplo. \perdominumfr}{Ô Dieu qui, chaque année, nous réjouissez par la fête de la bienheureuse Cécile, votre Vierge et Martyre, donnez-nous de suivre par l'imitation de sa pieuse vie, celle que nous honorons en cet office. \perdominumfr}
\feast{1123}{S. Clementis Papæ et Martyris}{23 novembre}
\saintindex{Cleméntem}
\saintindex{Clément}
\oratio{Deus, qui nos ánnua beáti Cleméntis Mártyris tui atque Pontíficis solemnitáte lætíficas: concéde propítius; ut cujus natalítia cólimus, virtútem quoque passiónis imitémur. \perdominumla}{}
\feast{1123}{S. Felicitatis Martyris}{23 novembre}
\saintindex{}
\saintindex{}
\oratio{Præsta, quǽsumus, omnípotens Deus: ut, beátæ Felicitátis Mártyris tuæ solémnia recenséntes, méritis ipsíus protegámur et précibus. \perdominumla}{Faites, nous vous le demandons, Dieu tout-puissant, que commémorant la solennité de votre bienheureuse Martyre Félicité, nous soyons protégés par ses mérites et ses prières. \perdominumfr}
\feast{1124}{S. Joannis a Cruce Confessoris et Ecclesiæ Doctoris}{24 novembre}
\saintindex{Joánnes}
\saintindex{Jean}
\oratio{Deus, qui sanctum Joánnem Confessórem tuum atque Doctórem perféctæ sui abnegatiónis et Crucis amatórem exímium effecísti: concéde; ut, ejus imitatióni júgiter inhæréntes, glóriam assequámur ætérnam. \perdominumla}{Ô Dieu, qui avez inspiré à saint Jean, votre Confesseur et Docteur, un rare amour de la croix et d’une parfaite abnégation de soi-même, faites que, nous appliquant sans cesse à l’imiter, nous obtenions la gloire éternelle. \perdominumfr}
\oratio{@:Oratio_
(sed rubrica tridentina)
@:Oratio_:s/ atque Doctorem//}{@:Oratio_
(sed rubrica tridentina)
@:Oratio_:s/ et Docteur//}
\feast{1124}{S. Andre Dung Lac and his companions, martyrs}{24 novembre}
\saintindex{}
\saintindex{}
\feast{1124}{S. Chrysogoni Martyris}{24 novembre}
\saintindex{}
\saintindex{}
\oratio{Adésto, Dómine, supplicatiónibus nostris: ut, qui ex iniquitáte nostra reos nos esse cognóscimus, beáti Chrysógoni Mártyris tui intercessióne liberémur. \perdominumla}{Agréez, Seigneur, nos supplications, afin que nous reconnaissant coupables à cause de notre iniquité, nous en soyons libérés par l’intercession du bienheureux Chrysogone, votre Martyr. \perdominumfr}
\feast{1125}{S. Catharinæ Virginis et Martyris}{25 novembre}
\saintindex{Catharínæ}
\saintindex{Catherine}
\oratio{Deus, qui dedísti legem Móysi in summitáte montis Sínai, et in eódem loco per sanctos Angelos tuos corpus beátæ Catharínæ Vírginis et Mártyris tuæ mirabíliter collocásti: præsta, quǽsumus; ut, ejus méritis et intercessióne, ad montem, qui Christus est, perveníre valeámus:
$Qui tecum}{Ô Dieu, qui avez donné votre loi à Moïse, sur le sommet du mont Sinaï, et avez fait miraculeusement transférer au même lieu, par vos saints Anges, le corps de la bienheureuse Catherine, votre Vierge et Martyre, accordez à notre demande que, par ses mérites et son intercession, nous puissions parvenir à cette montagne qui est le Christ :
$Qui tecum}
\feast{1126}{S. Silvestri Abbatis}{26 novembre}
\saintindex{}
\saintindex{}
\oratio{Clementíssime Deus, qui sanctum Silvéstrum Abbátem, sǽculi hujus vanitátem in apérto túmulo pie meditántem, ad erémum vocáre, et præcláris vitæ méritis decoráre dignátus es: te súpplices exorámus; ut, ejus exémplo terréna despiciéntes, tui consórtio perfruámur ætérno. \perdominumla}{Ô Dieu très clément, qui avez daigné appeler au désert saint Silvestre, Abbé, méditant avec piété la vanité de ce monde, devant un tombeau ouvert et illustrer sa vie de mérites éclatants, accordez-nous, nous Vous en supplions, de mépriser, à son exemple, les biens terrestres pour jouir de Votre éternelle compagnie.  \perdominumfr}
\feast{1126}{S. Petri Alexandrini Martyris}{26 novembre}
\saintindex{Petri}
\saintindex{Pierre}
\feast{1129}{In Vigilia S. Andreæ Apostoli}{29 novembre}
\saintindex{Andreas}
\saintindex{André}
\oratio{@Commune/C1v:Oratio 2 loco}{}
\feast{1129}{S. Saturnini Martyris}{29 novembre}
\saintindex{}
\saintindex{}
\oratio{Deus, qui nos beáti Saturníni Mártyris tui concédis natalício pérfrui: ejus nos tríbue méritis adjuvári. \perdominumla}{Ô Dieu, qui nous faites la grâce de nous réjouir de la naissance céleste du bienheureux Saturnin, votre Martyr, accordez-nous d’être secourus par ses mérites. \perdominumfr
}
\feast{1130}{S. Andreæ Apostoli}{30 novembre}
\saintindex{}
\saintindex{}
\oratio{Majestátem tuam, Dómine, supplíciter exorámus: ut, sicut Ecclésiæ tuæ beátus Andréas Apóstolus éxstitit prædicátor et rector; ita apud te sit pro nobis perpétuus intercéssor. \perdominumla}{Nous prions instamment et avec humilité, Seigneur, Votre Majesté, afin que le bienheureux Apôtre André, qui a été prédicateur et administrateur de Votre Église, soit au même degré, pour nous, un intercesseur perpétuel auprès de Vous. \perdominumfr}
\feast{1202}{S. Bibianæ Virginis et Martyris}{2 décembre}
\saintindex{Bibiána}
\saintindex{Bibiane}
\oratio{Deus, ómnium largítor bonórum, qui in fámula tua Bibiána cum virginitátis flore martýrii palmam conjunxísti: mentes nostras ejus intercessióne tibi caritáte conjúnge; ut, amótis perículis, prǽmia consequámur ætérna. \perdominumla}{Ô Dieu, dispensateur de tous les biens, qui, en votre servante Bibiane, à la fleur de la virginité, avez joint la palme du martyre, faites, par son intercession, que nos esprits vous soient unis par la charité, afin que, tout péril écarté, nous obtenions les récompenses éternelles. \perdominumfr}
\feast{1203}{S. Francisci Xaverii Confessoris}{3 décembre}
\saintindex{}
\saintindex{}
\oratio{Deus, qui Indiárum gentes beáti Francísci prædicatióne et miráculis Ecclésiæ tuæ aggregáre voluísti: concéde propítius; ut cujus gloriósa mérita venerámur, virtútum quoque imitémur exémpla. \perdominumla}{Ô Dieu qui avez voulu réunir à votre Église les peuples des Indes, par la prédication et les miracles du bienheureux François, accordez-nous miséricordieusement, d’imiter aussi les exemples de celui dont nous vénérons les glorieux mérites. \perdominumfr}
\feast{1204}{S. Petri Chrysologi Episcopi Confessoris et Ecclesiæ Doctoris}{4 décembre}
\saintindex{Petre Chrysóloge}
\saintindex{Pierre Chrysologue}
\oratio{Deus, qui beátum Petrum Chrysólogum Doctórem egrégium, divínitus præmonstrátum, ad regéndam et instruéndam Ecclésiam tuam éligi voluísti: præsta, quǽsumus; ut, quem Doctórem vitæ habúimus in terris, intercessórem habére mereámur in cælis. \perdominumla}{Ô Dieu, qui avez voulu que le bienheureux Pierre Chrysologue, Docteur éminent, fût choisi par désignation miraculeuse, pour gouverner et enseigner Votre Église, faites, nous Vous en prions, qu’après l’avoir eu comme Docteur de vie sur terre, nous méritions de l’avoir pour intercesseur dans les cieux. \perdominumfr}
\feast{1204}{S. Barbaræ Virginis et Martyris}{4 décembre}
\saintindex{Bárbaræ}
\saintindex{Barbe}
\oratio{}{Ô Dieu qui, parmi toutes les autres merveilles de Votre puissance, avez attribué la victoire du martyre même au sexe faible, accordez-nous, dans Votre bienveillance, qu’honorant la naissance céleste de la bienheureuse Barbe, Votre Vierge et Martyre, nous avancions vers Vous, aidés par ses exemples. \perdominumfr
}
\feast{1205}{S. Sabbæ Abbatis}{5 décembre}
\saintindex{Sabbæ}
\saintindex{Sabba}
\feast{1206}{S. Nicolai Episcopi et Confessoris}{6 décembre}
\saintindex{}
\saintindex{}
\oratio{Deus, qui beátum Nicoláum Pontíficem innúmeris decorásti miráculis: tríbue, quǽsumus; ut ejus méritis et précibus a gehénnæ incéndiis liberémur. \perdominumla}{Ô Dieu, qui avez glorifié le bienheureux Pontife Nicolas par d’innombrables miracles : accordez à notre demande que par ses mérites et ses prières, nous soyons préservés des feux de l’enfer. \perdominumfr}
\feast{1207}{S. Ambrosii Episcopi Confessoris et Ecclesiæ Doctoris}{7 décembre}
\saintindex{Ambrósium}
\saintindex{Ambroise}
\oratio{}{Ô Dieu, qui avez accordé à votre peuple le bienheureux Ambroise, comme ministre du salut éternel, faites, nous vous en prions, que l'ayant eu sur la terre, comme Docteur de vie, nous méritions de l’avoir comme intercesseur dans les cieux. \perdominumfr}
\feast{1208}{}{8 décembre}
\saintindex{}
\saintindex{}
\oratio{Deus, qui per immaculátam Vírginis Conceptiónem dignum Fílio tuo habitáculum præparásti: quǽsumus; ut qui ex morte ejúsdem Fílii tui prævísa, eam ab omni labe præservásti, nos quoque mundos ejus intercessióne ad te perveníre concédas. \pereumdemla}{Ô Dieu, qui, par l’immaculée Conception de la Vierge, avez préparé à votre Fils une demeure digne de lui, et qui, en prévision de la mort de ce même Fils, avez préservé celle-ci de toute souillure, accordez-nous, nous vous le demandons, d’arriver jusqu’à vous, purs nous aussi, par son intercession. \pereumdemfr}
\feast{1208}{}{8 décembre}
\saintindex{}
\saintindex{}
\oratio{Famulis tuis, quǽsumus Domine, cælestis gratiæ munus impertire: ut quibus beatæ Virginis partus exstitit salutis exordium, Conceptionis ejus votiva solemnitas pacis tribuat incrementum. \perdominumla}{À vos serviteurs, nous vous prions, Seigneur, accordez le don de la grâce céleste, afin que pour ceux pour qui la naissance de la Bienheureuse Vierge fut le commencement du salut, la solennité votive de sa Conception apporte une augmentation de paix. \perdominumfr}
\feast{1209}{}{9 décembre}
\saintindex{}
\saintindex{}
\feast{1209}{}{9 décembre}
\saintindex{}
\saintindex{}
\feast{1209}{S. Juan Diego}{9 décembre}
\saintindex{}
\saintindex{}
\feast{1210}{}{10 décembre}
\saintindex{}
\saintindex{}
\feast{1210}{S. Melchiadis Papæ et Martyris}{10 décembre}
\saintindex{Melchíadis}
\saintindex{Melchiade}
\feast{1211}{S. Damasi Papæ et Confessoris}{11 décembre}
\saintindex{Dámaso}
\saintindex{Dámase}
\oratio{Exáudi, Dómine, preces nostras: et interveniénte beáto Dámaso Confessóre tuo atque Pontífice, indulgéntiam nobis tríbue placátus et pacem. \perdominumla}{}
\feast{1211}{}{11 décembre}
\saintindex{}
\saintindex{}
\feast{1212}{}{12 décembre}
\saintindex{}
\saintindex{}
\feast{1212}{}{12 décembre}
\saintindex{}
\saintindex{}
\feast{1212}{}{12 décembre}
\saintindex{}
\saintindex{}
\oratio{Deus, qui sub beatíssimæ Vírginis Maríæ singulári patrocínio constitútos perpétuis benefíciis nos cumulári voluísti: præsta supplícibus tuis; ut, cujus hódie commemoratióne lætámur in terris, ejus conspéctu perfruámur in cælis. \perdominumla}{Ô Dieu, qui avez voulu nous combler de vos bienfaits grâce à la protection spéciale de la très sainte Vierge Marie, accordez à vos suppliants que nous nous réjouissions aujourd'hui de sa commémoration sur terre, et que nous jouissions de sa présence dans les cieux. \perdominumfr}
\feast{1213}{S. Luciæ Virginis et Martyris}{13 décembre}
\saintindex{Lúciæ Vírginis et Mártyris tuæ}
\saintindex{Lucie, votre Vierge et Martyre}
\oratio{@Commune/C7a}{}
\feast{1213}{}{13 décembre}
\saintindex{}
\saintindex{}
\feast{1213}{}{13 décembre}
\saintindex{}
\saintindex{}
\feast{1213}{}{13 décembre}
\saintindex{}
\saintindex{}
\feast{1214}{}{14 décembre}
\saintindex{}
\saintindex{}
\feast{1215}{}{15 décembre}
\saintindex{}
\saintindex{}
\feast{1216}{S. Eusebii Episcopi et Martyris}{16 décembre}
\saintindex{Eusébii}
\saintindex{Eusèbe}
\oratio{}{Ô Dieu, qui nous réjouissez par la fête annuelle de votre bienheureux Eusèbe, Martyr et Pontife, faites nous cette faveur que, célébrant sa naissance céleste, nous nous réjouissions aussi de sa protection. \perdominumfr}
\feast{1220}{In Vigilia S. Thomæ Ap.}{20 décembre}
\saintindex{Thomæ}
\saintindex{Thomas}

\end{document}